\bookpart{多模态处理}{Multimodal Processing}
\label{part:multimodal-processing}

\BookPartGuide{%
文字、图像和声音可以共同提供事件信息，也可能互相缺失或冲突。本部分研究这些信息如何连接，怎样形成有证据支持的理解结果，以及怎样制作满足条件和跨输出关系的媒体成果。

本部分使用第一部分的对齐与融合原理、第二部分建立的多模态模型接口，以及图像和视频任务形成的区域、时间与观测依据。多模态理解章把匹配、定位、读取和结构恢复用于描述与问答，保留答案的来源依据。多模态生成章研究视频、视觉配声、联合音视频、空间内容和文字语音响应，核对单项质量、条件遵循、时间同步与制作预算。}

\chapter{多模态理解}
\label{chap:vision-multimodal-understanding}

给一张街景照片，要求找出“右边拿杯子的人”，结果应落到一个可检查的区域。给一页票据，询问“总额是多少”，模型需要读出数字及其字段归属。给一段视频，追问“第二次敲门发生在讲话之前还是之后”，还需对齐事件、声音与时间。三个请求都可以用一句话回答，但支持这句话的证据并不相同：第一个需要区域，第二个需要字符串和字段关系，第三个需要带时间的事件记录。

本章要回答的问题是：怎样把文字、图像、视频、声音与页面结构联系起来，找到支持结果的证据，并给出可以复核的描述和答案。\ref{chap:vision-multimodal-alignment}章已经讲过跨模态对齐与融合的基本原理，包括典型相关分析、动态时间规整、注意力对齐、对比对齐以及早期、中间、晚期三种融合位置；\ref{chap:vision-image-features}章讲过边缘、局部特征、阈值分割和模板匹配等经典图像特征；\ref{chap:vision-multimodal-models}章讲过多模态模型怎样匹配表示、连接模态和解码语言；\ref{chap:vision-image-recognition}章和\ref{chap:vision-video-recognition}章讲过图像与视频中的检测、跟踪和检索。本章使用这些原理和接口，按请求所需的成果展开，用到时只回顾结论并指回原处。

路线由证据从粗到细的获取顺序决定。第一节先规定任务的输入、输出与评价指标，并说明多模态理解在对齐粒度、监督质量、语言先验、观测分辨率和证据组合上的困难。面对大量图片或视频，先要缩小候选范围，这是跨模态匹配与检索；候选找到后，要把查询落到区域、时间或控件上，这是跨模态定位；票据、表格和图表的内容写在像素里，还需读取文字并恢复结构；有了区域、文字和事件，才能选择内容形成描述，或者组合证据回答问题。最后一节讨论怎样分层核验这些结果。贯穿全章的是来源记录：每个判断依据哪幅图、哪一页、哪个区域或哪段时间。输入额外的OCR结果、字幕或真实数据表，会改变模型仍需完成的工作，因而也改变评价结论。

\section{多模态理解的问题与挑战}
\label{sec:vision-understanding-challenges}

开篇的三个请求有一个共同结构：系统收到一份或多份观测，例如图像、视频帧、音频或页面截图，同时收到一个用另一种模态表达的请求，例如一句话、一段声音或一幅示例图；它要输出一个能回到观测中检查的结果。按结果的形式，多模态理解的输出可以归为四类。第一类是排序列表，即从候选库中挑出相关项并给出名次；第二类是位置，即框、掩码、点击点或时间区间；第三类是字符串与结构，即读出的文字、表格的行列、表单的键值、图表的数据和图解的连接；第四类是自然语言，即描述和答案。不同输出能被检查的方式不同，评价指标也随之不同。后文各节比较模型时都要固定一种指标，因此这里先把指标定义清楚。

排序列表用召回类指标评价。设查询集合为$Q$，对查询$q$，相关候选的集合为$R_q$，模型排在前$K$名的候选集合为$S_q^K$。图文检索中最常用的是前$K$召回率（Recall@$K$，记作R@$K$）：
\begin{equation}
  \mathrm{R@}K=\frac{1}{|Q|}\sum_{q\in Q}
    \bm 1\bigl[S_q^K\cap R_q\neq\varnothing\bigr]\text{，}
  \label{eq:vision-understanding-recall-at-k}
\end{equation}
其中$\bm 1[\cdot]$为指示函数，条件成立取1，否则取0。它统计“前$K$名里至少出现一个正例”的查询比例。一个查询有多个正例时，也可以改用覆盖比例$|S_q^K\cap R_q|/|R_q|$，两种定义在多对多数据上给出不同数值，报告时必须写明。另一个常用量是\emph{中位名次}（median rank, MedR）：每个查询第一个正例的名次取中位数，越小越好。以三个查询为教学例，若第一个正例分别排在第1、4、12名，则R@1$=1/3$，R@5$=2/3$，R@10$=2/3$，MedR$=4$。R@$K$只关心正例是否进入前$K$名，名次在前$K$之内如何排列并不影响它，这正是后文把检索拆成召回与重排的依据。

位置输出用重叠度评价。框$B$与标注框$B^*$的\term{交并比}（intersection over union, IoU）定义为
\begin{equation}
  \operatorname{IoU}(B,B^*)=\frac{|B\cap B^*|}{|B\cup B^*|}\text{，}
  \label{eq:vision-understanding-iou}
\end{equation}
$|\cdot|$表示面积。指代表达定位通常报告IoU不小于0.5的样本比例，即Acc@0.5。时间区间用同一公式，只是把面积换成长度，称为时间IoU。例如框$[0,4]\times[0,2]$与标注框$[2,6]\times[0,2]$的交集面积为4，并集面积为12，IoU$=1/3$，按0.5的阈值记为错误；区间$[1,3)$与$[2,4)$的时间IoU同样是$1/3$。界面定位输出的是点击点，常用的口径是点是否落在目标控件框内；声源定位输出热图，要先按阈值取出区域，再与标注的声源区域比较。阈值本身是协议的一部分：同一组输出，在0.5与0.7两个阈值下可以给出差别很大的准确率。

字符串用编辑距离评价。两个字符串之间的\term{编辑距离}（edit distance），又称Levenshtein距离，是把一个串变成另一个串所需的最少替换、删除和插入次数。设参考串$\bm a=a_1\cdots a_n$与识别串$\bm b=b_1\cdots b_m$，记$d_{i,j}$为$a_1\cdots a_i$与$b_1\cdots b_j$之间的编辑距离。最后一步只有三种可能：删去$a_i$、插入$b_j$，或者让$a_i$与$b_j$对应（相同则不计代价，不同则计一次替换），于是
\begin{equation}
  d_{i,j}=\min\bigl\{d_{i-1,j}+1,\ d_{i,j-1}+1,\
    d_{i-1,j-1}+\bm 1[a_i\neq b_j]\bigr\}\text{，}
  \qquad d_{i,0}=i,\ d_{0,j}=j\text{。}
  \label{eq:vision-understanding-edit-distance}
\end{equation}
边界条件说明空串与长度为$i$的串之间需要$i$次插入或删除。这个递推是动态规划：$d_{i,j}$只依赖左、上、左上三个已经算好的格，按行填表即可，代价为$O(nm)$。动态规划把问题分解为重叠子问题、按依赖顺序求解的原理，可参考\BookChapterRef{chap:graph-theory}{卷一“图的表示与计算”章}中的最短路递推；编辑距离的递推相当于在$(n+1)\times(m+1)$的网格上，从左上角走到右下角的最短路径，横、竖每步代价为1，斜步代价为0或1。算法\ref{alg:vision-understanding-edit-distance}写出了计算步骤。

\begin{algorithm}[字符串编辑距离的动态规划计算]
  \label{alg:vision-understanding-edit-distance}
  \begin{algorithmic}[1]
    \Require 参考串$a_1\cdots a_n$，识别串$b_1\cdots b_m$
    \Ensure 编辑距离$d_{n,m}$
    \For{$i=0,1,\ldots,n$}
      \State $d_{i,0}\gets i$
    \EndFor
    \For{$j=0,1,\ldots,m$}
      \State $d_{0,j}\gets j$
    \EndFor
    \For{$i=1,\ldots,n$}
      \For{$j=1,\ldots,m$}
        \State 按式\eqref{eq:vision-understanding-edit-distance}取三种来源的最小值作为$d_{i,j}$
      \EndFor
    \EndFor
    \State \Return $d_{n,m}$；需要具体的编辑操作时，从$(n,m)$沿取得最小值的来源回溯到$(0,0)$
  \end{algorithmic}
\end{algorithm}

以教学票据编号为例，参考串为“A017”，识别器把数字0读成字母O并漏掉了1，输出“AO7”。表\ref{tab:vision-understanding-edit-example}按行填出$d_{i,j}$。右下角为2，对应一次替换（0读成O）和一次删除（漏掉1）。

\begin{table}[htbp]
  \centering
  \small
  \begin{tabular}{@{}lcccc@{}}
    \toprule
     & 空 & A & O & 7\\
    \midrule
    空 & 0 & 1 & 2 & 3\\
    A & 1 & 0 & 1 & 2\\
    0 & 2 & 1 & 1 & 2\\
    1 & 3 & 2 & 2 & 2\\
    7 & 4 & 3 & 3 & 2\\
    \bottomrule
  \end{tabular}
  \caption{参考串“A017”（行）与识别串“AO7”（列）的编辑距离表，为教学构造。每格取上方加1、左方加1、左上方加替换代价三者的最小值；例如第“0”行“O”列取左上方的0加1，得到1。右下角的2是两串的编辑距离。}
  \label{tab:vision-understanding-edit-example}
\end{table}

用编辑距离可以定义\term{字符错误率}（character error rate, CER）。若参考串含$N>0$个字符，最优编辑中的替换、删除和插入次数分别为$S,D,I$，则
\begin{equation}
  \mathrm{CER}=\frac{S+D+I}{N}\text{，}
  \label{eq:vision-cer}
\end{equation}
分子就是编辑距离。上例中CER$=2/4=0.5$。分母使用参考字符数，因此插入很多字符时CER可以超过1；空格、标点和大小写是否计入也要事先规定。以词为单位做同样的计算，得到词错误率。

答案中的字符串还有一种更宽容的指标。场景文字问答与文档问答常用\term{平均归一化编辑相似度}（average normalized Levenshtein similarity, ANLS）\scite{vision-stvqa2019}。设第$i$个问题有若干参考答案$a_{i,j}$，模型答案为$o_i$，归一化编辑距离$\operatorname{NL}$等于编辑距离除以两串中较长者的长度，取值在0到1之间。定义
\begin{equation}
  \mathrm{ANLS}=\frac1N\sum_{i=1}^{N}\max_j s(a_{i,j},o_i)\text{，}\qquad
  s(a,o)=\begin{cases}
    1-\operatorname{NL}(a,o), & \operatorname{NL}(a,o)<\tau,\\
    0, & \operatorname{NL}(a,o)\geq\tau,
  \end{cases}
  \label{eq:vision-understanding-anls}
\end{equation}
原文取$\tau=0.5$。阈值的用意是：错了一半以上的字符串更可能选错了文字，而不是读错了个别字符，此时不给部分分。例如参考答案“128.50”，回答“128.5”的编辑距离为1，$\operatorname{NL}=1/6$，得分约0.83；回答“100.00”需要三次替换，$\operatorname{NL}=3/6=0.5$，得分为0。ANLS宽容OCR的小错，但对金额这样的字段，少一位小数可能已经改变含义，因此仍要配合整串正确率。

表格结构需要比较两张表的行列和跨格关系，而不只是比较字符串。\term{树编辑距离相似度}（tree-edit-distance-based similarity, TEDS）把表格写成HTML树，计算
\begin{equation}
  \operatorname{TEDS}(T_a,T_b)=1-\frac{\operatorname{EditDist}(T_a,T_b)}{\max(|T_a|,|T_b|)}\text{，}
  \label{eq:vision-understanding-teds}
\end{equation}
其中$\operatorname{EditDist}$为把一棵树变成另一棵所需的最少节点插入、删除和替换代价，$|T|$为节点数\scite{vision-pubtabnet2020}。它把字符串编辑距离推广到树：节点是表格、行和单元格，单元格内容不同也计入替换代价。PubTables-1M另用GriTS系列指标，把表格看成单元格矩阵，在两个矩阵最相似的行列子结构上计算F分数形式的相似度，并分别考察拓扑、内容和位置\scite{vision-pubtables2022}。页面版面分析则把标题、正文、表格等区域当作检测对象，用\term{平均精度}（average precision, AP）评价。

视觉问答的答案通常很短。VQA数据集为每个问题收集10个人工答案，单个回答的准确率定义为$\min(\text{给出该答案的人数}/3,\,1)$，并在10个答案中任取9个形成的10个子集上取平均\scite{vision-vqa2015}。若10人中只有2人回答“蓝色”，模型也回答“蓝色”：在包含这2人的8个子集上得分$2/3$，在只含其中1人的2个子集上得分$1/3$，平均为$(8\times2/3+2\times1/3)/10=0.6$。这个定义承认答案的多样性，三人以上同意即得满分；但它只比较答案字符串，不检查答案依据的是图中哪一部分。图表问答中的数值答案常用放宽准确率：与正确值相差不超过5\%即算正确\scite{vision-chartqa2022}。

描述的评价要比较一句生成文字与若干参考描述。BLEU计算候选句$n$元组在参考句中出现的比例，即$n$元精确率$p_n$，同一个$n$元组的计数不超过它在任一参考句中出现的最多次数，以免重复一个高频词就得到高分；再用简短惩罚$\mathrm{BP}$抑制过短的句子\scite{papineni-etal-2002-bleu}：
\begin{equation}
  \mathrm{BLEU}=\mathrm{BP}\cdot\exp\Bigl(\sum_{n=1}^{N}w_n\log p_n\Bigr)\text{，}\qquad
  \mathrm{BP}=\begin{cases}1, & c>r,\\ e^{1-r/c}, & c\leq r,\end{cases}
  \label{eq:vision-understanding-bleu}
\end{equation}
其中$c$为候选长度，$r$为参考长度，常用$N=4$与$w_n=1/4$。CIDEr针对图像描述设计：它先把每个$n$元组按TF-IDF加权，即该$n$元组在句中出现越多权重越高，在整个数据集的描述中越普遍权重越低；再计算候选句与每句参考描述的加权向量余弦相似度并取平均，最后对$n=1,\ldots,4$等权求和\scite{vision-cider2015}。设$\bm g_n(c)$为候选句$c$在长度为$n$的全部$n$元组上的TF-IDF向量，$s_1,\ldots,s_m$为参考描述，则
\begin{equation}
 \begin{gathered}
  \mathrm{CIDEr}_n(c)=\frac1m\sum_{j=1}^{m}
    \frac{\bm g_n(c)^\top\bm g_n(s_j)}{\lVert\bm g_n(c)\rVert_2\,\lVert\bm g_n(s_j)\rVert_2}\text{，}\\
  \mathrm{CIDEr}(c)=\sum_{n=1}^{4}\frac14\,\mathrm{CIDEr}_n(c)\text{。}
 \end{gathered}
  \label{eq:vision-understanding-cider}
\end{equation}
IDF权重压低了“一个”“在”这类在各图描述中都会出现的词组，突出与当前图像相关的内容词。常用的CIDEr-D版本另加长度惩罚和计数截断，并把分数乘以10。两种指标衡量的都是与参考描述的措辞重合，并不直接检查描述中的对象是否真的出现在图中；参考描述没有提到但图中确实存在的内容，同样得不到分数。

表\ref{tab:vision-understanding-metrics}汇总这些指标及其检查范围。每个指标只检查一种成果，任务需要同时满足几种条件时，要分别报告，或者按\ref{sec:vision-understanding-grounded-evaluation}节的方式取交集。

\begin{table}[htbp]
  \centering
  \small
  \begin{tabularx}{\linewidth}{@{}p{14mm}p{20mm}XX@{}}
    \toprule
    输出 & 指标 & 检查内容 & 不能检查的内容\\
    \midrule
    排序列表 & R@$K$、MedR & 正例是否进入前$K$名、第一个正例的名次 & 正例在前$K$名内的顺序；关系是否由同一对象满足\\
    位置 & IoU、Acc@0.5、点命中 & 区域或区间与标注的重叠 & 选中的是否为查询所指对象之外的同类对象\\
    字符串 & CER、整串正确率 & 字符级差异 & 字符串属于哪个字段、单位是否正确\\
    表格与页面 & TEDS、GriTS、AP & 行列、跨格、区域类别 & 内容的数值运算是否正确\\
    答案 & VQA准确率、ANLS、放宽准确率 & 答案字符串或数值 & 答案依据的区域和时刻\\
    描述 & BLEU、CIDEr & 与参考描述的$n$元组重合 & 描述中的对象和关系是否有图像依据\\
    \bottomrule
  \end{tabularx}
  \caption{多模态理解的主要评价指标。中间一列是指标直接计算的内容，最后一列是同一指标无法发现的错误，需要另设检查。}
  \label{tab:vision-understanding-metrics}
\end{table}

指标定义了“做对”的标准，下面各小节讨论为什么做对很难。困难来自多模态数据本身：不同模态的对应关系有不同粒度，训练配对往往不精确，语言一侧的先验很强，视觉一侧的证据又可能小而分散，还常常需要多步组合。每个困难都有经典方法的处理方式，也有学习方法带来的改变。

\subsection{多粒度对应带来的对齐困难}
\label{sec:vision-understanding-challenge-granularity}

“穿红衣服的人正在骑车”与一幅街景之间的对应，可以在几个粒度上成立。整句与整幅图可以对应，说明这幅图大体相关；短语“红衣服的人”与图中某个人物区域对应，说明属性绑定到了哪个对象；视频中“搬起”这个词与若干帧对应，说明动作发生在哪段时间；敲门声与手碰到门的那几帧对应，说明声音和画面同步。这些粒度不能相互替代：整图与整句的匹配分数再高，也推不出红色属于骑车的人；短语与区域对齐了，也推不出动作发生在哪一秒。困难在于，系统必须在任务所需的粒度上保留对应关系，而粗粒度的表示往往已经丢掉了细粒度对应所需的位置和时间信息。

经典方法把对齐写成显式的优化问题，见\ref{chap:vision-multimodal-alignment}章。典型相关分析（\ref{sec:vision-alignment-cca}节）为两组成对的特征各学一个线性投影，使投影后的相关系数最大，从而得到一个公共空间；它处理的是整体粒度，每个样本只有一个向量，无法表达区域与短语的对应。动态时间规整（\ref{sec:vision-alignment-dtw}节）在单调和端点约束下，用动态规划找出两条序列之间累计代价最小的对应路径，适合字幕与帧、音频与视频这类保持先后顺序的序列；但它要求事先给出逐帧之间的距离，距离本身就要求两种特征可比。这些方法所用的视觉特征来自\ref{chap:vision-image-features}章的人工设计，例如HOG、SIFT和视觉词袋，它们描述外观，却不知道哪个外观对应哪个词。

学习方法改变的是两点。其一，特征和对齐一起学习：对比对齐（\ref{sec:vision-alignment-contrastive}节）在大规模配对数据上同时训练两个编码器，让配对的表示彼此靠近、不配对的彼此远离，可比性由数据学出，不再依赖人工设计的特征。其二，细粒度对应可以通过注意力表达（\ref{sec:vision-alignment-attention}节）：每个词对每个区域计算一个权重，粒度由参与注意力的单元决定。代价是，越细的粒度越需要相应的监督或自监督信号，例如短语与框的标注、音画同时出现的关系。本章\ref{sec:vision-understanding-retrieval}节在整体粒度上召回候选，再用交叉注意力或逐条件核验处理关系绑定；\ref{sec:vision-understanding-grounding}节分别处理短语与区域、声音与区域、音频与视频片段、指令与控件四种细粒度对应。

\subsection{弱配对与噪声监督带来的学习偏差}
\label{sec:vision-understanding-challenge-weak-pairing}

大规模多模态数据大多不是为当前任务标注的。网页图片的替代文字可能只写“图片1”，教学视频的旁白可能先说“接下来揉面”，几秒后手才开始动；视频的声音可能来自画面之外的人。这类数据只告诉我们“两种模态来自同一来源”，没有说明它们在细粒度上怎样对应，有时甚至并不对应。困难在于：用这样的配对训练，模型学到的可能是与目标无关的共现，例如“厨房背景”与“做饭”一词总是同时出现，模型便只凭背景判断，而不去看手上的动作。

经典做法依靠小规模人工标注：为图像标出类别、为短语标出框，再用固定类别的检测器和人工特征完成匹配。\ref{sec:vision-alignment-weak-pairing}节把这种只在整体层面成立、内部单元对应未知的配对称为弱配对，并区分了假正例与假负例两种噪声；在人工标注的范围内，监督是干净的，但类别和场景的覆盖受标注成本限制。

学习方法用规模和训练目标的设计来对抗噪声。规模大时，单条错配对梯度的影响被大量正确配对冲淡；目标函数可以只要求真实配对比错配高出一个间隔，满足后不再施压，减少对个别噪声样本的过拟合；还可以专门构造困难负例，例如同一视频中的另一句旁白、同一幅图中的背景区域，迫使模型区分细节而非场景。另一类做法是自举：用当前模型给出伪标注，再按可信度加权训练。这些方法降低了标注成本，也带来新的风险：伪标注来自模型自己，错误的高分区域可能被不断强化。本章的HowTo100M（\ref{sec:vision-understanding-howto}节）、声音来源定位（\ref{sec:vision-understanding-sound-source}节）和CRAFT（\ref{sec:vision-understanding-craft}节）分别代表了间隔排序与同源负例、同图困难负例和带置信度的伪标注三种处理方式；VisionTaPas（\ref{sec:vision-understanding-visiontapas}节）用启发式规则生成单元监督，也属于这一类，原论文报告了其中的噪声比例。

\subsection{语言先验带来的捷径与幻觉}
\label{sec:vision-understanding-challenge-language-prior}

问“天空是什么颜色”，不看图回答“蓝色”大多是对的；描述一张街景，写上“路边停着汽车”也常常是对的。语言中的统计规律，即\term{语言先验}（language prior），让模型可以不看视觉证据就给出看似合理的输出。这是一种\term{捷径学习}（shortcut learning）：模型利用了训练数据中与答案相关、却与任务本意无关的线索。它造成两种错误。一种是问题与常识不符时答错，例如阴天的照片仍被答成“蓝色”；另一种是\term{视觉幻觉}（visual hallucination），即输出图中没有依据的对象或属性。OCR也受同样的影响：模糊的“1”和“I”可能被语言上下文改写成常见词，票据编号这样的罕见字符串反而被“纠正”错。

经典流水线中，视觉判断和语言输出是分开的两步，语言部分只能处理视觉部分已经给出的检测结果或字符，先验的作用范围受到限制，但也无法利用语言帮助消歧。\ref{chap:vision-multimodal-alignment}章从融合角度讨论过同一问题：晚期融合（\ref{sec:vision-alignment-late-fusion}节）让各模态分别给出判断再汇总，每个模态的贡献可以单独检查；而模态不平衡（\ref{sec:vision-alignment-imbalance}节）说明，联合训练时更容易拟合的那个模态会主导优化，较弱的模态得不到充分学习。在视觉问答中，问题文字往往就是更容易拟合的一方。

学习方法，尤其是接入大型语言模型之后，语言先验更强了：模型能写出流畅、合乎常识的长段文字，却不一定逐句有视觉依据。应对的方向有三类：在数据上构造反例，例如为同一问题收集答案不同的相似图像；在输出中保留证据，例如让答案附带区域、单元或时间区间；在评价上单独检查，例如询问对象是否存在、替换或移除图像观察答案是否变化。本章\ref{sec:vision-understanding-description}节讨论描述中的幻觉，\ref{sec:vision-understanding-question}节讨论问答中的语言捷径，\ref{sec:vision-understanding-input-protocol}节给出检查模态贡献的方法。

\subsection{高分辨率与长输入带来的证据丢失}
\label{sec:vision-understanding-challenge-resolution}

决定答案的证据常常很小。专业软件的截图上，目标控件平均只占屏幕面积的0.07\%\scite{vision-screenseeker2025}；票据上“128.50”的小数点只有几个像素；一段一小时的视频里，第二次敲门只持续一秒。而视觉编码器的输入大小是固定的，例如224、384或448像素见方，视频也只能抽取有限的帧。把整页、整屏或整段视频压进这个预算，小字、小控件和短暂事件可能在编码之前就已经消失，后面的模型再强也无法恢复。

经典方法对这个问题有成熟的分析。\ref{sec:vision-formation-aliasing}节和\ref{sec:vision-formation-interpolation}节说明，缩小图像时若不先做低通滤波就会产生混叠，而做了低通滤波，小于新采样间隔的细节就被平滑掉了，缩小必然以丢失细节为代价。图像金字塔（\ref{sec:vision-transforms-pyramid}节）在多个尺度上分别处理，经典检测器用多尺度滑动窗口在细尺度上找小目标；视频中，关键帧选择（\ref{sec:vision-videofound-keyframes}节）从冗余的帧中挑出代表帧。这些方法的共同思想是：不在一个尺度上看完全部内容，而是按需要在合适的尺度上看局部。

学习方法沿用了这一思想，并把“看哪里”交给模型或检索来决定。一种做法是改变输入方式，例如按原纵横比切出尽可能多的图像块，而不是把页面压成固定正方形；另一种做法是多次观察，先在全局低分辨率上推测目标所在区域，再裁剪放大；对长视频，则先检索候选片段，或把历史压缩成记忆再按问题回读。代价是计算量随图像块数或观察次数增加，裁剪还可能丢掉判断所需的上下文。本章的ScreenSeekeR（\ref{sec:vision-understanding-screenseeker}节）和Pix2Struct（\ref{sec:vision-understanding-pix2struct}节）分别代表多次局部观察和可变分辨率输入，\ref{sec:vision-understanding-evidence}节讨论页面与长视频中按需获取证据的方法。

\subsection{多步推理带来的证据组合困难}
\label{sec:vision-understanding-challenge-composition}

有些问题不能由一条证据回答。问“甲的2月金额比1月多多少”，要先在表格中找到甲行，再读出两个单元的数值，最后做减法；问“第二次敲门是否在讲话之后”，要先找出所有敲门和讲话的时间区间，再比较第二次敲门的起点与讲话的终点。困难在于，每一步都可能出错，而错误会传递：读错一个数字，后面的计算再准确，答案也是错的；而且只看最终答案，无法知道错误出在哪一步。

经典流水线把每一步写成独立模块，用规则连接：检测图表元素，按刻度把像素换算成数值，再交给确定的程序计算。每一步的输出都可以检查，规则执行的计算也不会出错；但规则只覆盖设计时考虑到的图表样式和问题类型，遇到新样式就失效。

学习方法在“显式步骤”与“端到端”之间给出不同的折中。一种做法让模型只预测结构化的中间决定，例如选择哪些表格单元、执行哪种聚合操作，计算仍由规则完成，中间决定可以检查，但操作集合是预先固定的；另一种做法让模型直接生成答案文字，表达不受限制，却不再暴露中间步骤。本章\ref{sec:vision-understanding-evidence-operations}节讨论怎样用外显的操作组合多条证据，\ref{sec:vision-understanding-answering}节比较从选择、抽取、单元聚合到自由生成四种答案形式在可追溯性与表达范围上的取舍。

表\ref{tab:vision-understanding-challenge-map}把五项挑战与经典方法、学习方法及本章回答的位置对应起来。

\begin{table*}[tbp]
  \centering
  \small
  \begin{tabularx}{\linewidth}{@{}p{24mm}XXX@{}}
    \toprule
    挑战 & 经典方法的处理 & 学习方法的改变 & 本章回答的位置\\
    \midrule
    多粒度对应 & 典型相关分析、动态时间规整、人工特征 & 对比对齐联合学习特征，注意力表达细粒度对应 & 匹配与检索、跨模态定位\\
    弱配对与噪声 & 小规模人工标注、固定类别检测器 & 规模、间隔目标、困难负例、带置信度的伪标注 & HowTo100M、声源定位、CRAFT\\
    语言先验 & 视觉与语言分步处理，晚期融合便于检查 & 更强的语言模型加重先验；反例数据、证据输出与专门评价 & 视觉描述、视觉问答、模态贡献的核验\\
    高分辨率与长输入 & 抗混叠与金字塔、多尺度滑窗、关键帧 & 可变分辨率输入、裁剪放大搜索、检索与记忆 & 界面控件定位、页面解析、证据的获取\\
    证据组合 & 模块化流水线与确定规则 & 单元选择加聚合、自由生成 & 证据的多步组合、答案的形成\\
    \bottomrule
  \end{tabularx}
  \caption{多模态理解的五项挑战与对应处理。经典方法与学习方法并非互斥：后者常常沿用前者的思想，例如裁剪放大搜索延续了多尺度观察，单元聚合延续了规则计算。}
  \label{tab:vision-understanding-challenge-map}
\end{table*}

这些挑战沿着证据从粗到细的顺序依次变得突出，本章的路线也由此确定。检索只需要整体粒度的对应，弱配对的大规模数据在这一粒度上最有用，所以先讨论召回与重排；落到区域、时间和控件时，细粒度对应与分辨率预算成为主要困难，这是定位一节的内容；票据和图表中的证据要先从像素变成字符和结构，才能参与计算，于是读取与解析放在描述和问答之前；描述和问答直接面对语言先验，又需要组合前面各步取得的证据，因此放在后面，并以分层核验收尾。

\section{跨模态匹配与检索}
\label{sec:vision-understanding-retrieval}

\term{跨模态检索}（cross-modal retrieval）指用一种模态的查询，在另一种模态的候选集合中找出相关内容，例如用一句话找图片，或用一句话找视频片段。输入是查询和可能有数十万乃至上亿项的候选库，输出是按分数排列的候选编号。难点在于两个要求相互拉扯：候选库很大，每次查询只能负担很便宜的比较；而“穿红衣服的人正在骑车”这样的请求要求判断对象、属性与关系是否由同一个人满足，这又需要细致的联合计算。

常用的解法是把检索拆成先后三步。第一步规定什么候选算相关，否则模型给出的高分无从验收。第二步召回：查询和候选各自编码成向量，候选向量可以预先计算并缓存，查询到来后只做内积，从全库中取出前$K$个候选。第三步重排：只对这$K$个候选做更昂贵的联合判断，重新排列顺序。召回决定正确结果有没有机会被看到，重排决定它们最终排在哪里；两步的错误要分开检查。

\subsection{检索相关性的规定}
\label{sec:vision-understanding-relevance}

模型给出的分数只是相似程度，是否“找对”要由事先规定的相关性来判断。这一小节要解决的问题是：查询中的哪些条件必须同时成立，候选的单位是什么，正例关系怎样标注。这些规定不依赖具体模型，却决定了召回与重排的评价口径。

“穿红衣服的人正在骑车”同时规定了人物、衣服颜色和动作关系。画面中有红衣行人，也有另一人在骑车，只满足了部分条件；把红色归到自行车上，同样改变了请求。可把完整条件分解为对同一个对象$i$的判断
\[
  \operatorname{person}(i)\land\operatorname{redClothes}(i)
  \land\exists j\,[\operatorname{bicycle}(j)\land\operatorname{ride}(i,j)]\text{。}
\]
这个表达式是相关性协议的说明，不表示双编码器会显式计算这些谓词。它的作用是防止标注和验收把“出现了所有关键词”当成“关系全部成立”。

候选单位也要固定。检索整段视频时，只要片中某处出现目标事件，就可能被标为相关；检索一个两秒片段，则要求该片段包含所需证据。一段文字可能描述多幅图片，一幅图片也可以有多种正确描述，因此应用数据的正例关系通常是多对多。训练划分还应避免同一视频的相邻片段或同一图的近重复版本分散到训练集和测试集。

为诊断组合关系，可以在保持对象词汇的条件下改变请求，如“人搬起箱子”和“箱子挡住人”，或把“戴帽子”改为“没有戴帽子”。相关性应按改变后的条件重新标注。部分相关候选可以保留等级，但不能在一个实验中按完整匹配计分，在另一个实验中又按同主题计分。

\subsection{候选内容的召回}
\label{sec:vision-understanding-candidate-matching}

召回要从$N$个候选中取出前$K$个，$K$通常远小于$N$。为了让每次查询的代价只与一次查询编码和$N$次内积相当，候选必须能够脱离查询单独编码。设候选表示组成矩阵$\bm V\in\mathbb R^{N\times d}$，查询表示为$\bm u\in\mathbb R^{d}$，召回就是取$\bm V\bm u$中最大的$K$项；库很大时还可以用近似最近邻索引代替逐项内积。难点在于：独立编码时查询看不到候选，候选向量必须事先保留所有可能被问到的信息。

双编码召回是\ref{sec:vision-alignment-contrastive}节对比对齐的直接应用。那里的结论是：用批内其他样本作负例、以交叉熵区分真实配对，两个编码器会把配对的表示拉到同一空间中的相近位置，内积因而可以作为相关程度。它的经典前身是\ref{sec:vision-alignment-cca}节的典型相关分析：二者都为两种模态各学一个到公共空间的映射，但典型相关分析作用在固定的人工特征上，映射是线性的，目标是使配对投影的相关系数最大，没有显式的负例；对比对齐则端到端训练深层编码器，并要求真实配对胜过批内其他组合。从融合位置看，两个编码器在计算过程中互不可见，只在最后一步的内积处相遇，相当于放弃了\ref{sec:vision-alignment-intermediate-fusion}节所说的中间交互。这正是双编码能缓存候选向量的原因，也是它难以判断关系绑定的原因。

本小节的三种做法都采用这种双编码结构，区别在于“候选”由什么观测单位构成，以及训练时用什么来源的配对告诉模型哪些内容相关。CLIP以单幅图像为单位，用网页图文对训练；CLIP4Clip把视频看成若干抽样帧，先逐帧编码再沿时间汇总，用人工撰写的视频描述训练；HowTo100M同样以视频片段为单位，但不依赖人工描述，而是把教学视频里自动转写的旁白当作弱配对。三者对同一个难点给出不同回答：图文配对干净但不含时间，人工视频描述含时间但昂贵，旁白规模大但与画面并不总是同步。

\subsubsection{CLIP}
\label{sec:vision-understanding-clip}

CLIP的结构和训练已在\ref{sec:vision-models-contrastive}节展开：图像编码器和文字编码器分别把输入投影到同一个$d$维空间并做$L_2$归一化，训练时在一批真实图文对上计算余弦相似度矩阵，用行、列两个方向的交叉熵把对角线上的配对与批内其他组合区分开。这里用到的结论只有一条：训练后，内积越大表示图文越可能来自同一配对，但这个分数衡量的是批内相对选择，不是对每个条件逐项核对的结果。

用于检索时，先按模型约定裁剪、缩放和归一化图库中的每幅图，编码为$\bm v_i$并存成矩阵$\bm V$。文字查询$q$编码为$\bm u(q)$后，以
\[
  s_i=\bm u(q)^\top\bm v_i
\]
排序并返回图像编号。图像检索文字则反过来缓存候选描述$\bm u_j$，编码查询图像$\bm v$，比较$\bm v^\top\bm u_j$。两种方向共用相似度形式，但候选集合、正例标注和评价分母不同。查询文字可以直接使用用户原句，也可以套用“一张……的照片”这样的模板；模板改变文字编码器看到的上下文，应在比较实验中固定。

假设查询向量为$(1,0)^\top$，三个图像向量为$(0.8,0.6)^\top$、$(0.6,0.8)^\top$和$(0,1)^\top$，均已归一化。教学排序为甲、乙、丙，分数依次为0.8、0.6、0。若人工核对发现乙才满足“红衣人物骑车”的全部关系，甲仅有相关对象，那么模型把部分相关候选排到了完整正例之前。将分数转成softmax并不会补上这个关系判断。这是CLIP召回最典型的失败：整图向量能记住“有红色、有人、有自行车”，却不一定记住红色属于谁。

是否微调取决于目标图库与预训练数据的差距。直接使用原始权重时，不需要目标任务的标注；若用目标领域的图文对微调，通常沿用式\eqref{eq:vision-clip-loss}的双向目标，同时更新两个编码器和投影。此时要注意多对多的正例关系：一个查询的全部已知正例应同时参与目标，或在采样时排除已知正例之间的“假负例”。召回数$K$是主要超参数：$K$越大，后续重排越有机会看到正确候选，重排代价也随$K$线性增加。

\subsubsection{CLIP4Clip}
\label{sec:vision-understanding-clip4clip}

CLIP4Clip的帧编码、时间聚合和训练过程已在\ref{sec:vision-video-clip4clip}节展开。它从CLIP权重出发，把按固定规则抽取的视频帧逐帧编码，再用均值或带时间位置的序列模块汇总成一个视频向量，用视频—描述对上的双向交叉熵微调编码器和聚合器。对检索而言，它把CLIP的“单幅图像”候选换成了“抽样帧的汇总”，因而能比较文字与一段时间内的内容。

使用时，按训练时的帧率和帧数预先编码候选视频或片段，收到查询后与查询向量比较，输出视频编号及分数。抽帧数、片段长度和聚合方式都是需要随结果报告的超参数：抽帧过稀会漏掉短暂动作，片段过长会让目标事件在均值中被稀释。图像已经提供的对象证据与视频增加的时间证据，应分别检查。

对同一图库，依次输入“箱子”“人搬起箱子”“穿红衣服的人搬起箱子”，请求增加了动作和参与者条件。排序改变可能来自真正的时间关系，也可能只来自某帧出现了红衣人物。可以定位命中片段，查看谁接触箱子、箱子是否离开支撑面，以及动作顺序。若帧聚合主要保留整段共同出现的内容，交换“人”和“箱子”的语义角色未必产生足够大的得分变化。

命中整段视频仅完成了候选选择，动作起止仍需时间定位。抽样漏掉搬起瞬间时，下游文字模型可能根据前后场景作合理猜测，但不能把该瞬间写成已经观察到的事实。

\subsubsection{HowTo100M}
\label{sec:vision-understanding-howto}

CLIP4Clip的训练依赖人工为视频片段撰写描述，这类数据的规模受标注成本限制，覆盖的动作种类也有限。教学视频提供了另一种联系：讲解者一边操作一边说明，语音识别生成的字幕行自带起止时间。HowTo100M利用这一点，从约122万段带旁白的教学视频中切出约1.36亿个“片段—旁白”记录，覆盖两万多种任务，不需要人工撰写描述\scite{vision-howto2019}。

这种配对是\emph{弱配对}：旁白可能提前说“接下来揉面”，也可能省略手上正在做的动作，字幕的时间区间与动作并不严格对齐。HowTo100M的主要想法是用规模抵消单条记录的噪声，并在训练中专门加入同一视频内部的错配，让模型必须区分同一厨房里的不同步骤，而不能只凭背景判断“这是做饭视频”。

模型由两个独立的编码分支组成，具体配置见表\ref{tab:vision-understanding-howto-config}。视频一侧使用预先提取、训练时不更新的特征：以ImageNet上训练的ResNet-152每秒取一帧提取外观特征，以Kinetics上训练的16帧ResNeXt-101三维卷积网络每秒约取1.5个特征描述运动；两种特征各自沿片段时间做最大池化，再拼接为一个4096维向量$\bm x_{\mathrm v}$。时间最大池化保留片段内每一维的最强响应，代价是丢掉了这些响应出现的先后顺序。文字一侧以在GoogleNews语料上预训练的word2vec词向量表示旁白中的词，再经一个浅层一维卷积网络得到4096维的$\bm x_{\mathrm t}$；论文没有给出这个卷积网络的层数和卷积核大小。两侧各接一个\term{门控嵌入单元}（gated embedding unit）：它先做线性变换，再让输出的每一维乘上一个由自身计算的0到1之间的系数，按当前内容调节各维的贡献。论文把它写成线性层后接上下文门控。设输入为$\bm x$，有
\begin{equation}
  \bm z=\bm W_1\bm x+\bm b_1,\qquad
  g(\bm x)=\bm z\odot\sigma(\bm W_2\bm z+\bm b_2)\text{，}
  \label{eq:vision-howto-gate}
\end{equation}
其中$\sigma$为逐元素sigmoid，$\odot$为逐元素乘积。原设置中输入维数与输出维数$d$都是4096，因此$\bm W_1,\bm W_2\in\mathbb R^{4096\times4096}$；两侧参数不共享。输出经$L_2$归一化后，以余弦相似度$s_{i,j}$比较第$i$段视频和第$j$条旁白。门控与\ref{sec:vision-alignment-attention}节的注意力都在计算权重，区别在于门控的权重作用在特征维度上，由同一模态的内容决定，并不在两个模态的元素之间建立对应；跨模态的比较仍然只发生在最后的余弦相似度。

\begin{table}[htbp]
  \centering
  \small
  \begin{tabularx}{\linewidth}{@{}p{16mm}XXp{14mm}@{}}
    \toprule
    组件 & 输入 & 结构 & 输出\\
    \midrule
    外观特征 & 每秒1帧 & ImageNet预训练ResNet-152，冻结 & 沿时间最大池化\\
    运动特征 & 16帧片段，每秒约1.5个 & Kinetics预训练ResNeXt-101，冻结 & 沿时间最大池化\\
    视频向量 & 两种特征 & 拼接 & 4096维\\
    文字向量 & 旁白词序列 & word2vec词向量，浅层一维卷积（规模论文未给出） & 4096维\\
    门控嵌入 & 4096维 & 式\eqref{eq:vision-howto-gate}，两侧各一个 & 4096维\\
    \bottomrule
  \end{tabularx}
  \caption{HowTo100M文本—视频嵌入模型的组件配置，依据原论文。两种视觉特征各自的维数论文未分别给出，只给出拼接后的4096维。}
  \label{tab:vision-understanding-howto-config}
\end{table}

训练希望真实配对得分比错配至少高出间隔$\delta>0$。对一个配对$i$及负例$j$，两个方向的损失为
\begin{equation}
  L_{i,j}=[\delta+s_{i,j}-s_{i,i}]_+
          +[\delta+s_{j,i}-s_{i,i}]_+\text{，}
  \qquad [a]_+=\max(0,a)\text{。}
  \label{eq:vision-howto-ranking}
\end{equation}
第一项替换旁白，第二项替换视频。各项已经满足间隔时不再产生惩罚，这与交叉熵始终推动正例概率趋近1不同。原设置取$\delta=0.1$。原实现用批次内的其他记录形成负例，其中一半来自同一段原视频，另一半来自其他视频；否则数量更多的跨视频负例会让模型只学会区分场景。每个批次取32段视频、每段64个片段—旁白对，用学习率$10^{-4}$的Adam训练；预提取的视觉特征保持不变。

以$\delta=0.1$、$s_{i,i}=0.7$、$s_{i,j}=0.65$、$s_{j,i}=0.4$为教学数值，损失为$0.05+0=0.05$。若$j$是同一厨房里的另一条旁白，第一项要求模型把两个步骤分开，只认出厨房背景还不够。

训练后的两个映射可以直接用于新数据：对查询编码，按片段分数排序，再把片段时间还原到原视频。原论文既在YouCook2、CrossTask等教学视频数据上不经微调直接检索和定位步骤，也在MSR-VTT等人工描述数据上以HowTo100M预训练为起点再微调。排序目标鼓励区别配对，却不会自动消除旁白和实际动作之间的错位；\ref{sec:vision-alignment-contrastive}节介绍的MIL-NCE就是针对这一点的后续改进，它把时间上相邻的若干段旁白都当作候选正例；查询内容若是教学视频很少涉及的场景，例如街头监控，弱配对数据的覆盖也帮不上忙。

三种召回方式的差别可以放在同一张表中对照（表\ref{tab:vision-understanding-recall-compare}）。从CLIP到CLIP4Clip，改变的是候选单位：把单幅图换成多帧汇总，换来了对时间内容的比较能力，代价是抽帧与聚合会稀释短暂事件。从CLIP4Clip到HowTo100M，改变的是配对来源：用自动旁白代替人工描述，换来了规模，代价是配对噪声和领域局限。三者都只能给出整体相似度，关系是否由同一对象满足，要留给重排处理。

\begin{table}[htbp]
  \centering
  \small
  \begin{tabularx}{\linewidth}{@{}p{14mm}XXX@{}}
    \toprule
    比较项 & CLIP & CLIP4Clip & HowTo100M\\
    \midrule
    候选单位 & 单幅图像 & 抽样帧汇总的片段 & 字幕时间区间对应的片段\\
    配对来源 & 网页图文对 & 人工视频描述 & 自动转写的旁白\\
    视觉编码 & 图像编码器 & CLIP图像编码器逐帧编码后聚合 & 预提取的二维与三维卷积特征\\
    训练目标 & 双向交叉熵 & 双向交叉熵 & 双向间隔排序\\
    典型失败 & 属性与对象错绑 & 短暂动作被稀释 & 旁白与动作错位\\
    \bottomrule
  \end{tabularx}
  \caption{三种召回方式的对照。三者都预先编码候选、以内积排序；表中各行说明它们在候选单位、配对监督和失败情形上的差别。}
  \label{tab:vision-understanding-recall-compare}
\end{table}

\subsection{候选内容的重排}
\label{sec:vision-understanding-reranking}

召回给出$K$个候选，其中可能有同主题而关系不同的内容，例如红衣行人和蓝衣骑车者同框的图片。重排的输入是查询和这$K$个候选，输出是新的顺序，难点是让查询直接作用于候选的局部内容，判断条件是否由同一对象满足。因为只处理$K$项，可以负担每项一次完整的联合计算。

按\ref{chap:vision-multimodal-alignment}章的融合位置来看，重排是在召回放弃的地方补上中间融合（\ref{sec:vision-alignment-intermediate-fusion}节）：查询的词元与候选的视觉单元在网络内部通过交叉注意力相互读取，每个词可以把注意力放到它所描述的区域上，关系绑定才有机会在计算中体现。中间融合的代价是表示不能离线缓存，每个查询—候选对都要重新计算一次，所以只能用于召回之后的少量候选。

重排有两条路线，区别在于关系判断是否由一个学到的分数隐式完成。基于配对判定的重排把查询和候选一起送入交叉注意力网络，输出一个“是否匹配”的概率，BLIP-2的图文匹配头是这种做法。基于定位与读取的重排把查询拆成可检查的条件，调用定位和读取接口取得区域证据，再按条件满足情况排序。前者一次计算、无需拆解查询，后者每个条件都有可核对的依据，但依赖拆解规则和下游接口。两条路线都无法找回召回阶段已经漏掉的正例。

\subsubsection{基于配对判定的重排}
\label{sec:vision-understanding-blip-retrieval}

BLIP-2的Q-Former已在\ref{sec:vision-models-blip2}节展开：32个可学习查询通过交叉注意力读取图像编码器输出，形成固定数量的图像查询表示；第一阶段预训练用同一组参数、三种注意力掩码学习三个目标。其中图文对比目标让查询与文字互不可见，得到可以独立缓存的表示；图文匹配目标让查询与文字双向交互，由二分类头判断图文是否配对。检索正好用到这两种模式：前者负责召回，后者负责重排。

使用时，先用图像查询表示与文字表示的相似度召回有限候选。重排阶段把一个候选文字与图像查询一同送入Q-Former，并采用允许二者互看的注意力掩码；每个查询输出经二分类头得到匹配logit，取平均后得到配对概率$p_1$，按$p_1$重新排列候选。这条路径只用到Q-Former，不需要大型语言模型逐词生成。由于每个候选都需要一次联合前向计算，文字检索图像时代价是$K$次Q-Former计算，而召回只需一次查询编码。

原论文的检索适配共同微调图像编码器和Q-Former，并保留对比、匹配及图像条件文字生成三个目标；这不同于第一阶段冻结图像编码器的预训练。其推理设置先取128个候选，再按配对匹配分数重排。冻结范围和候选数都应随所用实验设置报告，不能从预训练说明中直接沿用。

重排能够让“同一画面里的红衣人物是否就是骑车者”影响联合计算，但仍需关系数据检验：匹配头的训练负例多数来自不同图片，未必覆盖属性在同图内换绑的情形。若图库有三个完整正例，召回集仅含其中两个，那么无论重排多准确，完整正例的召回率最多为$2/3$。这个上限来自候选已经被删除，与后续模型有多少参数无关。

\subsubsection{基于定位与读取的重排}
\label{sec:vision-understanding-constraint-check}

配对概率是一个数，它不说明哪个条件成立、哪个不成立。对需要严格满足条件的查询，可以使用一个教学组合流程：从请求中列出需要核对的对象、属性、关系、计数或文字条件$c_1,\ldots,c_m$；对高分候选调用后文的定位和读取接口取得区域证据；再对每个条件记录状态$z_k\in\{\text{满足},\text{不满足},\text{证据不足}\}$。排序时先比较“不满足”的条件数，少者在前；再比较“证据不足”的条件数；仍相同时才用召回分数。这条规则不新增模型训练，也不属于CLIP或BLIP-2的原始输出。

例如，甲图得分0.8，但定位后发现红衣人物在步行，骑车者穿蓝衣，记录为“动作绑定不满足”，排到后面。乙图得分0.6，相关人物和自行车位置、接触关系均可见，三项条件都满足，排在前面。若车身遮住腿部，图像又无法区分骑行和推行，应把这一条记为证据不足，而不是用较高相似度补足不可见部分。

这一流程的超参数是参与核验的候选数和各接口的阈值，失败主要来自两处：拆解规则漏掉了条件，或者定位与读取接口本身出错。结果可以同时返回候选编号、匹配分数、目标区域及各条件状态，使排序错误能被追溯到具体环节。

两种重排可以按同一查询和同一召回集比较（表\ref{tab:vision-understanding-rerank-compare}）。基于定位与读取的重排把配对判定的单个分数拆成逐条件的证据记录，换来了可解释性和对“证据不足”的显式表达，代价是更多的接口调用，以及对查询拆解规则的依赖；配对判定不需要拆解，但它犯错时只能从数据和模型训练上找原因。

\begin{table}[htbp]
  \centering
  \small
  \begin{tabularx}{\linewidth}{@{}p{16mm}XX@{}}
    \toprule
    比较项 & 基于配对判定的重排 & 基于定位与读取的重排\\
    \midrule
    每个候选的计算 & 一次联合前向计算 & 每个条件一次定位或读取\\
    输出 & 配对概率 & 各条件的状态与区域证据\\
    训练需求 & 需要图文匹配训练 & 复用已有接口，不新增训练\\
    能否表达证据不足 & 不能，只有一个分数 & 能，单独记录\\
    主要失败 & 匹配头未见过的换绑情形 & 拆解遗漏条件、接口误差\\
    \bottomrule
  \end{tabularx}
  \caption{两种重排的对照。两者都在同一召回集上工作，都不能找回召回阶段漏掉的正例。}
  \label{tab:vision-understanding-rerank-compare}
\end{table}

\section{跨模态定位}
\label{sec:vision-understanding-grounding}
\label{sec:vision-understanding-grounding-query}

检索给出“哪幅图相关”，而“右边拿杯子的人”需要落到画面中的某个对象。\term{视觉定位}（visual grounding）把文字或其他查询对应到视觉对象、空间区域或时间位置。输入是一份观测和一个查询，输出是框、掩码、点或时间区间；输出的位置必须仍属于原始观测，才能用于复核、裁剪或后续操作。

定位之前要先解析查询，确定目标是什么、靠什么参照来区分。在“右边拿杯子的人”中，“人”是目标类别，“拿杯子”是人与杯子的关系，“右边”还依赖参照系。若两个人都在拿杯子，“右边”可能指图像右侧，也可能指说话者右手方向。系统应把默认参照系写入协议；本章未另说明时，图像横轴向右、纵轴向下。多图输入还需指定“图2中的人物”，页面输入需带页码，视频输入需带时刻或区间。“他旁边的杯子”依赖此前已经确定的“他”，不能仅按最近出现的名词替换。无法唯一确定时，可以返回候选集合或请求补充参照条件；查询目标根本不存在，也应允许输出空结果。

输出形式也随任务规定。框表示外接范围，掩码描述像素归属，点击点必须位于可操作控件内，区间描述事件持续时间。它们的几何关系可以辅助转换，但不能仅凭输出了四个数就认定完成框定位，也不能把框中心总是可点击当作普遍条件。

查询与被定位的对象可以属于不同模态，本节按定位对象展开四种情形。文字短语定位到图像区域时，监督最充分，可以端到端训练检测式模型。声音来源定位到图像区域时，没有框标注，只能从音画同时出现的关系中自监督学习。音画事件定位到时间区间时，要比较两条时间序列。界面控件按指令定位到点击位置时，截图分辨率高、目标极小，难点转向观测的分辨率预算。四种情形的共同问题是：定位需要保留空间或时间分辨率，把整图压成一个向量再与整句比较，无法知道哪个位置支持哪段表达。

\subsection{文字短语的区域定位}
\label{sec:vision-understanding-grounding-match}

输入是一幅图像和一段文字，文字可以是类别名、短语，也可以是\term{指代表达}（referring expression），即用属性和关系唯一指明某个对象的描述，如“右侧戴帽子的人”。输出是若干框，每个框附带它对应的文字位置。难点有两个：文字词汇是开放的，不能事先列出类别表；同一句话可能提到多个对象，模型要知道每个框对应句中哪几个词。

两种代表模型都沿用DETR式集合预测（\ref{sec:vision-recognition-detr}节）：一组对象查询经解码得到框，训练时用二分匹配把预测与标注一一对应。它们的差别在于文字在何处、以何种方式进入检测计算。MDETR在编码器入口把图像特征与文字词元拼成一个序列，此后每一层都联合计算，并让每个对象查询预测它指向哪些词元位置。Grounding DINO保留两条分支，在特征增强阶段交换信息，再用文字相关性挑选初始化对象查询，并以更大规模的检测与定位数据训练，面向开放类别检测。

用\ref{chap:vision-multimodal-alignment}章的术语说，两者都属于中间融合（\ref{sec:vision-alignment-intermediate-fusion}节）：各模态先由自己的编码器处理，再在网络内部交互。MDETR的交互方式是把两种特征拼接成一个序列后做自注意力，形式上接近\ref{sec:vision-alignment-early-fusion}节的特征拼接，但拼接发生在各自编码之后，而且拼接后还有多层注意力，每个词元都能对每个视觉单元分配权重；Grounding DINO则用两个方向的交叉注意力交换信息，两条分支始终保持各自的表示。这两种结构都让词与区域之间的注意力权重成为\ref{sec:vision-alignment-attention}节所说的软对齐，而短语—框标注直接监督了这种对齐的结果。

\subsubsection{MDETR}
\label{sec:vision-understanding-mdetr}

在MDETR之前，常见的文字定位系统先用一个固定类别的检测器提出候选框，再让文字与这些框匹配。检测器若没有为“行李架”这类不在类别表中的对象提出候选，后续匹配无从补救；检测器训练时也不知道当前句子要找什么。MDETR的主要想法是让文字从一开始就调制检测：图像与文字在同一个Transformer编码器中交换信息，对象查询直接读取融合后的表示产生框，因而候选框的生成本身就以文字为条件\scite{vision-mdetr2021}。

模型的数据流如下。图像经卷积骨干得到$h\times w$的空间特征网格，用$1\times1$卷积降到$d$维后展平为$hw$个视觉单元，并加上二维位置编码以保留空间信息。文字经预训练的RoBERTa编码为$l$个词元表示，再线性投影到同样的$d$维。两组单元沿序列方向拼接成$(hw+l)\times d$的矩阵，送入联合Transformer编码器，自注意力在视觉单元和词元之间不加区分地交换信息。解码器有$N$个可学习的对象查询，读取编码器输出后，每个查询产生一个$d$维表示$\bm o_i$，再接两个预测头：框头输出归一化的中心与宽高，词元头输出一个在全部词元位置及额外“无对象”位置上的概率分布$p_{i,\cdot}$。词元位置的上限为$L=256$。

表\ref{tab:vision-understanding-mdetr-config}列出各组件的规模。论文正文给出了文字编码器、骨干、查询数、温度与训练轮数；联合编码器与解码器的层数、宽度和头数，正文只说沿用DETR，具体数值取自官方实现的默认配置，表中分别注明。以ResNet-101为例作一个教学计算：骨干最后一级特征图的步长为32，若输入图像为$800\times1333$像素，特征图约为$25\times42$，即约1050个视觉单元；加上至多256个词元，联合编码器处理的序列长度约为1300。论文正文未给出预训练时的图像尺寸，这里的尺寸只用于说明序列长度的量级。

\begin{table}[htbp]
  \centering
  \small
  \begin{tabularx}{\linewidth}{@{}p{18mm}XX@{}}
    \toprule
    组件 & 配置 & 来源\\
    \midrule
    图像骨干 & ResNet-101（ImageNet预训练，冻结批归一化），或EfficientNet-B3（约1200万参数）、B5（约3000万参数） & 论文\\
    文字编码器 & RoBERTa-base：12层，宽度768，12头 & 论文\\
    投影与位置 & 两侧线性投影到$d=256$；图像加二维位置编码 & 投影见论文，$d$见官方实现\\
    联合编码器 & 6层，宽度256，8头，前馈维数2048，丢弃率0.1 & 官方实现\\
    解码器 & 6层，100个对象查询，宽度与头数同编码器 & 查询数见论文，层数见官方实现\\
    预测头 & 框：中心与宽高；词元：256个位置加“无对象” & 论文\\
    对比对齐 & 投影到64维，温度$\tau=0.07$ & 维数见官方实现，温度见论文\\
    预训练 & 约130万对图文、约20万幅图像；40轮，等效批大小64 & 论文\\
    参数量 & 论文未给出 & —\\
    \bottomrule
  \end{tabularx}
  \caption{MDETR的组件配置。官方实现指作者发布代码中的默认参数；论文正文未给出的数值不作推测。}
  \label{tab:vision-understanding-mdetr-config}
\end{table}

联合编码器沿用DETR的Transformer块，默认在残差相加之后做层归一化。设输入序列为$\bm Z\in\mathbb R^{(hw+l)\times256}$，位置编码为$\bm P$，单个编码器块计算
\[
  \bm Z'=\operatorname{LN}\bigl(\bm Z+\operatorname{MSA}(\bm Z+\bm P,\,\bm Z+\bm P,\,\bm Z)\bigr),\qquad
  \bm Z''=\operatorname{LN}\bigl(\bm Z'+\operatorname{FFN}(\bm Z')\bigr)\text{，}
\]
其中$\operatorname{MSA}(\cdot,\cdot,\cdot)$的三个参数依次为查询、键和值，位置编码只加在查询和键上；$\operatorname{FFN}$为$256\to2048\to256$的两层全连接网络。解码器块在自注意力和前馈网络之间多一层交叉注意力，以对象查询为查询、以编码器输出为键和值。对象查询读到的编码器输出已经混合了文字信息，所以框的产生本身是以文字为条件的。

训练数据是带有“短语—框”对应的图文记录，原论文从Flickr30k、COCO的指代表达数据和Visual Genome中汇集约130万对图文。集合匹配沿用DETR，在预测框与真实框之间寻找一对一分配，再计算框的$L_1$损失和广义IoU损失。文字对应则允许多对多：一个短语可以指两个人，同一个人也可能在句中被不同短语指代。

第一个文字目标是\term{软词元预测}（soft token prediction）：不要求模型预测类别名，而是预测当前对象在句中对应的词元位置分布。设一个已匹配对象对应的词元位置集合为$T_i^+$，目标把概率质量均匀分给这些位置，其交叉熵为
\begin{equation}
  L_{\mathrm{token},i}
  =-\frac{1}{|T_i^+|}\sum_{j\in T_i^+}\log p_{i,j}\text{。}
  \label{eq:vision-mdetr-token}
\end{equation}
例如“红衣／人物”占两个词元位置，目标分别为$1/2$，不是要求模型任选一个位置得到概率1。未匹配的对象查询则把概率集中到“无对象”位置。实际对应建立在分词后的位置上，不能用中文字符数代替模型词元数。

软词元预测约束的是输出位置，并没有直接约束表示之间的距离。MDETR另外把对象表示$\bm o_i$与编码器输出的词元表示分别投影到低维并归一化，得到$\hat{\bm o}_i$与$\hat{\bm t}_j$，计算对象到文字的对比项
\begin{equation}
  L_{\mathrm{o\to t},i}
  =-\frac{1}{|T_i^+|}\sum_{j\in T_i^+}
    \log\frac{\exp(\hat{\bm o}_i^\top\hat{\bm t}_j/\tau)}
    {\sum_k\exp(\hat{\bm o}_i^\top\hat{\bm t}_k/\tau)}\text{。}
  \label{eq:vision-mdetr-contrast}
\end{equation}
文字到对象方向按相同方式对齐，再取两个方向的平均；没有正对应的项不参与构造。原论文按文献惯例取温度$\tau=0.07$。总目标为框的$L_1$与广义IoU损失、软词元损失和对比对齐损失的加权和，官方实现默认的权重依次为5、2、1、1，共同更新骨干、文字编码器、联合编码器和解码器；预训练时骨干与Transformer的学习率按论文附录分别为$10^{-4}$与$10^{-5}$，35轮后降为十分之一；文字编码器的学习率先线性预热到$5\times10^{-5}$，再线性衰减。两个文字目标使用相同标注，但作用不同：一个规定“输出指向哪里”，一个规定“表示彼此靠近”。

预训练后，MDETR可以直接用于短语定位：对句中每个目标短语，把各对象查询在该短语词元上的概率相加，按这个分数筛选框并恢复原图坐标。用于指代表达理解、GQA式问答等下游任务时，原论文在相应数据上继续微调。若图中有红衣人物甲和蓝衣人物乙，输入“红衣人物拿着杯子，蓝衣人物站在旁边”，模型可以同时输出两个人物框及对应短语位置；若任务只要求拿杯者，还需排除作为参照出现的乙。框能支持区域检查，像素掩码需要另加分割路径；复杂关系能否正确消歧，仍取决于训练与测试数据是否覆盖相应关系。

\subsubsection{Grounding DINO}
\label{sec:vision-understanding-grounding-dino}

Grounding DINO的跨模态特征增强、语言引导的查询选择和跨模态解码已经在\ref{sec:vision-recognition-grounding-dino}节展开。与MDETR相比，它保留图像与文字两条分支，在多个阶段做双向交叉注意力，并按候选位置与词元的相似度挑选初始化对象查询；输出同样是框和框对各词元的匹配得分。这里讨论把它用于文字短语定位时，输入怎样组织、输出怎样解码。

输入文字有两种组织方式。类别词列表按所用接口的分隔约定写成“人物 . 帽子 .”这样的串，每个类别是一段独立的短语；完整指代表达则作为一句话输入，保留属性和关系，不随意拆成互不相关的词。把“人物”“戴帽子的人”“右侧戴帽子的人”分别作为查询，可以观察属性和参照条件对结果的影响。使用已发布检查点时，通常不为每次查询更新权重，并保留其分词、图像尺度和阈值协议；原论文在指代表达任务上另做了微调，两种用法的结果不应混为一谈。

解码时，先取每个框在全部词元上的最高得分，以框阈值筛出候选；再对每个保留的框，以文字阈值选出得分足够高的词元，拼接恢复短语。两道阈值作用不同：降低框阈值会增加候选区域，降低文字阈值会扩大附着在框上的文字范围。二者都不能把“右侧”一定指向当前目标作为逻辑保证，返回的短语也未必是连续、完整的语法单位。对“右侧戴帽子的人”，应核对帽子属于哪个人物、“右侧”的参照是否一致，以及是否存在第二个满足条件的候选。需要概念对应的实例掩码时，可以使用\ref{sec:vision-recognition-sam3}节的SAM 3接口，分别检查概念存在与实例范围；框和掩码的组合增加了成果类型，并没有免除关系核验。

两者在同一任务上的差别集中在文字进入检测的位置。MDETR从编码器入口开始联合计算，词元与视觉单元在每一层都能直接交互，适合需要逐词对应的短语定位，但每次查询都要重新运行整个编码器，训练数据也局限于带短语对应的图文集。Grounding DINO把文字作用分散到特征增强、查询选择和解码三处，并利用大规模检测数据，换来更强的开放类别检测能力；代价是长句中的关系仍主要靠词元匹配得分表达，阈值解码可能截断短语。两者都输出框而非掩码，也都需要另外检查参照关系。

\subsection{声音来源的区域定位}
\label{sec:vision-understanding-sound-source}

一帧画面中同时出现说话者和旁观者，问题是“声音来自画面中哪里”。输入是一帧图像和与之同时的一段音频，输出是一张与图像网格对齐的热图，高值区域表示\term{声源定位}（sound source localization）的结果。难点在于几乎没有带框的声源标注，能大量获得的只有“这段声音与这段画面来自同一视频”的配对关系。因此这一任务通常采用自监督路线：让模型从音画是否同时出现中学习，再从学到的空间相似度中读出位置。路线之间的区别在于负例怎样构造。

经典方法把这个问题看成像素与声音之间的相关分析。发声物体的运动与声音的强弱同步变化，例如说话者的嘴唇开合与语音能量起伏一致，因此可以在一段视频上，统计每个像素的亮度变化与音频特征随时间的相关程度，相关最强的位置就是声源。Pixels that Sound把\ref{sec:vision-alignment-cca}节的典型相关分析用于这一问题，并要求视觉一侧的投影稀疏，使高权重集中在少数像素上\scite{vision-pixelssound2005}。这类方法不需要标注，但依赖声音与像素变化的同步，对静止的发声物体或背景中同样在动的物体都会失效。学习方法把像素换成深层特征，把时间相关换成音画是否同时出现的对比判别，从大量视频中学到“什么样子的东西发出什么样的声音”，因而能定位单帧中静止的发声物体。

此前常见的学习做法把整幅图与声音作对比：同一视频的图和声为正例，不同视频的为负例。这种训练只要求区分“有人说话的画面”和“另一种场景”，并不要求模型指出画面中哪个位置在发声，定位能力只是附带出现。Localizing Visual Sounds the Hard Way的主要想法是把同一幅图中的背景区域也当作负例：既然画面里只有一部分在发声，模型就应让声音与这部分的相似度高于与同图其他部分的相似度\scite{vision-visual-sounds2021}。

模型有两条编码分支。图像经卷积网络保留空间网格，得到各位置的特征$\bm f_{u,v}\in\mathbb R^d$；音频先转为频谱图，再经另一卷积网络并全局池化为单个向量$\bm a\in\mathbb R^d$。对非零的$\bm a$与$\bm f_{u,v}$，逐位置计算余弦相似度
\begin{equation}
  S_{u,v}=\frac{\bm a^\top\bm f_{u,v}}
                 {\lVert\bm a\rVert_2\lVert\bm f_{u,v}\rVert_2}\text{，}
  \label{eq:vision-sound-map}
\end{equation}
得到一张$h\times w$的热图。声音编码器输出整段表示，因此热图描述的是这段声音与当前图像的空间联系，不是每个音频采样点的位置。这个计算与\ref{sec:vision-alignment-attention}节的注意力打分形式相同：音频向量相当于查询，各空间位置的视觉特征相当于键，只是这里不做softmax归一化，而直接把余弦相似度当作定位分数。

表\ref{tab:vision-understanding-lvs-config}给出原论文的配置。两条分支都用ResNet-18：图像分支保留空间网格，输入$224\times224\times3$时输出$14\times14\times512$的特征；音频分支把3秒音频转成$257\times300$的幅度谱图，经ResNet-18得到$17\times13\times512$的特征图，再做最大池化得到512维向量。由于两侧通道数相同，式\eqref{eq:vision-sound-map}可以直接在512维上计算。

\begin{table}[htbp]
  \centering
  \small
  \begin{tabularx}{\linewidth}{@{}p{14mm}XXX@{}}
    \toprule
    分支 & 输入 & 主干 & 输出\\
    \midrule
    图像 & $224\times224\times3$，不裁剪 & ResNet-18，ImageNet预训练 & $14\times14\times512$特征网格\\
    音频 & 3秒片段的$257\times300$幅度谱图 & ResNet-18 & $17\times13\times512$，最大池化为512维\\
    热图 & 两侧输出 & 逐位置余弦相似度 & $14\times14$\\
    \bottomrule
  \end{tabularx}
  \caption{Localizing Visual Sounds the Hard Way的结构配置。论文未给出音频采样率；测试时使用完整长度的谱图。}
  \label{tab:vision-understanding-lvs-config}
\end{table}

训练只有音画配对时，真实声源掩码不可用。方法用当前热图本身构造较可信的正区域和同图负区域，用两道阈值$\epsilon_+>\epsilon_-$和温度$\tau$得到软掩码：
\[
  m^+_{u,v}=\sigma\!\left(\frac{S_{u,v}-\epsilon_+}{\tau}\right),
  \qquad
  m^-_{u,v}=1-\sigma\!\left(\frac{S_{u,v}-\epsilon_-}{\tau}\right)\text{。}
\]
高分位置获得较大正权重，低分位置获得较大负权重，中间区域对两侧的贡献较小；由于使用sigmoid，中间带不是严格为零。对批次中的第$i$个视频，正响应$P_i$是以$m^+$加权的平均相似度；负响应$N_i$由两部分相加：以$m^-$加权的同图背景平均相似度，以及该声音与批内其他视频图像各位置相似度之和除以网格大小$hw$。训练目标为
\begin{equation}
  L=-\frac1k\sum_{i=1}^{k}\log\frac{\exp(P_i)}{\exp(P_i)+\exp(N_i)}\text{，}
  \label{eq:vision-sound-loss}
\end{equation}
其中$k$为批次大小。它提高正响应相对于两类负响应的得分，同时更新图像和音频编码器。原设置取$\epsilon_+=0.65$、$\epsilon_-=0.4$、$\tau=0.03$：温度很小，sigmoid接近阶跃函数，相似度高于0.65的位置几乎全部算作正区域，低于0.4的几乎全部算作负区域。训练使用学习率$10^{-4}$的Adam，批大小256，数据取自Flickr SoundNet和VGGSound中1万或14.4万对图像与音频；原论文还为评价构建了VGG-SS测试集，表中列出5158段视频、220个类别，正文概称约五千段、两百多类。

使用时，把热图上采样到原图尺寸，按所定阈值得到区域，再与标注的可见声源区域比较。同图困难负例使训练必须区分发声者和视觉上同样显著的旁观者。不过，正负区域来自模型自己的热图，这种自举仍可能把错误的高分区域当成正例并不断强化。若声音来自画外，热图峰值仍会指向某个可见对象；处理这种情况需要另设“无可见声源”判定和相应数据，原热图接口不能给出保证。

\subsection{音画事件的时间定位}
\label{sec:vision-understanding-av-time}

空间声源回答“画面哪里”，时间定位回答“另一模态的哪一段与这段内容同步”。输入是一段音频查询和一段较长的视频（或反过来），输出是视频中与之对应的时间区间。难点在于两条序列的单元要能直接比较，而训练时只知道哪些音画片段同时发生。Audio-Visual Event Localization一文为此构建了AVE数据集，其中的视频包含同时可见又可听的事件，并提出从音频定位视觉片段及其反方向的跨模态定位任务\scite{vision-ave2018}。

该文的主要想法是学习一个音画距离：同步的音频片段与视觉片段在共同空间中距离小，错配的距离大，定位时沿时间滑动寻找累计距离最小的窗口。AVE数据集含4143段10秒视频，覆盖28类事件；每段视频按1秒切成10个单元，每个单元标注事件类别或背景。

模型称为音画距离学习网络，结构见表\ref{tab:vision-understanding-avdln-config}。视觉一侧，每个1秒单元取16帧RGB图像，用ImageNet上预训练的VGG-19提取pool5特征图，再在16帧上取平均，得到一个$7\times7\times512$的特征图；音频一侧，用在AudioSet上预训练的类VGG音频网络为每个1秒单元提取128维表示。两种特征都不再更新。两侧各接一个两层全连接网络，投影到共同空间后以欧氏距离$D=\lVert\bm r_{\mathrm v}-\bm r_{\mathrm a}\rVert_2$比较，其中$\bm r_{\mathrm v}$、$\bm r_{\mathrm a}$为两侧的投影结果。论文没有给出全连接层的宽度，也没有说明$7\times7\times512$的视觉特征在进入全连接层之前怎样展平或池化。

\begin{table}[htbp]
  \centering
  \small
  \begin{tabularx}{\linewidth}{@{}p{14mm}XXX@{}}
    \toprule
    分支 & 输入 & 特征提取 & 投影\\
    \midrule
    视觉 & 每秒16帧RGB & ImageNet预训练VGG-19的pool5，16帧平均为$7\times7\times512$ & 两层全连接，宽度论文未给出\\
    音频 & 每秒一段音频 & AudioSet预训练的类VGG网络，128维 & 两层全连接，宽度论文未给出\\
    距离 & 两侧投影 & — & 欧氏距离，对比损失间隔$m=2.0$\\
    \bottomrule
  \end{tabularx}
  \caption{AVE音画距离学习网络的结构配置，依据原论文。两种特征提取网络均为预训练后固定的模块。}
  \label{tab:vision-understanding-avdln-config}
\end{table}

若$y=1$表示同步配对，$y=0$表示错配，训练目标为
\begin{equation}
  L=yD^2+(1-y)[m-D]_+^2\text{，}\qquad m>0\text{。}
  \label{eq:vision-av-distance}
\end{equation}
同步对被拉近，错配对在距离小于$m$时被推开，超过$m$后不再受罚；原设置取$m=2.0$。标签来自训练配对关系，不是模型使用时预先知道目标时间。

设音频查询含$l$个时间单元，候选视频含$T$个时间单元。沿视频滑动同样长度的窗口，选择累计距离最小的起点
\begin{equation}
  \hat r=\argmin_{1\leq r\leq T-l+1}
    \sum_{j=1}^{l}D(\bm v_{r+j-1},\bm a_j)\text{。}
  \label{eq:vision-av-sliding}
\end{equation}
这里$\bm v_k$、$\bm a_j$表示已投影的片段特征。若单元时长为$\Delta$，且视频从0秒开始，第$r$个单元覆盖$[(r-1)\Delta,r\Delta)$，输出区间为$[(\hat r-1)\Delta,(\hat r+l-1)\Delta)$。

图\ref{fig:vision-av-temporal-match}用二维投影的教学数据走通三次窗口比较。取音频向量$(0,0)$、$(1,0)$，四段视频向量$(2,0)$、$(0,0.1)$、$(1,0.2)$、$(3,0)$。三个窗口的累计距离约为3.005、0.300和3.020，故选中第2、3单元；当$\Delta=1$秒时，区间为$[1,3)$秒。

\begin{figure}[htbp]
  \centering
  % 教学欧氏距离：声音(0,0),(1,0)，视频(2,0),(0,.1),(1,.2),(3,0)。
\begin{tikzpicture}[x=1cm,y=1cm,font=\figurefont\small]
  \node[anchor=west] at (0,3.5) {逐单元距离$D(\bm v_k,\bm a_j)$};
  \foreach \x/\k in {2.2/1,4.2/2,6.2/3,8.2/4}
    \node at (\x,2.7) {视频$\k$};
  \node[anchor=east] at (1.1,1.8) {声音1};
  \node[anchor=east] at (1.1,0.7) {声音2};
  \foreach \x/\y/\v in {
    2.2/1.8/2.000,4.2/1.8/0.100,6.2/1.8/1.020,8.2/1.8/3.000,
    2.2/0.7/1.000,4.2/0.7/1.005,6.2/0.7/0.200,8.2/0.7/2.000}{
    \draw[BookRule,line width=0.5pt,fill=BookPaper]
      (\x-0.9,\y-0.4) rectangle (\x+0.9,\y+0.4);
    \node at (\x,\y) {$\v$};
  }
  \draw[BookTeal,line width=1.2pt] (3.3,1.4) rectangle (5.1,2.2);
  \draw[BookTeal,line width=1.2pt] (5.3,0.3) rectangle (7.1,1.1);
  \node[anchor=west] at (0,-0.2) {窗口累计距离：};
  \node[align=center] at (1.7,-1.0) {视频1--2\\$3.005$};
  \node[align=center,text=BookTeal] at (4.8,-1.0) {视频2--3\\$\mathbf{0.300}$};
  \node[align=center] at (7.9,-1.0) {视频3--4\\$3.020$};
  \node[anchor=west] at (0,-2.0) {选中$[1,3)$秒；时间单元长1秒};
\end{tikzpicture}

  \caption{音频查询与视频窗口对应的教学计算。每列是一个视频时间单元；窗口内按相对次序逐项相加，粗框标出累计距离最小的窗口。向量和距离均为构造值，不是模型测量结果。}
  \label{fig:vision-av-temporal-match}
\end{figure}
\FloatBarrier

式\eqref{eq:vision-av-sliding}可以看成\ref{sec:vision-alignment-dtw}节动态时间规整的一个受限特例。动态时间规整允许一条序列的一个单元对应另一条序列的多个单元，用动态规划在所有单调路径中找累计代价最小的一条，适合两条序列速度不同的情形，例如同一句话说得快慢不同。滑动窗口只允许对角线路径：音频的第$j$个单元固定对应窗口内的第$j$个视频单元，于是只需比较$T-l+1$个起点。这个限制在这里是合理的，因为音频与视频来自同一时钟，事件的持续时间在两种模态中相同；若查询来自另一段录音，节奏与目标视频不同，就需要放开对应路径。原论文的评价只在事件短于10秒的测试视频上进行，并要求预测的片段与真实片段完全一致才算正确。

重复敲门可能在两处形成同样小的距离，最小值不再唯一；较粗的单元又限制边界精度。原工作另外研究逐秒预测“同时可见且可听的事件类别”，其监督是事件标签，输出是时间标签序列。该任务与上述查询对应的监督及输出均不同，不能把两套结果相互替代。

\subsection{界面控件的指令定位}
\label{sec:vision-understanding-context}

在界面上执行“保存当前文档”，需要把指令落到截图中的某个控件。\term{界面定位}（GUI grounding）的输入是一张截图和一条自然语言指令，输出是点击点或控件框。它与自然图像定位的难点不同：截图常有上千像素宽，目标控件可能只占极小面积；界面上文字密集，可能有两个“保存”按钮，一个属于当前对话框，另一个属于背景窗口。整屏缩小后，决定控件区别的小字可能消失。

两种代表做法分别从模型和观测两端解决这个问题。SeeClick通过大规模界面定位数据训练视觉语言模型，让它一次读完整张截图就直接生成坐标；它改进的是模型本身的定位能力。ScreenSeekeR不训练模型，而是把一次全屏读取改成多次局部搜索：先由规划模型推测目标可能所在的区域，再裁剪放大交给定位模型；它改进的是模型看到的观测。

\Needspace{60mm}
\subsubsection{SeeClick}
\label{sec:vision-understanding-seeclick}

此前的界面智能体多依赖网页HTML或系统无障碍树等结构化文本来确定控件，这些信息在许多应用中不可得，或者冗长难用；而通用视觉语言模型虽然能读截图，却很少学过把指令落到精确坐标。SeeClick的主要想法是只用截图，并把坐标写成文字，让语言模型像生成普通答案一样生成位置；所需的大量训练记录则从网页和移动界面自动收集\scite{vision-seeclick2024}。

模型以Qwen-VL为基础，结构沿用不变，各组件规模见表\ref{tab:vision-understanding-seeclick-config}\scite{vision-qwenvl2023}。截图缩放到$448\times448$，视觉编码器是以OpenCLIP的ViT-bigG初始化的ViT，按步长14切块，得到$32\times32=1024$个图像块特征；适配器是一层随机初始化的交叉注意力，用256个可学习查询读取这些特征，并在查询与键上加入二维绝对位置编码，把视觉信息压缩为256个单元；这些单元放在指令文字前面送入以Qwen-7B初始化的语言模型。输出是一段文字，例如\code{click (0.49, 0.40)}，两个数是保留两位小数、取值在$[0,1]$的比例，分别表示相对于截图宽、高的位置。这与Qwen-VL本身用$[0,1000)$整数表示框坐标的约定不同；SeeClick把数字当作普通文字处理，不增加专门的位置词元。由于位置以文字形式生成，模型不需要新增回归头，指令跟随和坐标预测共用同一个解码器。

\begin{table}[htbp]
  \centering
  \small
  \begin{tabularx}{\linewidth}{@{}p{16mm}XXp{12mm}@{}}
    \toprule
    组件 & 结构 & 输入与输出 & 参数量\\
    \midrule
    视觉编码器 & ViT，以OpenCLIP ViT-bigG初始化，切块步长14 & $448\times448$截图 \ensuremath{\to} 1024个图像块特征 & 1.9B\\
    适配器 & 单层交叉注意力，256个可学习查询，二维绝对位置编码 & 1024个特征 \ensuremath{\to} 256个视觉单元 & 0.08B\\
    语言模型 & 以Qwen-7B初始化的解码器 & 256个视觉单元与指令词元 \ensuremath{\to} 坐标文字 & 7.7B\\
    合计 & — & — & 9.6B\\
    \bottomrule
  \end{tabularx}
  \caption{SeeClick所用Qwen-VL的组件规模。参数量取自Qwen-VL论文，合计为原文给出的约数；视觉编码器的层数与宽度沿用ViT-bigG，SeeClick论文未另行给出。}
  \label{tab:vision-understanding-seeclick-config}
\end{table}

训练数据约100万条，混合三类来源：从约30万个网页自动提取的可见文字及其悬停提示与对应位置，约43万条；移动界面的控件描述、界面摘要等数据，约42万条；以及约14.5万条LLaVA一般视觉语言指令数据，用于保持原有能力。每条记录可以组织成“文字找位置”或“位置说文字”两种方向。监督是目标输出序列的交叉熵，模型逐词预测位置文本，不是用欧氏距离回归坐标。训练时用LoRA适配视觉编码器和语言模型，即在原权重旁加入低秩增量矩阵，只更新这些增量；约训练1万步，相当于一轮，全局批大小64，初始学习率$3\times10^{-5}$并按余弦衰减。论文同时建立了ScreenSpot评测集，含600余张截图、1200余条指令，覆盖iOS、Android、macOS、Windows和网页界面。

使用时，把截图和指令按训练格式输入，解析输出中的两个数，检查是否落在$[0,1]$之内，再乘以截图尺寸。若截图宽1600、高1000像素，输出$(0.49,0.40)$对应连续坐标$(784,400)$；实际点击还需与显示缩放和截图坐标一致。数字生成损失并不直接等于像素误差，“0.49”与“0.94”只差一个词元的顺序，像素距离却很远。定位评价可检查点是否落在目标控件框内；点落在正确按钮上，不能说明该按钮当前可用，更不能说明保存已经成功，多步操作还需在新截图中确认结果。

\subsubsection{ScreenSeekeR}
\label{sec:vision-understanding-screenseeker}

SeeClick一次读取整张截图，输入分辨率受模型限制。专业软件的高分辨率屏幕上，目标控件只占屏幕面积的极小比例，压缩到固定输入尺寸时可能只剩几个像素。ScreenSpot-Pro为此收集了专业应用的高分辨率截图作为评测基准：1581条指令，每条对应一张不同的截图，来自23个应用，分辨率多在1080p以上，目标平均只占截图面积的0.07\%；同一论文提出的ScreenSeekeR是一种搜索方法，二者名称不能互换\scite{vision-screenseeker2025}。ScreenSeekeR的主要想法是利用界面的层级布局知识缩小搜索区域：保存按钮通常在工具栏，工具栏通常在窗口上方，于是可以先找较大的区域，再逐级放大。

方法组合两个已有模型，不做训练。原实现以GPT-4o作为规划模型，以OS-Atlas-7B作为定位模型。规划模型读取截图和指令，用文字给出目标本身以及可能邻近的控件或区域；定位模型把这些文字描述落到框。候选区域的评分不直接采用规划模型的口头置信度，而是让定位框的中心为候选区域投票：将中心$(x,y)$转换为区域内部归一化位置$(x',y')$，在区域内时给予
\begin{equation}
  s=\exp\!\left(-\frac{(x'-0.5)^2+(y'-0.5)^2}{2\sigma^2}\right)
  \label{eq:vision-screenseeker-vote}
\end{equation}
的贡献，区域外贡献为0，原设置取$\sigma=0.3$。将所有定位框的贡献相加，再用非极大值抑制减少重复候选。靠近区域中心的多条线索可以提高该区域的优先级，但这种几何投票仍可能支持错误窗口。

搜索按评分依次进入裁剪区域，在其中继续提出位置或进行直接定位，直到规划模型确认找到目标、区域小于所定尺寸或到达最大深度。原设置在区域足够小时直接调用定位模型，这一尺寸超参数取1280像素；最小区域尺寸与最大搜索深度作为算法输入出现，论文未给出具体数值。这种逐级缩小的搜索，正是\ref{sec:vision-understanding-challenge-resolution}节所说多尺度观察的延续：经典方法在图像金字塔的每一层上滑窗，这里改由规划模型根据布局知识决定只放大哪几个区域。每次裁剪都应保存到父图的坐标变换，再组合回全屏。例如，父裁剪左上角为$(400,200)$，其内再裁出左上角$(100,50)$、宽高$200\times100$的区域；若最终归一化点为$(0.25,0.60)$，全屏点为$(400+100+0.25\times200,\ 200+50+0.60\times100)=(550,310)$。若中途调整过尺寸，还要先逆转该次缩放。

在ScreenSpot-Pro上，同一个OS-Atlas-7B定位模型单次读取全屏时的平均准确率为18.9\%，接入ScreenSeekeR的搜索后为48.1\%。两个数字使用同一定位模型和同一评测集，差别来自观测方式的改变，这与“小控件在整屏缩放中丢失”这一失败解释一致。

两种做法的比较见表\ref{tab:vision-understanding-gui-compare}。ScreenSeekeR保留了定位模型的单次读取方式，把预算从一次全屏读取改为多次局部搜索，换来了小控件的有效分辨率，代价是多次调用大模型的时间与费用，以及对规划模型布局知识的依赖。裁剪虽然放大了目标，也可能删去“当前窗口”等消歧线索，因此递归状态仍需保留整屏上下文与候选所属位置。两者都只对当前截图产生定位，历史动作和界面状态变化应作为下一次观测重新进入系统。

\begin{table}[htbp]
  \centering
  \small
  \begin{tabularx}{\linewidth}{@{}p{16mm}XX@{}}
    \toprule
    比较项 & SeeClick & ScreenSeekeR\\
    \midrule
    改进对象 & 模型的定位能力 & 模型看到的观测\\
    训练 & 界面定位数据上LoRA适配 & 不训练，组合已有模型\\
    每条指令的调用 & 一次生成 & 多次规划、定位与裁剪\\
    高分辨率处理 & 整屏缩放到固定输入 & 逐级裁剪放大\\
    主要失败 & 小控件在缩放中丢失 & 规划误导、裁剪丢失上下文\\
    \bottomrule
  \end{tabularx}
  \caption{两种界面定位做法的对照。两者输出都是截图坐标，评价口径都是点是否落在目标控件内。}
  \label{tab:vision-understanding-gui-compare}
\end{table}

\section{文字读取与结构化解析}
\label{sec:vision-understanding-reading}

定位找到了票据区域，但“看见数字”和“知道数字属于哪个字段”是两回事。表格中，即使所有字符都读对，合并表头归属错误仍会改变含义；图表中，柱子的高度只有对照刻度和图例才成为数值。本节的输入是含文字或符号的图像，输出是字符串以及它们之间的关系：阅读顺序、表格的行列、字段的键值、图表的数据记录、图解的连接。

本节按被恢复的对象展开，对象的结构由简到繁。字符串的读取是其余各项的基础，要回答“这里写了什么”。页面结构的解析在字符串之上恢复区域类型和阅读关系。表格、表单、图表和图解各有专门的结构：表格靠行列与跨格，表单靠键与值，图表靠几何编码与刻度，图解靠节点与有向连线。后四种对象的恢复大多复用前两项的模型，再加上对象特有的结构规则，因此不再按模型细分。

\subsection{字符串的读取}
\label{sec:vision-understanding-ocr}

\term{光学字符识别}（optical character recognition, OCR）把文字图像转换为机器可处理的字符串。输入是一幅场景图或页面图，输出是若干文字实例，每个实例带一个框或多边形以及一个字符串。完整流程通常分为两步：文字检测找出哪些区域包含文字，文字识别把每个区域读成字符串。位置正确而内容读错、内容正确但属于另一行，是两种不同错误，分别出现在这两步。

经典的文档OCR用\ref{chap:vision-image-features}章的工具完成这两步。检测阶段先二值化：文档中墨迹与纸张的灰度分布通常呈两个峰，Otsu方法（\ref{sec:vision-features-otsu}节）选择使两类类间方差最大的阈值，光照不均时改用局部自适应阈值；随后做连通分量标记（\ref{sec:vision-formation-connectivity}节），每个分量大致是一个字符或字符的一部分，再按分量的位置和大小把它们聚成行，例如沿水平方向统计墨迹像素数，用投影的谷值切分行与字。识别阶段把每个字符图像缩放到固定大小，与字库模板做归一化互相关（\ref{sec:vision-features-template-matching}节），或提取HOG等梯度特征（\ref{sec:vision-features-hog}节）后交给分类器。这条流水线在白纸黑字、版式规整的扫描件上有效，因为它依赖两个条件：前景与背景可以用一个阈值分开，字符彼此分离而且排成直线。场景文字两个条件都不满足：招牌背景复杂，阈值分不开文字与纹理；字符可能粘连、弯曲或带透视变形，连通分量与直线投影都会失效。下面的两种学习模型保留了这条流水线“先找部件、再组合、再识别”的结构，但把每个部件换成了从数据中学到的判断。

两步各有难点。检测要处理弯曲、倾斜、长短悬殊的文字行，直接回归一个矩形框可能包含大量背景，也可能把一行切断。识别要处理模糊、手写和罕见字符串，还要决定语言上下文能在多大程度上帮助判断字形。CRAFT回答检测的难点：先找单个字符，再判断相邻字符是否相连，由下而上拼出任意形状的文字。TrOCR回答识别的难点：用预训练的图像Transformer读取文字行，再用预训练的文字Transformer逐词生成字符串。两者处在流程的不同环节，常被串联使用。

\subsubsection{CRAFT}
\label{sec:vision-understanding-craft}

弯曲招牌中的字符未必落在同一条直线上，以整词为单位回归框的检测器难以给出贴合的边界；长文字行的感受野也可能不够覆盖整行。CRAFT的主要想法是把检测拆成两个局部判断：哪里是字符中心，哪里是两个相邻字符之间的连接。局部判断只需看字符尺度的邻域，组合起来却能描出任意长度、任意弯曲的文字行\scite{vision-craft2019}。

模型是一个全卷积网络。输入为任意尺寸的$H\times W\times3$图像，编码器采用带批归一化的VGG-16，解码器逐级上采样并通过跳跃连接融合编码器各层的特征，结构类似U-Net；卷积网络的基本运算可参考\BookChapterRef{chap:convolutional-networks}{卷五的卷积神经网络章}。输出是两张稠密图：\emph{字符区域图}$S_{\mathrm r}$给出每个位置作为字符中心的程度，\emph{关联图}$S_{\mathrm a}$给出每个位置作为相邻字符连接处的程度。

论文只给出了骨干和跳跃连接的示意，表\ref{tab:vision-understanding-craft-layers}的逐层配置依据作者发布的官方实现。编码器取VGG-16-BN的前若干卷积组，在四个分辨率上取出跳跃特征$\bm s_1,\ldots,\bm s_4$；末端用一个步长为1的最大池化和一个扩张率为6的$3\times3$卷积代替原VGG的全连接层fc6，再接$1\times1$卷积代替fc7，在不降低分辨率的情况下扩大感受野，使长文字行的上下文也能被看到。解码器有四个上采样块，每块的计算顺序是：把上一级输出用双线性插值放大到下一个跳跃特征的尺寸，与该跳跃特征沿通道拼接，经$1\times1$卷积、批归一化、ReLU压缩通道，再经$3\times3$卷积、批归一化、ReLU输出。第一块的输入不需要放大，直接拼接fc7的输出与$\bm s_4$。最后的分类头由三层$3\times3$卷积和两层$1\times1$卷积组成，输出两个通道，分别作为$S_{\mathrm r}$与$S_{\mathrm a}$。由于最浅的跳跃特征位于$1/2$分辨率，输出图的边长是输入的一半。

\begin{table*}[tbp]
  \centering
  \small
  \begin{tabularx}{\linewidth}{@{}p{24mm}Xp{24mm}p{14mm}p{10mm}p{22mm}@{}}
    \toprule
    阶段或层组 & 层类型 & 卷积核与步长 & 输出通道 & 重复 & 输出尺寸\\
    \midrule
    卷积组1 & 卷积、批归一化、ReLU & $3\times3$，1 & 64 & 2 & $H\times W$\\
    卷积组2 & 最大池化；卷积、批归一化、ReLU & $2\times2$，2；$3\times3$，1 & 128 & 1；2 & $H/2$，取出$\bm s_1$\\
    卷积组3 & 最大池化；卷积、批归一化、ReLU & $2\times2$，2；$3\times3$，1 & 256 & 1；2 & $H/4$，取出$\bm s_2$\\
    卷积组4 & 卷积；最大池化；卷积、批归一化、ReLU & $3\times3$，1；$2\times2$，2；$3\times3$，1 & 256；512 & 1；1；2 & $H/8$，取出$\bm s_3$\\
    卷积组5 & 卷积；最大池化；卷积、批归一化、ReLU & $3\times3$，1；$2\times2$，2；$3\times3$，1 & 512；512 & 1；1；2 & $H/16$，取出$\bm s_4$\\
    代替fc6、fc7 & 最大池化；扩张卷积；卷积 & $3\times3$，1；$3\times3$扩张6，1；$1\times1$，1 & 1024 & 各1 & $H/16$\\
    上采样块1 & 拼接fc7输出与$\bm s_4$；卷积；卷积 & $1\times1$；$3\times3$ & 512；256 & 1 & $H/16$\\
    上采样块2 & 放大并拼接$\bm s_3$；卷积；卷积 & $1\times1$；$3\times3$ & 256；128 & 1 & $H/8$\\
    上采样块3 & 放大并拼接$\bm s_2$；卷积；卷积 & $1\times1$；$3\times3$ & 128；64 & 1 & $H/4$\\
    上采样块4 & 放大并拼接$\bm s_1$；卷积；卷积 & $1\times1$；$3\times3$ & 64；32 & 1 & $H/2$\\
    分类头 & 卷积与ReLU；卷积与ReLU；卷积 & $3\times3$；$1\times1$；$1\times1$ & 32、32、16；16；2 & 3；1；1 & $H/2\times W/2\times2$\\
    \bottomrule
  \end{tabularx}
  \caption{CRAFT的逐层结构，依据官方实现。输出尺寸只写高度，宽度按相同比例变化；例如输入$768\times768$时，输出为$384\times384\times2$。官方实现在原VGG第2组最后一个卷积之后、第3至5组倒数第二个卷积之后取出跳跃特征（取在批归一化之后、ReLU之前）；组内剩余的卷积在下一段开头执行，表中按执行顺序写在下一行，原VGG第5组的最后一个卷积未被使用。}
  \label{tab:vision-understanding-craft-layers}
\end{table*}

训练目标需要两张图的标注，而大多数真实数据只有词级框。合成数据集SynthText提供字符框，可把二维高斯模板通过透视变换映射到每个字符框内，形成区域目标$S_{\mathrm r}^*$；相邻两个字符框各取上下三角形中心，连成的四边形内再放高斯模板，形成关联目标$S_{\mathrm a}^*$。对只有词级框和字符串的真实数据，原方法用当前模型自己生成字符级伪标注：裁出词区域，预测区域图，用分水岭算法（\ref{sec:vision-features-watershed}节）切分出若干字符框，再把坐标还原到原图。分水岭把区域图看成地形，每个字符中心是一个盆地，从盆地向外“注水”直到相邻盆地相遇，相遇处就是字符之间的分界。

伪标注可能切错，因此要衡量可靠性。设一个词的真实长度为$l(w)$，切分得到的字符数为$l^{\mathrm c}(w)$，置信度定义为
\begin{equation}
  s_{\mathrm{conf}}(w)=\frac{l(w)-\min\bigl(l(w),\,|l(w)-l^{\mathrm c}(w)|\bigr)}{l(w)}\text{。}
  \label{eq:vision-craft-conf}
\end{equation}
切出的字符数与真实长度一致时置信度为1，相差越多越低。若$s_{\mathrm{conf}}<0.5$，就放弃切分结果，改为按字符数均分词区域，并把置信度置为0.5。词内各像素取所属词的置信度，组成置信图$S_{\mathrm c}$，合成数据上取1。训练损失为逐像素加权的平方误差
\begin{equation}
  L=\sum_{p}S_{\mathrm c}(p)\Bigl(\bigl\lVert S_{\mathrm r}(p)-S_{\mathrm r}^*(p)\bigr\rVert_2^2
   +\bigl\lVert S_{\mathrm a}(p)-S_{\mathrm a}^*(p)\bigr\rVert_2^2\Bigr)\text{。}
  \label{eq:vision-craft-loss}
\end{equation}
置信度低的伪标注贡献较小，模型不会被明显切错的字符框带偏。原论文先在SynthText上训练5万次迭代，再在各真实数据集上微调2.5万次迭代，微调时按1比5的比例混入合成数据，并以1比3的正负像素比例做困难样本挖掘；各阶段都用Adam优化。论文没有给出学习率、批大小和训练图像尺寸。

推理时，将区域图超过阈值$\tau_{\mathrm r}$或关联图超过阈值$\tau_{\mathrm a}$的位置标为1，得到二值图；做连通分量标记（\ref{sec:vision-formation-connectivity}节），每个分量就是一个文字实例，再为其拟合最小面积的旋转矩形，或沿字符中心构造适合弯曲文字的多边形。论文没有给出两个阈值的取值，官方实现默认以0.4作为区域图和关联图的二值化阈值，并只保留区域图最大值超过0.7的分量。这一步与经典流水线中对二值化墨迹做连通分量分组的做法相同，区别在于被分组的是学到的字符响应和关联响应：字符之间即使在像素上粘连或断开，只要关联图的响应正确，分组就仍然正确。教学招牌“青山路”若三个字符响应之间有两条正确关联，会形成一个文字实例；若“路”与旁边门牌之间也出现强关联，分组就可能错误合并。此时问题发生在文字检测，后续识别器再准确也无法恢复正确的实例边界。

\subsubsection{TrOCR}
\label{sec:vision-understanding-trocr}

文字区域确定后，还要把裁出的文字图像读成字符串。此前的识别器通常以卷积网络提取特征，接循环网络逐列建模，再用连接时序分类等方法输出字符，需要时另接外部语言模型。这种结构的视觉与语言部分难以直接利用大规模预训练的Transformer。TrOCR的主要想法是把识别写成图像到文字的序列转换：编码器和解码器分别用已有的图像Transformer和文字Transformer初始化，二者之间只新增交叉注意力，由识别数据学会连接\scite{vision-trocr2022}。

输入文字图像$\bm X_{\mathrm{text}}$被缩放到$384\times384$，切成$16\times16$的图像块，共$24\times24=576$块；每块展平后线性投影为$d$维，加上位置表示，送入从BEiT或DeiT初始化的ViT编码器，输出$576\times d$的视觉表示。解码器从RoBERTa或MiniLM初始化，按标准编码器—解码器结构读取已生成的词元，并以交叉注意力读取视觉表示，预测下一个子词词元；交叉注意力的计算可参考\BookChapterRef{chap:transformer}{卷五的Transformer章}。这一固定尺寸设置把长文字行压成正方形，与后文Pix2Struct保持纵横比的设计不同。

论文给出了三种规模，见表\ref{tab:vision-understanding-trocr-config}。三种规模的输入都是$384\times384$、图像块$16\times16$，编码器序列为576个图像块，另加一个[CLS]词元，用DeiT初始化时再加一个蒸馏词元。解码器层数少于用来初始化的文字模型时，取该模型的后一半层，例如从24层的RoBERTa-Large取后12层；文字模型原本只有编码器，没有交叉注意力，这部分参数随机初始化。

\begin{table*}[tbp]
  \centering
  \small
  \begin{tabularx}{\linewidth}{@{}p{16mm}p{14mm}p{12mm}Xp{18mm}Xp{14mm}@{}}
    \toprule
    规模 & 图像块 & 词元数 & 编码器（初始化；层数、宽度、头数） & 解码器初始化 & 解码器（层数、宽度、头数） & 参数量\\
    \midrule
    TrOCR-Small & $16\times16$ & 576 & DeiT-Small；12层，384，6头 & MiniLM & 6层，256，8头 & 62M\\
    TrOCR-Base & $16\times16$ & 576 & BEiT-Base；12层，768，12头 & RoBERTa-Large & 12层，1024，16头 & 334M\\
    TrOCR-Large & $16\times16$ & 576 & BEiT-Large；24层，1024，16头 & RoBERTa-Large & 12层，1024，16头 & 558M\\
    \bottomrule
  \end{tabularx}
  \caption{TrOCR的三种规模，依据原论文的实验设置。词元数为编码器的图像块数，不含特殊词元；MLP倍率论文未单独给出，沿用各初始化模型。}
  \label{tab:vision-understanding-trocr-config}
\end{table*}

编码器的单个块采用ViT的写法（\ref{sec:vision-ivm-vit}节），层归一化放在子层之前：
\[
  \bm z'=\bm z+\operatorname{MSA}(\operatorname{LN}(\bm z)),\qquad
  \bm z''=\bm z'+\operatorname{MLP}(\operatorname{LN}(\bm z'))\text{。}
\]
解码器的单个块依次做三步：带因果掩码的自注意力，使第$j$个位置只能看到已生成的词元；以编码器输出为键和值的交叉注意力；前馈网络。每步都带残差连接和层归一化。文字图像的信息只通过交叉注意力进入解码器，而解码器的自注意力承载语言上下文，这两个通道的相对强弱决定了后文所说语言先验的影响程度。

以文字图像$\bm X_{\mathrm{text}}$和目标词元$y_1,\ldots,y_m$为一条记录，训练最小化
\begin{equation}
  L_{\mathrm{OCR}}
  =-\sum_{j=1}^{m}\log p_{\bm\theta}
    (y_j\mid y_{<j},\bm X_{\mathrm{text}})\text{。}
  \label{eq:vision-trocr-loss}
\end{equation}
训练分阶段进行：第一阶段在合成的印刷文字行上预训练，共6.84亿行；第二阶段用规模较小、分别面向任务构建的数据继续预训练，印刷体约330万行，手写体约1790万行，场景文字则用MJSynth与SynthText约1600万幅；最后在票据（SROIE）、手写（IAM）和场景文字基准上微调。全部参数参与更新，批大小2048，学习率$5\times10^{-5}$。输出用BPE与SentencePiece分词，得到子词词元。

使用时从起始符逐词解码，直到结束符或长度上限；用束搜索保留若干候选，原设置的束宽为10，最终去掉特殊词元并恢复字符串。这里的预测单元是子词，一步未必恰好对应一个字符。模糊字形“1”和“I”可能在视觉上相近，语言上下文会影响选择：对常见地名，这有时有助于消歧；对票据编号或未见专名，强语言先验也可能把正确的罕见字符串改成常见词。因此应保留原裁剪和识别结果用于核对。

CRAFT与TrOCR分工明确，比较时要放在各自的环节中看。CRAFT以字符和连接为单位，换来了对弯曲和长文字行的适应，代价是依赖伪标注质量，且字符间距很大或很小时容易断开或误连；它不输出字符串。TrOCR把识别交给预训练的编码器和解码器，换来了语言上下文和手写体的处理能力，代价是固定输入尺寸和语言先验可能改写罕见字符串；它不给出全页文字位置、双栏阅读顺序或表格归属。二者串联时，检测阶段的合并或截断错误会原样传给识别阶段。

\subsection{页面结构的解析}
\label{sec:vision-understanding-page}

一页文档上的文字可能属于标题、正文、表格、脚注或公式。把左栏第一行和右栏第一行交错拼接，即使每个字都正确，仍会得到错误句子。这一小节的输入是页面图像，输出是带类型的页面元素、文字所属的字段或区域，以及元素之间的阅读关系。难点在于：语义既依赖文字内容，也依赖二维位置和视觉样式，三者需要在同一个表示中联系起来。

两条路线对“文字从哪里来”给出不同回答。LayoutLMv3依赖外部OCR提供词及其框，把文字、二维布局和图像块共同编码，再接分类或抽取任务头输出结构；它擅长在已读出的文字上做判断。Pix2Struct不使用外部OCR，直接从像素生成结构化文字序列；它省去OCR环节，但要求视觉编码器自己保留小字。

\subsubsection{LayoutLMv3}
\label{sec:vision-understanding-layoutlm}

此前的文档预训练模型通常用卷积网络或检测器提取页面图像特征，并为文字和图像设计不同的预训练目标，例如文字一侧做掩码语言建模，图像一侧做像素重建。两种目标的粒度不一致，难以让模型学到“这个词对应页面上哪块图像”。LayoutLMv3的主要想法是统一掩码：文字和图像块都用离散单元表示、都做掩码预测，再加一个专门的词块对齐目标，把文字与它所在的图像块联系起来；图像一侧也不再使用卷积骨干，只做线性投影\scite{vision-layoutlmv32022}。

文字一侧的输入来自外部OCR：每个词经分词得到词元，词元嵌入由RoBERTa初始化，再加上一维序列位置和二维布局位置。布局位置把所属文字段的框坐标按页面尺寸归一化，分别嵌入横纵坐标与宽高；同一文字段内的词共享一个框。图像一侧把页面缩放到$224\times224$，切成$16\times16$的块，共196块，各块线性投影后加上一维位置嵌入。两组单元拼接后送入Transformer，输出每个词元和每个图像块的上下文表示。这样，“128.50”的表示可以同时参考附近的“总额”及其页面位置。表\ref{tab:vision-understanding-layoutlmv3-config}给出两种规模的配置。文字序列至多512个词元，用BPE分词；嵌入层和Transformer以RoBERTa的权重初始化。图像块只加一维可学习位置嵌入，论文报告加二维位置嵌入没有带来改善。

\begin{table*}[tbp]
  \centering
  \small
  \begin{tabularx}{\linewidth}{@{}p{22mm}p{12mm}Xp{8mm}p{10mm}p{8mm}p{14mm}p{12mm}@{}}
    \toprule
    规模 & 图像块 & 词元数 & 层数 & 宽度 & 头数 & MLP倍率 & 参数量\\
    \midrule
    LayoutLMv3-Base & $16\times16$ & 至多512个文字词元，196个图像块 & 12 & 768 & 12 & 4（3072） & 133M\\
    LayoutLMv3-Large & $16\times16$ & 至多512个文字词元，196个图像块 & 24 & 1024 & 16 & 4（4096） & 368M\\
    \bottomrule
  \end{tabularx}
  \caption{LayoutLMv3的两种规模，依据原论文。图像输入为$224\times224$；MLP倍率一栏括号内为前馈层中间维数。}
  \label{tab:vision-understanding-layoutlmv3-config}
\end{table*}

单个Transformer块沿用RoBERTa的结构，在残差相加之后做层归一化，但自注意力的打分加入了位置偏置。设第$i$、$j$个单元的查询与键为$\bm q_i$、$\bm k_j$，每个头的维数为$d_h$，打分为
\[
  e_{i,j}=\frac{\bm q_i^\top\bm k_j}{\sqrt{d_h}}+b^{\mathrm{1D}}_{i,j}+b^{\mathrm{2D}}_{i,j}\text{，}
\]
其中$b^{\mathrm{1D}}_{i,j}$由两个单元在序列中的相对距离决定，$b^{\mathrm{2D}}_{i,j}$由两个文字段在页面上的横、纵相对距离决定，二者都是可学习的偏置，沿用LayoutLMv2的做法。随后$\bm Z'=\operatorname{LN}(\bm Z+\operatorname{MSA}(\bm Z))$，$\bm Z''=\operatorname{LN}(\bm Z'+\operatorname{FFN}(\bm Z'))$。二维相对偏置让“同一行右侧相邻”这类版面关系直接影响注意力，而不必完全依靠绝对坐标嵌入去推断。

预训练使用约1100万页IIT-CDIP文档图像，目标由三项相加：
\begin{equation}
  L=L_{\mathrm{MLM}}+L_{\mathrm{MIM}}+L_{\mathrm{WPA}}\text{。}
  \label{eq:vision-layoutlmv3-loss}
\end{equation}
文字掩码项$L_{\mathrm{MLM}}$以片段方式遮去30\%的文字词元，片段长度服从参数为3的泊松分布，用交叉熵预测被遮词元。图像掩码项$L_{\mathrm{MIM}}$按块状区域遮去约40\%的图像块，预测这些块的离散视觉编码；编码由一个离散变分自编码器给出，它以文档图像模型DiT的图像分词器初始化，词表含8192个视觉词元，而不是直接回归原始像素。\term{词块对齐}（word-patch alignment, WPA）项只作用于未被遮的文字词元：若该词对应的图像块也未被遮，标为“对齐”，否则标为“未对齐”，由一个两层感知机以二元交叉熵预测。预训练共50万步，批大小2048。若一行金额文字位于图像块7和8，遮去这两块而保留词元时，对齐标签就变为“未对齐”；模型要答对，必须知道这行文字在页面上的位置对应哪些图像块。这个目标利用OCR位置建立的联系，不是判断任意一段文字与整幅图是否语义一致，也不能保证OCR本身读对。

下游使用时在预训练表示上加任务头并微调全部参数。表单理解按词元或文字段预测“问题”“答案”“标题”“其他”等类别；票据信息抽取预测字段标签；文档问答在词元上预测答案的起止位置；文档分类读取整体表示。版面分析是另一种配置：不输入文字，只把LayoutLMv3的图像部分作为检测骨干，从基础版的第4、6、8、12层取出特征，经改变分辨率的模块变成多尺度特征，再接带特征金字塔的Cascade R-CNN检测器，输出标题、表格、图片等区域框。两种配置的输入不同，LayoutLMv3因此不能被理解为一个无条件输出整页HTML及阅读顺序的解码器；能恢复什么，由输入、任务头和监督决定。

\subsubsection{Pix2Struct}
\label{sec:vision-understanding-pix2struct}

LayoutLMv3式方法在使用时必须先运行OCR，OCR的错误和耗时都会传入下游；同时，把页面缩放到固定正方形会让长条网页或宽表格的小字失真。Pix2Struct的主要想法有两点：一是不用OCR，直接从像素生成结构化文字；二是按输入预算可变地缩放图像，尽量保留原纵横比。它的预训练任务是把网页截图解析成简化的HTML，网页自带的结构为“从像素到结构”提供了大量免费监督\scite{vision-pix2struct2023}。

编码器是ViT。给定最多$n_{\max}$个$16\times16$图像块的预算，图像按原纵横比放大或缩小，使切出的块数不超过预算且尽可能多；块的位置用行号和列号两套二维绝对位置嵌入表示，因此不同形状的页面都能被正确编码。原论文预训练时预算为2048块。解码器是Transformer文字解码器，读取编码器输出并逐词生成目标序列。论文把这种结构称为基于ViT的图像编码器—文字解码器，并把整体类比为视觉上的T5：各种任务都化为“输入图像、输出文字”。两种规模的配置见表\ref{tab:vision-understanding-pix2struct-config}。

\begin{table*}[tbp]
  \centering
  \small
  \begin{tabularx}{\linewidth}{@{}p{22mm}p{12mm}Xp{16mm}p{10mm}p{10mm}p{12mm}p{12mm}@{}}
    \toprule
    规模 & 图像块 & 编码器块数上限 & 层数 & 宽度 & 头数 & MLP倍率 & 参数量\\
    \midrule
    Pix2Struct-Base & $16\times16$ & 预训练2048；微调4096（信息图问答6144） & 编码器12，解码器12 & 768 & 论文未给出 & 论文未给出 & 282M\\
    Pix2Struct-Large & $16\times16$ & 预训练2048；微调3072 & 18（编码器与解码器的划分论文未给出） & 1536 & 论文未给出 & 论文未给出 & 1.3B\\
    \bottomrule
  \end{tabularx}
  \caption{Pix2Struct的两种规模，依据原论文。块数上限是可变分辨率输入的预算，实际块数随页面纵横比变化；预训练时解码器序列长度为128个词元。}
  \label{tab:vision-understanding-pix2struct-config}
\end{table*}

可变分辨率的计算可以用一个教学例说明。设页面为$1600\times900$像素，预算为2048块。保持纵横比缩放$\alpha$倍后，块网格为$\lfloor1600\alpha/16\rfloor\times\lfloor900\alpha/16\rfloor$；要求乘积不超过2048，约有$100\alpha\times56.25\alpha\leq2048$，即$\alpha\leq0.603$。取$\alpha=0.6$得到$60\times33=1980$块，页面被缩到约$960\times540$像素。若改用$224\times224$的固定正方形，横向缩小约7倍、纵向缩小约4倍，字形还被拉扁。二维位置嵌入用行号和列号分别查表，所以$60\times33$与$33\times60$两种网格都能被正确编码。

预训练数据来自C4语料中网址对应的约8000万张网页截图及其HTML。原方法保留可见内容相关的节点，压缩冗长的嵌套，只保留节点的文字与图像替代信息，再选取能放入目标长度的子树；截图上标出该子树的区域，并以随机片段遮去其中一半文字，训练解码器生成完整的简化结构。训练目标是目标序列的逐词交叉熵，形式与式\eqref{eq:vision-trocr-loss}相同。直接从解析任务开始训练不稳定，因此先做阅读热身：把BooksCorpus的文本以随机字体和颜色渲染成图，输入上限128块，训练3万步读出文字，再进入截图解析。截图解析阶段，基础版再训练27万步、批大小2048，大型版训练17万步、批大小1024。

下游任务以同一接口微调。问题等文字条件不作为单独的文字输入，而是渲染在图像顶部作为标题区，与页面一同被编码；指代任务则把框直接画在截图上。论文在四个领域的九个任务上微调：插图类的ChartQA、AI2D、OCR-VQA，界面类的RefExp、Widget Captioning、Screen2Words，自然图像类的TextCaps，以及文档类的DocVQA和InfographicVQA。使用时必须遵循相应检查点的预处理和输出约定。若教学页面具有“发票—编号A017—总额128.50”的层次，可以用嵌套节点表示从属关系；解析器应检查节点是否闭合、文字是否遗漏及金额属于哪个字段。字符串符合语法，只说明可以读入结构，不说明内容忠实于页面。改变块预算会改变可保留的视觉细节，解码更长的字符串不能恢复已经因缩小而不可辨认的小数点。

页面元素恢复后，还要按既定关系组织阅读顺序。双栏正文应在本栏顺读，再转入下一栏；跨栏标题和所属公式要与相应内容连接。OmniDocBench同时提供完整文档解析和文字、版面、表格、公式等模块的评价\scite{vision-omnidocbench2025}。这些分项使整页误差能够定位到具体结构，而非仅用一串文字相似度概括。

两种路线的对照见表\ref{tab:vision-understanding-page-compare}。Pix2Struct把LayoutLMv3依赖的外部OCR换成了模型自己的读取，把固定正方形输入换成了可变分辨率，换来了无需OCR的单一接口和对网页结构的利用；代价是所有文字都要从像素中读出，小字能否保留完全取决于块预算，且生成的结构串需要额外的语法与内容检查。LayoutLMv3在OCR可靠时能直接利用准确的文字和坐标，输出也天然对应到具体词元，便于追溯来源。

\begin{table}[htbp]
  \centering
  \small
  \begin{tabularx}{\linewidth}{@{}p{14mm}XX@{}}
    \toprule
    比较项 & LayoutLMv3 & Pix2Struct\\
    \midrule
    文字来源 & 外部OCR的词与框 & 模型从像素读取\\
    图像输入 & 固定$224\times224$ & 按块预算可变缩放，保持纵横比\\
    预训练 & 文字与图像块掩码预测、词块对齐 & 截图到简化HTML的生成\\
    输出方式 & 任务头：分类、跨度、检测 & 自回归生成结构化文字\\
    来源追溯 & 输出对应到OCR词元及其框 & 需另行对齐生成内容与页面\\
    主要失败 & OCR错误直接传入 & 缩放丢失小字、结构串不忠实\\
    \bottomrule
  \end{tabularx}
  \caption{两种页面解析路线的对照。比较以同一页面、同一结构化任务为前提。}
  \label{tab:vision-understanding-page-compare}
\end{table}

\subsection{表格结构的解析}
\label{sec:vision-understanding-tables}
\label{sec:vision-understanding-pubtables}

表格中，一个数的含义来自它所在的行、列和表头。输入是页面或裁剪的表格图像，输出是逻辑表：每个单元的行列索引、跨度、内容和来源框。难点在于行列边界往往没有画线，表头可能跨多行多列，仅按OCR字符串的出现顺序排列会丢失这些关系。

有框线的表格可以用经典方法恢复。二值化后，用一个很长的水平线段作结构元素做形态学开运算（\ref{sec:vision-transforms-binary-morphology}节），只有长度超过结构元素的水平笔画能够保留，文字笔画被去掉，剩下的就是横线；竖线同理。也可以在边缘图上用Hough变换（\ref{sec:vision-features-hough}节）检测直线：每个边缘点在参数空间中为所有经过它的直线投票，票数集中的参数就是图中的长直线。横线与竖线的交点给出单元格的角点。这类方法在框线完整时可靠，但许多表格只有表头下的一条横线，甚至完全没有框线，行列边界只能从文字的对齐和间距推断；跨格表头也不会在框线上留下额外线索。学习方法改为直接预测行、列和跨格的区域，不再依赖框线是否存在。

PubTables-1M的做法是把表格结构当作一组可检测的对象\scite{vision-pubtables2022}。此前的表格数据集标注常不一致，例如同一个跨两列的表头，有的标为一格，有的标为一格文字加一格空白；模型因而学到不一致的结构。该工作从PubMed Central开放获取论文中收集近百万张表格，并对标注做\emph{规范化}：按统一规则修正过度切分，使同一种表头含义对应稳定的结构。主要想法是：行、列、表头和跨格单元都是有明确边界的矩形，用目标检测就能预测它们，再由几何关系组装成逻辑表。

模型分为两个沿用DETR的检测器（\ref{sec:vision-recognition-detr}节）：一个在页面上定位整张表格，另一个在裁剪的表格图像中预测六类对象，即整表、行、列、列表头、投影行表头和跨格单元。\term{投影行表头}（projected row header）是横跨所有列、用于划分后续行组的标题行。两个检测器都以框和类别为监督，训练沿用DETR的二分匹配、类别交叉熵、$L_1$和广义IoU损失。数据集用于结构识别的部分含947642张表格，其中52.7\%至少含一个跨格单元。

两个检测器的结构相同，只有对象查询数不同，见表\ref{tab:vision-understanding-pubtables-config}。查询数取得略多于训练样本中对象数的最大值：表格检测取15个，结构识别取125个，因为一张表的行、列、表头和跨格对象远多于一页上的表格数。图像经ResNet-18主干得到步长32的特征图，用$1\times1$卷积降到256维并展平，加上位置编码后送入6层编码器；125个对象查询经6层解码器读取编码器输出，每个查询再由线性分类头输出6类对象或“无对象”，由三层感知机输出归一化的框。

\begin{table}[htbp]
  \centering
  \small
  \begin{tabularx}{\linewidth}{@{}p{14mm}XX@{}}
    \toprule
    部分 & 配置 & 输入与输出\\
    \midrule
    主干 & ResNet-18，ImageNet预训练，前几层冻结 & $H\times W\times3$ \ensuremath{\to} $H/32\times W/32\times512$\\
    投影 & $1\times1$卷积到256维，加位置编码 & 展平为$(HW/1024)\times256$\\
    编码器 & 6层，宽度256，8头，前馈2048，层归一化前置 & 序列长度不变\\
    解码器 & 6层；表格检测15个查询，结构识别125个查询 & $15\times256$或$125\times256$\\
    预测头 & 类别（结构识别为6类加“无对象”）与框 & 每个查询一组类别概率和框\\
    训练 & 20轮；结构识别初始学习率$5\times10^{-5}$，“无对象”类权重0.4，不用辅助解码损失 & —\\
    \bottomrule
  \end{tabularx}
  \caption{PubTables-1M所用DETR的配置。主干、层数、查询数、轮数、学习率与“无对象”权重见论文附录；宽度、头数、前馈维数、层归一化位置与辅助损失设置取自官方代码的结构识别配置文件。}
  \label{tab:vision-understanding-pubtables-config}
\end{table}

图\ref{fig:vision-table-structure}用四行三列的教学表说明组装过程。前两行为表头：“项目”跨两行，“金额（元）”跨第2、3列，其下分别为1月、2月。后两行是甲的$(12,18)$与乙的$(8,10)$。行列对象相交先产生12个基本格；两个跨格对象指出哪些基本格应合并，合并后形成10个逻辑单元；列表头对象标出前两行的功能。跨格表头使每个金额都能回溯到月份和单位。

\begin{figure*}[tbp]
  \centering
  % 构造表：4行3列；两个跨格对象各合并2格，得到10个逻辑单元。
\begin{tikzpicture}[x=1cm,y=1cm,font=\figurefont\small,
  arr/.style={-{Stealth[length=2mm]},BookInk,line width=0.9pt}]
  \node[anchor=west,font=\figurefont\bfseries\small] at (0,4.0)
    {行列交叉与跨格对象};
  \node[anchor=west,font=\figurefont\bfseries\small] at (9,4.0)
    {恢复后的逻辑表};
  \foreach \dx in {0,9}{
    \begin{scope}[xshift=\dx cm]
      \fill[BookBlue!12] (0,1.6) rectangle (2,3.2);
      \fill[BookTeal!12] (2,2.4) rectangle (6,3.2);
      \draw[BookInk,line width=0.9pt] (0,0) rectangle (6,3.2);
      \draw[BookRule,line width=0.5pt] (0,0.8) -- (6,0.8);
      \draw[BookRule,line width=0.5pt] (0,1.6) -- (6,1.6);
      \draw[BookRule,line width=0.5pt] (2,0) -- (2,3.2);
      \node[fill=BookBlue!12,inner sep=2pt] at (1,2.4) {项目};
      \node[fill=BookTeal!12,inner sep=2pt] at (4,2.8) {金额（元）};
      \foreach \x/\y/\v in {
        3/2/1月,5/2/2月,1/1.2/甲,3/1.2/12,5/1.2/18,
        1/.4/乙,3/.4/8,5/.4/10}
        \node at (\x,\y) {\v};
    \end{scope}
  }
  \draw[BookMuted,dashed,line width=0.5pt] (4,0) -- (4,3.2);
  \draw[BookMuted,dashed,line width=0.5pt] (0,2.4) -- (6,2.4);
  \draw[BookBlue,line width=1.1pt] (0,1.6) rectangle (2,3.2);
  \draw[BookTeal,line width=1.1pt] (2,2.4) rectangle (6,3.2);
  \node[fill=BookBlue!12,inner sep=2pt] at (1,2.4) {项目};
  \node[fill=BookTeal!12,inner sep=2pt] at (4,2.8) {金额（元）};
  \draw[BookRule,line width=0.5pt] (13,0) -- (13,2.4);
  \draw[BookRule,line width=0.5pt] (11,2.4) -- (15,2.4);
  \draw[arr] (6.3,1.6) -- (8.6,1.6);
  \node[align=center] at (7.45,2.35) {按跨格对象\\合并};
  \node at (3,-0.6) {12个基本格};
  \node at (12,-0.6) {10个逻辑单元};
\end{tikzpicture}

  \caption{行列检测到逻辑单元的教学恢复。左侧虚线表示行列交叉形成的基本格；右侧根据两处跨格对象合并表头。表头与单元归属是结构成果，字符内容来自另行读取；图中的数值供后续问答复算。}
  \label{fig:vision-table-structure}
\end{figure*}
\FloatBarrier

推理还要解决结构冲突，例如两个跨格对象同时声称拥有同一基本格。原实现通过抑制冲突对象或调整范围，使行、列及跨格关系可以组成逻辑表；之后把OCR或PDF中带位置的文字分配给单元，输出行列索引、跨度、内容和来源框。不能只输出检测框就声称已经恢复了HTML表格。

表格结构正确与字符串正确应分别评价。若甲2月的“18”被读成“13”，结构可以完全正确而后续差值错误；若“金额”被当作只覆盖1月，所有字符正确也仍损失单位归属。原工作不覆盖跨页表格，连续两页的表头重复和行连接需要额外处理。

\subsection{表单字段的解析}
\label{sec:vision-understanding-fields}

票据的“编号”“日期”“总额”不是规则的二维表格，而是一组键与值。输入是页面图像及其OCR结果，输出是字段记录：每个字段的名称、原文、规范化值和来源框。难点在于键与值的空间关系随版式变化，同一页上还可能出现多个相似金额。

常用做法分两步。第一步复用LayoutLMv3（\ref{sec:vision-understanding-layoutlm}节）的页面表示，按文字段预测角色，例如“问题”（即键）、“答案”（即值）、“标题”和“其他”，这需要相应的字段监督并微调任务头。第二步连接键与值：可以在成对文字段的表示上另训一个关系头，也可以使用清楚写出的规则；后者不能无条件归为LayoutLMv3原有结构。

设教学票据的读取结果含“编号A017”“日期2026-04-03”“已付100.00”“总额128.50”。一种简单规则是在每个键的同行右侧寻找值，要求日期符合日期格式、金额具有数值格式，并避免同一值被互不相容的字段重复占用。这条规则可以在该版式下恢复记录，但遇到“值在键下方”的模板就需调整；学到的关系头能适应更多版式，却需要成对的连接标注。

字段结果应保留原文字、规范化值和区域来源。例如，原文“128.50”对应数值128.50及单位元；不能因另有“已付100.00”就将两者都归为总额。

\subsection{图表数据的解析}
\label{sec:vision-understanding-charts}
\label{sec:vision-understanding-chart-reading}

图表用位置、长度、颜色或角度表达数值。输入是图表图像，输出是数据表：每个数值连同它所属的系列和类别。读取图表需要反向还原这些编码，再把数值对应到标签和系列；只恢复一个高度而不知道它属于哪个月份，仍不能回答月份比较。

ChartQA使用的数据表恢复管线在ChartOCR基础上扩展\scite{vision-chartqa2022}。它先用关键点检测网络找出绘图区、条形端点、折线顶点及文字区域，用CRAFT（\ref{sec:vision-understanding-craft}节）检测并读取轴标签和图例；条形由左上和右下关键点恢复，折线由连接点恢复，饼图由扇区角度恢复。随后依据纵轴标签估计刻度，再用颜色和位置邻近关系把数值关联到横轴标签与图例。这一管线由检测、读取和几何规则组成，数值换算本身不需要学习。

在线性纵轴上，已知像素纵坐标$y_0,y_1$分别对应数值$v_0,v_1$，则位置$y$的读数为
\begin{equation}
  v(y)=v_0+\frac{y-y_0}{y_1-y_0}(v_1-v_0)\text{，}
  \qquad y_1\neq y_0\text{。}
  \label{eq:vision-chart-calibration}
\end{equation}
这里纵坐标向下增加，因此较大的数值通常对应较小的$y$。公式来自两个锚点间的线性插值，不能用于未变换的对数轴。

沿用表格中甲的金额，设图表0元刻度位于$y_0=340$，20元刻度位于$y_1=140$，两根柱顶分别为220和160像素。代入得到$0+\frac{220-340}{140-340}\times20=12$元和$0+\frac{160-340}{140-340}\times20=18$元。检测结果还需依据横轴位置对应1月、2月，再依据图例把它们归到甲系列。图\ref{fig:vision-chart-reading}突出这个“像素—刻度—记录”的联系。

\begin{figure}[htbp]
  \centering
  % 数值与table-structure-recovery相同；像素轴 y=340-10v。
\begin{tikzpicture}[x=1cm,y=1cm,font=\figurefont\small,
  arr/.style={-{Stealth[length=2mm]},BookInk,line width=0.8pt}]
  \node[anchor=west] at (0,4.8) {金额（元）};
  \draw[arr] (0.8,0) -- (0.8,4.4);
  \draw[arr] (0.8,0) -- (5.0,0);
  \foreach \y/\v in {0/0,2/10,4/20}{
    \draw[BookMuted,line width=0.5pt] (0.7,\y) -- (0.8,\y);
    \node[anchor=east] at (0.65,\y) {\v};
    \draw[BookRule,dashed,line width=0.5pt] (0.8,\y) -- (4.7,\y);
  }
  \draw[BookBlue,line width=0.9pt,fill=BookBlue!18]
    (1.65,0) rectangle (2.35,2.4);
  \draw[BookBlue,line width=0.9pt,fill=BookBlue!18]
    (3.35,0) rectangle (4.05,3.6);
  \node[anchor=south] at (2,2.5) {$y=220$};
  \node[anchor=south] at (3.7,3.7) {$y=160$};
  \node[anchor=north] at (2,-0.1) {1月};
  \node[anchor=north] at (3.7,-0.1) {2月};
  \node[anchor=north] at (2.9,-0.7) {甲系列};
  \node[anchor=west,font=\figurefont\bfseries\small] at (5.5,3.7)
    {恢复记录};
  \node[anchor=west] at (5.5,2.7) {甲，1月，12元};
  \node[anchor=west] at (5.5,1.7) {甲，2月，18元};
  \node[anchor=west,align=left] at (5.5,0.5)
    {像素纵坐标向下\\金额沿纵轴向上};
\end{tikzpicture}

  \caption{线性纵轴的教学读数。轴上标注金额，柱顶旁标注原图像素纵坐标；右侧记录同时保留系列、月份、数值与单位。两根柱的数值与图\ref{fig:vision-table-structure}中甲行一致。}
  \label{fig:vision-chart-reading}
\end{figure}
\FloatBarrier

若纵轴是对数刻度，应在线性关系成立的$\log v$上插值，再取指数恢复$v$；若柱从非零基线开始，只比较柱长又会改变读数。颜色近似、图例遮挡和折线交叉会造成系列归属错误，需要独立于数值误差检查。将“18”误读成“13”后，后续计算器仍能准确算出$13-12=1$，却无法得到真实差值6；算术正确不能修复此前丢失的视觉证据。

\subsection{图解关系的解析}
\label{sec:vision-understanding-diagrams}
\label{sec:vision-understanding-diagram-graph}

流程图中的箭头表示有向关系，位置相邻的两个框未必相连。输入是图解图像，输出是有向图：节点及其文字，边及其方向和标签。难点在于连线可能交叉、弯折或断开，节点之间的关系只能通过追踪线迹来确定，不能由相对位置推断。

一种教学组合方法先用区域检测取得节点框，在各框内识别文字；随后检测线段和箭头尖端，沿连续线迹追踪连接；再依据端点与节点边界的接触，将线迹归到源节点和目标节点。节点检测可用框监督，连接预测可用端点或关系监督；在规则线稿中，也可以明确定义阈值、线段连接和箭头形状规则而不训练模型。

设节点集合为$V=\{A,B,C,D\}$，目标输出为有向图$\mathcal G=(V,E)$。若两条连线$A\to D$和$B\to C$在中间交叉，交点处没有连接点标记，恢复结果仍应只有这两条边；按“距离最近”接线可能误造$A\to C$和$B\to D$。断线、遮盖或缺少箭头时，应保留方向未知或多个候选，不能仅由页面上下顺序决定方向。

输出记录可以写成“源节点A，边标签通过，目标节点D”，并带原始线迹和端点位置。结构检查包括是否遗漏边、方向是否反转、边标签是否属于当前连接。Pix2Struct在图解问答上的微调说明，生成式模型可以直接回答关于图解的问题，但这并不等于输出了可检查的关系图；需要显式图结构时，仍要有相应监督和验证。

\section{视觉描述与内容概括}
\label{sec:vision-understanding-description}

定位和读取保留具体证据，描述则从中选择一部分组成语言。输入是图像、区域或视频，输出是一段文字；难点在于选什么内容、选出的内容能否在观测中找到依据。语言解码沿用\ref{sec:vision-models-language-generation}节的条件生成，这里要补足的是内容如何被选择，以及选择后如何保持与来源一致。

本节按描述覆盖的对象展开。全局场景的描述概括整幅图，要在有限长度内挑出主要对象与关系；局部对象的描述只讲指定区域，要在放大细节的同时保留与周围的关系；事件与步骤的描述面对视频，要决定哪些重复和细节可以压缩，哪些次数与先后必须保留。覆盖范围越窄，描述对证据精度的要求越高。

\subsection{全局场景的描述}
\label{sec:vision-understanding-global-caption}

\term{图像描述}（image captioning）以整幅图为输入，生成一句或一段概括性的文字。街景短描述可能只需指出“几辆车停在路边”；详细描述则需要补充车辆、人物、招牌与相互位置。不同长度并不意味着观察事实不同，但会让更多属性和关系进入输出，也就增加了需要核验的内容。

生成每个词时，模型都要决定从图像的哪里取信息。两条路线对“取信息的单位”给出不同回答。基于检测区域注意力的描述先用目标检测器把图像拆成一组对象区域，再在生成每个词时对这些区域做注意力；视觉单位是对象，语言解码器是为描述任务专门训练的循环网络。基于视觉语言模型的描述用通用视觉编码器的图像块或查询表示作为视觉单位，接入预训练的大型语言模型，用指令控制描述的范围和长度。

\subsubsection{基于检测区域注意力的描述}
\label{sec:vision-understanding-bottomup}

早期带注意力的描述模型在卷积特征图的均匀网格上分配注意力。一个对象常被网格切成若干块，模型很难把注意力恰好放在整个对象上；而描述中的名词大多对应完整对象。Bottom-Up/Top-Down的主要想法是把注意力的候选从网格单元换成检测器提出的对象区域：“自底向上”指区域由图像检测产生，“自顶向下”指当前语言状态决定读取哪些区域\scite{vision-bottomup2018}。

自底向上的部分是一个Faster R-CNN检测器，骨干为在ImageNet上预训练的ResNet-101，再在Visual Genome上训练预测1600类对象和400类属性；属性预测迫使区域特征保留颜色、材质等细节。候选框的非极大值抑制取IoU阈值0.7，各类别检测结果的抑制取0.3；保留任一类别置信度超过0.2的区域，每幅图最多100个，区域数随图像复杂程度变化。原论文还报告，初步实验中直接取前36个区域效果几乎相同。每个区域的特征经平均池化为2048维向量$\bm v_i$。描述模型训练期间，这些图像特征固定不变。

经典检测器（例如\ref{sec:vision-features-dpm}节的DPM）同样先确定候选窗口、再逐个判断类别；这里的差别在于区域特征是为对象和属性分类联合学习的深层特征，并直接作为后续语言模型的视觉单位，而不只是输出一个类别标签。从\ref{chap:vision-multimodal-alignment}章的角度看，检测器把注意力的候选限制在少数对象区域上，相当于一种“硬”注意力；后面的语言模型再在这些区域上做软注意力。

自顶向下的部分由两层LSTM组成，LSTM的门控结构可参考\BookChapterRef{chap:recurrent-networks}{卷五的循环神经网络章}。第一层注意力LSTM读取三样输入：前一时刻语言LSTM的输出、所有区域的平均特征，以及前一个词的嵌入，形成当前状态$\bm h_t^{\mathrm{att}}$。它据此计算
\begin{equation}
  \begin{aligned}
    e_{i,t}&=\bm w^\top\tanh(\bm W_{\mathrm v}\bm v_i
                             +\bm W_{\mathrm h}\bm h_t^{\mathrm{att}}),\\
    \alpha_{i,t}&=\frac{\exp(e_{i,t})}{\sum_{k=1}^{K}\exp(e_{k,t})},
    \qquad
    \hat{\bm v}_t=\sum_{i=1}^{K}\alpha_{i,t}\bm v_i\text{。}
  \end{aligned}
  \label{eq:vision-caption-attention}
\end{equation}
各区域特征与注意力状态先映射到相同维数，$\bm w$再把每个组合压成标量。归一化权重形成当前视觉信息$\hat{\bm v}_t$，它与注意力状态一起输入第二层语言LSTM；后者的输出经线性层和softmax预测下一个词。这是\ref{sec:vision-alignment-attention}节软对齐的直接应用：查询是语言状态，键和值是区域特征；区别只在打分形式，这里用一个带tanh的小网络计算分数，而不是缩放点积。

原论文给出了各层的宽度，表\ref{tab:vision-understanding-updown-config}按数据流列出每层的输入与输出维数。两层LSTM的隐藏单元数$M$都是1000，注意力层的隐藏维数$H$为512，词嵌入维数$E$为1000；词表由COCO训练描述中出现至少五次的词构成，共10010个。

\begin{table}[htbp]
  \centering
  \small
  \begin{tabularx}{\linewidth}{@{}p{18mm}XXp{14mm}@{}}
    \toprule
    层 & 输入 & 计算 & 输出维数\\
    \midrule
    区域特征 & 图像 & Faster R-CNN（ResNet-101），平均池化 & $K\times2048$\\
    注意力LSTM & 上一时刻语言LSTM输出（1000）、区域平均特征（2048）、上一词嵌入（1000），拼接为4048维 & LSTM，$M=1000$ & 1000\\
    注意力 & 区域特征与注意力LSTM状态 & 式\eqref{eq:vision-caption-attention}，$\bm W_{\mathrm v}\in\mathbb R^{512\times2048}$，$\bm W_{\mathrm h}\in\mathbb R^{512\times1000}$，$\bm w\in\mathbb R^{512}$ & 2048\\
    语言LSTM & 加权视觉特征（2048）与注意力LSTM输出（1000），拼接为3048维 & LSTM，$M=1000$ & 1000\\
    输出层 & 语言LSTM输出 & 线性层与softmax & 10010\\
    \bottomrule
  \end{tabularx}
  \caption{Up-Down描述模型的逐层维数。宽度取自原论文的实现细节；拼接维数由各输入维数相加得到。}
  \label{tab:vision-understanding-updown-config}
\end{table}

对于“一个人走过停着的汽车”，预测人物相关词时可以更多读取人物区域，描述汽车时再读取车辆区域。若两区域的打分为$(\log3,0)$，权重就是$(3/4,1/4)$；这表示当前计算中两种特征的贡献，并不是“这个词有75\%的概率由该区域证明”。

训练分两个阶段。先以图像及参考描述最小化逐词交叉熵，更新两个LSTM、注意力和输出层；原设置的学习率从0.01线性降到0，共6万次迭代，批大小100。随后用\term{自评序列训练}（self-critical sequence training）直接优化描述评价指标CIDEr：对同一幅图，采样生成一句$\bm w^{\mathrm s}$，再用贪心解码生成一句$\hat{\bm w}$作为基线，两句的CIDEr分数分别为$r(\bm w^{\mathrm s})$与$r(\hat{\bm w})$，参数沿
\begin{equation}
  -\bigl(r(\bm w^{\mathrm s})-r(\hat{\bm w})\bigr)\nabla_{\bm\theta}\log p_{\bm\theta}(\bm w^{\mathrm s})
  \label{eq:vision-scst}
\end{equation}
的方向下降。采样句比贪心句得分高时，提高它的概率，反之降低。这样不必让评价指标可导，也能直接改善它；代价是模型可能学会迎合指标偏好的措辞。CIDEr的定义见式\eqref{eq:vision-understanding-cider}；它对每句采样都能给出一个分数，正好用作这里的回报。使用时依赖已生成的词前缀逐步选择区域并以束搜索解码，原设置的束宽为5，并禁止同一个词连续出现两次。检测器遗漏的对象难以通过区域路径被描述，固定词表和训练分布又限制了可以表达的内容。注意力图能展示计算读取了哪里，但不能独立证明每个生成词都正确。

\subsubsection{基于视觉语言模型的描述}
\label{sec:vision-understanding-vlm-caption}

BLIP-2、Qwen3-VL等模型把视觉表示接入已有语言模型，连接方式见\ref{sec:vision-models-blip2}节和\ref{sec:vision-models-qwen3vl}节。它们的共同结论是：视觉编码器输出的表示经连接模块投影为语言模型可读的输入，语言模型在这些输入和文字提示的条件下逐词生成，预训练使其具备通用的图像描述和指令跟随能力。

用于描述时，输入由图像和描述指令组成。可以直接使用已有检查点，按其支持的提示格式要求“一句话概括”或“详细描述”；也可以在目标描述数据上微调。BLIP-2的描述微调是一个例子：语言模型保持冻结，以“a photo of”作为生成的起始提示，更新Q-Former和图像编码器，目标是参考描述的逐词交叉熵。描述指令改变输出选择，并没有增加图像中的事实。

同一张照片要求“一句话概括”时，可以覆盖主体活动；要求“说明右侧两辆车的差别”时，必须读取局部颜色、形状和相对位置。生成后可以将内容拆成“存在两辆车”“右车为蓝色”“左车位于前方”等判断，逐项回查。若某一细节因分辨率不足不可见，应缩小描述范围，或补充局部观测。这类模型的典型失败是\ref{sec:vision-understanding-challenge-language-prior}节所说的视觉幻觉，即输出了图像中没有依据的对象或属性；语言模型的先验越强，越可能补出“街景中通常会有”的内容。

两条路线在同一街景描述任务上的差别见表\ref{tab:vision-understanding-caption-compare}。视觉语言模型把检测区域换成通用的图像块或查询表示，把专门训练的LSTM换成预训练语言模型，换来了开放词汇、可由指令控制的长度与范围，以及不依赖检测器类别表的能力；代价是计算量大得多，视觉单位不再对应可检查的对象框，生成内容与图像区域的对应也更难追溯。

\begin{table}[htbp]
  \centering
  \small
  \begin{tabularx}{\linewidth}{@{}p{14mm}XX@{}}
    \toprule
    比较项 & 基于检测区域注意力的描述 & 基于视觉语言模型的描述\\
    \midrule
    视觉单位 & 检测器给出的对象区域 & 图像块或查询表示\\
    语言解码 & 专门训练的两层LSTM & 预训练大型语言模型\\
    训练 & 描述数据上交叉熵，再优化CIDEr & 大规模预训练，可选描述微调\\
    输出控制 & 由训练数据决定的固定风格 & 指令控制长度与范围\\
    可追溯性 & 注意力权重指向对象区域 & 需另行回查区域\\
    主要失败 & 检测遗漏、词表受限 & 视觉幻觉、细节因压缩丢失\\
    \bottomrule
  \end{tabularx}
  \caption{两种全局描述路线的对照。两者输入同一幅图，输出一段描述；评价都需要同时检查覆盖与事实性。}
  \label{tab:vision-understanding-caption-compare}
\end{table}

流畅度、覆盖和事实性是不同维度。短句可以忠实却遗漏关键内容，长句可以覆盖广却加入不存在的对象。评价时先固定期望描述范围，再判断哪些必要信息被覆盖、哪些输出得不到图像支持。

\subsection{局部对象的描述}
\label{sec:vision-understanding-region-caption}
\label{sec:vision-understanding-crop-caption}

全局压缩常弱化小物体细节。局部描述的输入是一幅图和一个指定区域，输出是关于该区域对象的文字。它可以提高目标区域的有效分辨率，但对象与参照物的关系可能跨越裁剪边界，因此需要同时保留来源与必要上下文。

一种组合流程先按查询取得区域框（\ref{sec:vision-understanding-grounding}节），裁剪后记录“源图编号、原图边界、缩放比例”，再将局部图和整图按模型支持的多图接口组织为输入。若模型本来接受框或区域标记，也可以使用其原生接口。训练区域描述器需要对应的“区域—描述”记录，不能把整图描述训练无条件当作已经学过所有区域指代。

假设两辆外观相近的车分别位于左右两侧，局部裁剪显示右车有车顶行李架；整图则显示它位于另一辆车后方。描述“右侧带行李架的车停在另一辆车后面”同时使用了局部外观和全局位置。若只保留右车裁剪，“在后面”的参照就被删掉了。结果应明确当前描述属于哪幅源图的哪个对象，不能把局部图中的左、右重新当作全图关系。

核验包括区域是否选对、局部属性是否读对、引用的关系是否仍有上下文支持。增加裁剪数量会增加读取成本，也可能产生相互矛盾的局部判断，因此应记录裁剪预算，避免只比较最终文字长度。

\subsection{事件与步骤的描述}
\label{sec:vision-understanding-event-caption}
\label{sec:vision-understanding-hierarchical-caption}

视频描述的输入是一段视频，输出是按时间组织的事件或步骤文字。概括会舍弃重复和细节，但事件次数与先后可能正是后续问题需要的信息。压缩的对象和保留条件应由请求决定，不能默认所有重复事件都可合并。

可以先用\ref{chap:vision-video-recognition}章的取样与片段表示，为每段保留实际时间，再交给支持视频或交错帧输入的语言接口形成描述，例如\ref{sec:vision-models-qwen3vl}节的接口。长视频也可以先生成局部描述，再按问题组织总体概括；这是一种分层组合，摘要应仍可追溯到片段。

设教学视频在$[1,2)$秒和$[5,6)$秒各发生一次敲门，$[3,4)$秒有人讲话。问“发生了什么”可以概括为“人物两次敲门，中间说了一句话”；问“第二次敲门是否在讲话之后”则必须保留第二次事件和讲话结束位置。只存储“人物敲门并讲话”会丢掉次数与顺序，即使该短描述本身没有虚构对象。

片段合并应避免把原顺序按主题重新排列，也不能把同一事件在相邻窗口中的重复描述当成两次事件。旁白、字幕和可见动作分别标明来源：字幕出现“我开门了”，不证明镜头中已经呈现开门。保存局部到总体的对应，可以在后续问答时重新取得被概括省略的证据。

\section{视觉问答与证据推理}
\label{sec:vision-understanding-qa}

\term{视觉问答}（visual question answering, VQA）依据图像、视频或文档等视觉内容回答问题。输入是视觉内容和一个问题，输出是答案，最好同时附带支持答案的证据。它可能只需读出一个日期，也可能需要连接多个时刻的事件。难点在于答案正确并不说明证据正确：模型可能凭问题中的语言规律猜对，也可能从错误的单元读出恰好相同的数。

前几节的检索、定位和读取可以成为问答的组成步骤。本节按回答一个问题的先后环节展开：先确定问题要求什么形式的答案；再按缺少的事实获取证据；需要时把多条证据组合起来；然后以选择、抽取、单元操作或生成的方式形成答案；最后核验答案与证据是否一致。

\subsection{答案要求的确定}
\label{sec:vision-understanding-question}

同一句“多少”可能要求不同的操作。存在问题要求判断是否有目标，属性问题要求给出颜色或状态，计数问题要求定义可数对象和范围，空间问题要求确定参照系，文档问题还需明确字段和单位。对“总额是多少”，应区分读取票据上已经写出的总额，还是依据若干金额自行求和；两种请求使用的操作不同，所需证据也不同。

问答模型容易利用问题中的语言规律猜测答案，例如“天空是什么颜色”多半答“蓝色”。VQA v2为同一问题收集答案不同的相似图像，使仅凭问题作答的捷径受到限制\scite{vision-vqav22017}。对同一问题更换图像后，答案是否随相关证据改变，是有用的诊断，但需要确认新图像确实改变了答案条件。

GQA利用场景图及问题程序构造组合问题，并检查一致性、视觉依据和答案合理性\scite{vision-gqa2019}。它能够系统改变属性、关系和组合深度，但这种构造也规定了题目分布，不能代表所有真实用户问题。面向连续对话时，还要保留图像编号、当前指代和已有回答的证据范围，不能只把此前答案作为不容质疑的事实。

\subsection{证据的获取}
\label{sec:vision-understanding-evidence}

确定答案要求之后，下一步是取得回答所缺的事实。获取证据应由当前缺少什么决定：读数问题需要足够分辨率，事件先后问题需要相应时间片段，多页文档问题则可能先要确定页面。每次增加输入，都应能够说明它补上了哪项信息。

证据所在的位置决定了获取方式，本小节按两种输入条件展开。页面与图像中的证据在空间上分布，获取靠页面检索、区域定位和局部放大，主要预算是分辨率和页数。长视频与多源输入中的证据在时间上分布，还可能分属画面、声音和字幕，获取靠时间检索、记忆压缩和跨来源对齐，主要预算是帧数和上下文长度。

\subsubsection{页面与图像中的证据获取}
\label{sec:vision-understanding-evidence-composition}

对“表格中甲的2月金额是多少”，若已给清晰单页，可以直接定位目标单元并读取；若给一份长文档，需要先找到包含该表的页面；若直接提供真实数据表，页面搜索和像素读取已由外部完成。三种条件的最终答案都可能是18元，但视觉工作量不同。

一种无需新增训练的组合规则是：先按标题或内容检索候选页（\ref{sec:vision-understanding-retrieval}节），再读取表格及其行列表头（\ref{sec:vision-understanding-tables}节），按“甲、2月”定位单元，最后返回数值、单位和来源。每条中间记录保留源文件或图像编号、页码、框或时间区间、原文字和处理方式。若目标页未被召回，应补充检索；若小数点看不清，应取得更清晰的裁剪。重复要求模型“仔细思考”并没有增加这两类缺失证据。

图\ref{fig:vision-qa-evidence}将同一教学表中的差值问题拆成可核验操作。这里的选择与计算是外显组合规则，各模型只承担已经说明的定位、读取或生成职责，不假设它们内部一定沿相同步骤推理。

\begin{figure}[htbp]
  \centering
  % 教学组合规则；不表示某个VLM内部必然具有该操作轨迹。
\begin{tikzpicture}[x=1cm,y=1cm,font=\figurefont\small,
  fact/.style={draw=BookBlue,fill=BookBlue!10,line width=0.9pt,
    minimum width=44mm,minimum height=13mm,align=center},
  operation/.style={draw=BookTeal,fill=BookTeal!10,line width=0.9pt,
    minimum width=50mm,minimum height=11mm,align=center},
  arr/.style={-{Stealth[length=2mm]},BookInk,line width=0.9pt,
    rounded corners=1mm}]
  \node[anchor=west] at (0,5.7) {问题：甲2月比1月多多少？};
  \node[fact] (jan) at (2.2,4.4)
    {甲，1月：12元\\页1，行3，列2};
  \node[fact] (feb) at (7.8,4.4)
    {甲，2月：18元\\页1，行3，列3};
  \node[operation] (subtract) at (5,2.25) {$18\text{元}-12\text{元}=6\text{元}$};
  \coordinate (join) at (5,3.2);
  \draw[BookInk,line width=0.9pt,rounded corners=1mm]
    (jan.south) -- (2.2,3.2) -- (join);
  \draw[BookInk,line width=0.9pt,rounded corners=1mm]
    (feb.south) -- (7.8,3.2) -- (join);
  \draw[arr] (join) -- (subtract.north);
  \node[operation,minimum width=60mm,fill=white] (answer) at (5,0.3)
    {甲2月比1月多6元};
  \draw[arr] (subtract.south) -- (answer.north);
\end{tikzpicture}

  \caption{带来源的表格问答教学流程。两个输入事实都指向同一页、同一行及不同月份单元；差值计算保留操作数与单位。任一读数或行列归属出错，都应返回相应证据检查。}
  \label{fig:vision-qa-evidence}
\end{figure}
\FloatBarrier

\subsubsection{长视频与多源输入中的证据获取}
\label{sec:vision-understanding-video-evidence}

视频问题需要的片段可能分散在很长的历史中。可以利用\ref{chap:vision-video-recognition}章的取样、检索和MovieChat式压缩记忆获得候选，再重新读取与问题相关的时刻。压缩历史保存了便于搜索的线索，并不保证保留每次手部动作或每个短暂出现的物体。

声音、字幕和图像应各自保留时间来源。画外有人说“放好了”，画面中的手却仍在移动物体，这可能来自说话者不同、语言含义不同或音画错位。支持音画输入的模型可采用\ref{sec:vision-models-qwen-omni}节的接口，具体区间对应可按\ref{sec:vision-understanding-av-time}节检查，不能只把不同来源文本拼在一起就认定它们同时描述同一事件。

LongVideoBench把视频与字幕共同作为较长上下文，要求利用所指片段的细节回答问题\scite{vision-longvideobench2024}。解释其结果时应分别记录帧数、视频时长、字幕及候选选项，避免把由字幕直接给出的答案全部归于视觉记忆。多选题的正确答案也没有同时验证空间框和精确时间区间。

两种条件下的获取都可以写成“缺什么、补什么”的循环，差别在于补充的代价。页面与图像中，补一次裁剪或一页的代价较小，主要风险是放大后丢掉上下文；长视频中，重新读取一段需要再解码视频帧，压缩记忆又可能已经丢掉所需细节，主要风险是证据在获取前就已不可恢复，以及多个来源在时间上错位。

\subsection{证据的多步组合}
\label{sec:vision-understanding-evidence-operations}

有些问题不能由单条证据回答，例如跨图比较数量，或判断两个事件的先后。输入是多条已核对的事实，输出是由它们推出的结论；难点在于每一步操作都要有明确的输入和依据，否则中间错误会被后续步骤掩盖。一种做法是把已核对的事实交给明确的操作规则执行，而不是要求模型在一次生成中隐式完成全部推理。

例如，两张图分别可见3辆和5辆符合条件的车，先分别保存对象编号和区域，再计算差值$5-3=2$。若问题问的是“新增了几辆”，只比较数量不足以回答，还需要跨图身份对应：增加2辆同时离开1辆，会使数量只增加1辆。

对前述敲门视频，可以将事件列表组织为“敲门1：$[1,2)$；讲话：$[3,4)$；敲门2：$[5,6)$”。由第二次敲门起点5大于讲话终点4，可回答“在讲话结束之后”。若两段存在重叠，则需要改用起点、终点或主要持续时段等更明确的关系，而不能强行给出唯一先后。

这种组合可以复用模型和计算工具，不需要虚构一个新的端到端训练目标。每步保留选取对象、读取事实、操作与结果，解释才有可核验的依据。NExT-QA通过时间关系、动作解释与场景问题考察更复杂的视频问答\scite{vision-nextqa2021}。其中“为什么”的标注属于对视频动作的语义解释，不等同于通过干预识别因果效应；仅凭“讲话之后敲门”无法证明讲话导致敲门。

\subsection{答案的形成}
\label{sec:vision-understanding-answering}

证据齐备后，还要把它们转成答案。输入是问题和证据的表示，输出是一个答案；难点在于答案的形式差别很大：可能是常见词“蓝色”，可能是只在当前图像中出现的地名，可能是文档原文中的一段日期，也可能是需要计算的差值。

不同的答案形式对应不同的输出计算，四种代表模型各采用一种。LoRRA从“固定答案词表加当前图像的OCR字符串”中选择一项，解决罕见词不在词表中的问题。DocVQA基线从文档的OCR文字序列中抽取一段连续跨度，答案保留原文拼写。VisionTaPas在数据表中选择单元并执行聚合操作，答案由所选单元算出。VL-T5逐词生成答案，输出不受候选集合限制。从前到后，输出空间越来越开放，而答案与证据之间的对应越来越隐蔽。

\subsubsection{LoRRA}
\label{sec:vision-understanding-lorra}

招牌上的罕见地名可能不在训练答案词表中。此前的VQA模型把问答当作在固定答案集合上的分类，词表外的答案永远无法输出，即使文字清清楚楚写在图里。TextVQA论文为此建立了需要读取图中文字的问答数据集，在OpenImages的28408幅图像上收集了45336个问题；并提出LoRRA，其主要想法是把当前图像的OCR字符串作为动态候选加入答案空间，使模型可以直接指向刚刚读到的内容\scite{vision-textvqa2019}。

模型有三个部分，结构见表\ref{tab:vision-understanding-lorra-config}。问题截取至多14个词，由GloVe词向量（词表取VQA 2.0问题中约7.7万个高频词）和带自注意力的LSTM编码为向量$\bm q$。视觉部分沿用VQA模型Pythia：网格特征取自预训练ResNet-152的res-5c块，经平均池化为2048维；区域特征取自在Visual Genome对象与属性上训练的Faster R-CNN的fc6层，训练时微调fc7层。两种特征在问题条件下分别做注意力加权平均，再拼接。读取部分先用外部OCR系统Rosetta取得至多50个OCR词，以FastText嵌入表示，使词表外的词也有向量；再以问题为条件对这些OCR词做注意力。外部OCR不与回答模块共同训练。论文没有给出注意力层和融合层的隐藏维数。

\begin{table}[htbp]
  \centering
  \small
  \begin{tabularx}{\linewidth}{@{}p{14mm}XX@{}}
    \toprule
    部分 & 输入与编码 & 输出\\
    \midrule
    问题 & 至多14个词；GloVe词向量，带自注意力的LSTM & 问题向量$\bm q$\\
    视觉 & ResNet-152 res-5c网格特征（2048维）与Faster R-CNN fc6区域特征；以$\bm q$为条件的注意力 & 两种加权特征拼接，与$\bm q$逐元素相乘\\
    读取 & Rosetta识别的至多50个OCR词；FastText嵌入；以$\bm q$为条件的注意力 & 加权特征与50个注意力权重拼接，与$\bm q$逐元素相乘\\
    答案 & 两路结果拼接；两层前馈网络 & $3996+50$个logit\\
    \bottomrule
  \end{tabularx}
  \caption{LoRRA的结构，依据TextVQA论文。固定答案词表含训练集中至少出现两次的3996个答案；隐藏维数论文未给出。}
  \label{tab:vision-understanding-lorra-config}
\end{table}

视觉与读取两路各自与问题表示融合，融合方式是逐元素相乘。按\ref{chap:vision-multimodal-alignment}章的分类，这是中间融合（\ref{sec:vision-alignment-intermediate-fusion}节）中最简单的乘性交互：问题向量的每一维像一个开关，放大或抑制对应维度的视觉特征，代价是只能表达同一维度之间的交互，不能像双线性融合那样让任意两维相互作用。两路融合结果拼接后送入答案模块。设固定答案数为$n$，当前OCR候选数为$m$，答案模块输出$n+m$个logit：前$n$个位置对应固定答案，后$m$个对应当前图像的OCR字符串，选中第$n+j$个位置就复制第$j$个字符串。读取部分的注意力只给出加权平均后的OCR特征，会丢失“哪个位置的词得了高分”这一信息；原实现因此把按原候选顺序排列的注意力权重也拼进特征，使答案模块知道应复制第几个候选。

训练使用逐候选的二元交叉熵，而不是要求唯一位置的softmax交叉熵。原因是同一答案可能同时出现在固定词表和OCR候选中，两个位置都应获得正监督。若固定词表没有“青山路”，OCR已正确读出它，模型仍可通过动态位置回答；若OCR把它切成“青山”和“路”两个候选，原LoRRA一次只能复制一个候选，不能输出多词答案；原论文统计，TextVQA中有8.9\%的问题只能由多个OCR词组合回答。OCR读错时，复制机制会原样复制错误。训练使用AdaMax优化器，批大小128，共2.4万次迭代，评价用\ref{sec:vision-understanding-challenges}节定义的VQA准确率。

\subsubsection{DocVQA基线}
\label{sec:vision-understanding-docvqa}

文档问答的答案往往已经出现在原文中，例如日期、编号或机构名。从固定词表中选择不适合这类答案，复制单个OCR词又不够表达多词答案。DocVQA数据集在12767页文档图像上收集了5万个问题，主要指标是ANLS（式\eqref{eq:vision-understanding-anls}）\scite{vision-docvqa2021}。论文同时测试了LoRRA和另一个文字问答模型M4C：原设置的LoRRA测试ANLS约为0.11，去掉固定答案词表后大幅下降，说明它主要依靠固定词表作答；M4C在动态OCR候选增加到500个时约为0.39，去掉检测器给出的对象特征后反而略有提高，说明为自然图像设计的对象区域在文档上用处不大。二者都不如下面的阅读理解基线。论文中表现最好的文字基线，主要想法是借用阅读理解中的\term{抽取式问答}（extractive question answering）：答案是上下文中的一段连续文字，模型只需预测这段文字的起点和终点，答案因此保留原有拼写。

基线把商业OCR软件识别出的词按从左上到右下的顺序、以空格连接成一个字符串作为上下文，与问题拼接后送入BERT编码。论文比较了三个预训练检查点：bert-base-uncased，以全词掩码预训练的bert-large，以及在其基础上又在SQuAD 1.1上微调过的bert-large-squad。BERT-base为12层、宽度768、12头，约1.1亿参数；BERT-large为24层、宽度1024、16头，约3.4亿参数\scite{devlin2019}。每层都是后置层归一化的标准Transformer编码器块。DocVQA论文没有给出最大序列长度和长文档的滑窗步长。设上下文词元表示为$\bm h_j$，两个线性头分别给出开始分数和结束分数，在全部上下文位置上做softmax，得到开始概率$p_{\mathrm s}(j)$和结束概率$p_{\mathrm e}(j)$。它利用的是OCR文字，不能因此被描述为BERT直接从像素理解布局。

对训练跨度$(s^*,e^*)$，目标为
\begin{equation}
  L_{\mathrm{span}}=-\log p_{\mathrm s}(s^*)-\log p_{\mathrm e}(e^*)\text{。}
  \label{eq:vision-docvqa-span}
\end{equation}
表现最好的配置采用已在SQuAD阅读理解数据上微调过的BERT-large，再在DocVQA上继续微调6轮，批大小8，学习率$2\times10^{-5}$；它在测试集上的ANLS为0.665，准确率为55.77\%。不经DocVQA微调、直接使用bert-large-squad时，测试ANLS只有0.475，说明阅读理解能力需要适应文档OCR文字的分布。数据只给答案字符串，原论文沿用TriviaQA的做法，用答案在序列中的首次匹配位置近似构造训练跨度；重复的日期或编号可能使这一弱标注指向错误位置。使用时在满足$s\leq e$和长度上限的组合中选择$p_{\mathrm s}(s)p_{\mathrm e}(e)$最大的跨度，再由词元到OCR字符串的映射恢复答案。

在教学票据中，序列包含“编号A017日期2026-04-03已付100.00总额128.50”，问题“票据日期是什么”应指向日期对应的跨度。若OCR已把03读成08，抽取模块即使选对跨度也会输出错误日期；若双栏顺序错乱，连续跨度又可能混入另一栏文字。要在结果中返回来源，可额外保留OCR词元到框的映射；LayoutLMv3等页面模型（\ref{sec:vision-understanding-layoutlm}节）以同样的跨度目标微调，只是把布局和图像也纳入了编码。

\subsubsection{VisionTaPas}
\label{sec:vision-understanding-visiontapas}

图表问题常需要计算，例如“甲的2月比1月多多少”。抽取式和选择式答案都无法表达一个不在页面上的差值。ChartQA数据集包含约九千六百个人工撰写的问题和约两万三千个由图表摘要生成的问题，其中人工问题大量涉及比较与算术\scite{vision-chartqa2022}。该论文提出的VisionTaPas的主要想法是沿用表格问答模型TaPas的做法：答案由“选哪些单元”和“对它们做什么聚合”两部分决定，模型只需预测这两部分，计算交给确定的规则；同时加入图表图像，使问题提到颜色或形状时也能利用视觉信息。

模型有三个编码部分，配置见表\ref{tab:vision-understanding-visiontapas-config}。图表图像切成图像块，经线性投影并加上一维可学习位置嵌入，送入12层的ViT-Base编码器，以ViT原论文发布的权重初始化\scite{dosovitskiy2020vit}；ChartQA论文没有给出输入分辨率和图像块大小。问题与按行展平的数据表以“[CLS] 问题 [SEP] 表格”的格式送入TaPas编码器。TaPas在BERT的基础上为每个词元加入所在行、列、数值排序以及上一轮答案等嵌入\scite{vision-tapas2020}，这里使用在WikiTQ表格问答数据上微调过的TaPas-Base权重；数据表可以是真实表，也可以由\ref{sec:vision-understanding-charts}节的管线从图表恢复。

两路输出进入四个跨模态块，每块包含视觉分支和文字表格分支。设第$\ell$块的视觉输入为$\bm H^{(\ell)}$，文字表格输入为$\bm Z^{(\ell)}$，视觉分支的计算顺序为
\[
 \begin{aligned}
  \bm H'&=\bm H^{(\ell)}+\operatorname{CA}(\bm H^{(\ell)},\bm Z^{(\ell)}),\\
  \bm H''&=\bm H'+\operatorname{MSA}(\bm H'),\\
  \bm H^{(\ell+1)}&=\bm H''+\operatorname{FFN}(\bm H'')\text{，}
 \end{aligned}
\]
其中$\operatorname{CA}(\bm A,\bm B)$表示以$\bm A$为查询、以$\bm B$为键和值的多头交叉注意力；每个子层都带层归一化。文字表格分支对称地以$\bm Z^{(\ell)}$为查询读取$\bm H^{(\ell)}$。两个分支的交叉注意力并行计算，读取的都是上一块的输出。这是\ref{sec:vision-alignment-intermediate-fusion}节中基于交叉注意力的中间融合：图像块可以根据表格中提到的系列名调整表示，表格单元也可以读取对应柱子的颜色和位置。末端在文字表格分支上接TaPas原有的两个头：单元选择头为每个单元给出被选概率，聚合头在若干操作中选择一个。

\begin{table}[htbp]
  \centering
  \small
  \begin{tabularx}{\linewidth}{@{}p{16mm}XX@{}}
    \toprule
    组件 & 配置 & 初始化\\
    \midrule
    视觉编码器 & ViT-Base，12层；分辨率与块大小论文未给出 & ViT原论文权重\\
    文字表格编码器 & TaPas-Base，12层，宽度768 & 在WikiTQ上微调的TaPas\\
    跨模态编码器 & 4块，每块两个分支：交叉注意力、自注意力、前馈层 & 论文未说明\\
    输出头 & 单元选择头；聚合头（无聚合、求和、计数、平均、差值、比值及固定答案类） & 前四种操作沿用TaPas\\
    训练 & ChartQA上30轮，学习率$10^{-5}$，批大小64 & —\\
    \bottomrule
  \end{tabularx}
  \caption{VisionTaPas的配置，依据ChartQA论文正文与附录。TaPas-Base的宽度沿用BERT-base；跨模态块的头数与宽度论文未给出。}
  \label{tab:vision-understanding-visiontapas-config}
\end{table}

原TaPas的聚合操作有无聚合、求和、计数和平均。ChartQA增加了差值和比值，并把“是”“否”等固定答案作为额外类别。设被选中的数值为$v_1,\ldots,v_k$，选择“求和”时返回$\sum_i v_i$，选择“计数”时返回$k$，选择“无聚合”时直接输出单元内容，选择“差值”时返回两个所选数值之差。对于“甲的2月金额比1月多多少”，正确选择是甲行的18和12，并选择差值，得到$18-12=6$元。单元选择指向具体来源，使人能够发现模型是否误用了乙行的10。

训练时，原TaPas可以只用最终答案作弱监督，让模型自己找到能算出答案的单元和操作。对新增的差值和比值，ChartQA改用启发式方法从答案和数据表直接推出候选单元：已知答案6，就在表内寻找差为6的数对。可能有多组数对满足同一差值，原论文人工检查一百例，发现约四分之一的这类监督有误。使用时选取概率超过阈值的单元，执行聚合头选出的操作。若请求“甲乙2月之和比1月之和多多少”，需要$(18+10)-(12+8)=8$元，这超出了单次聚合的表达，不能仅凭操作集中同时有“求和”和“差值”就认定支持嵌套计算。

\subsubsection{VL-T5}
\label{sec:vision-understanding-vlt5}

VisionTaPas的答案受限于操作集合，不能输出解释或任意格式的文字。VL-T5原本是一个把多种视觉语言任务统一为文字生成的模型：在T5编码器—解码器上，把图像区域特征和问题文字一起作为编码器输入，由解码器逐词生成答案\scite{vision-vlt5-2021}。ChartQA用它作为生成式基线，主要想法是不规定答案形式，让模型直接写出答案。

VL-T5以T5-Base为骨干\scite{raffel2020t5}，编码器和解码器各12层，宽度768，约2.2亿参数。原模型用在Visual Genome上训练的Faster R-CNN为每幅图取36个区域，每个区域的嵌入是四项之和：区域特征经线性层投影，框坐标经线性层投影，图像编号嵌入（多图输入时区分来自哪幅图），以及区域编号嵌入；区域编号嵌入与文字词表中新增的36个视觉哨兵词元共享，使模型可以用文字指称某个区域。文字词元与区域嵌入拼接成一个序列送入编码器，编码器块由自注意力和全连接层组成，解码器块另加交叉注意力。原模型在COCO与Visual Genome的约18万幅图像、918万个图文对上做多任务预训练。

文字输入由问题与序列化数据表组成。视觉输入沿用原VL-T5的接口，是一组对象区域特征及其框坐标；但自然图像检测器的对象类别不适合图表，ChartQA另以ResNet-101为骨干、以COCO权重初始化，训练了识别条形、折线、扇区、图例和坐标轴标签等15类图表元素的Mask R-CNN，提取这些元素的区域特征，数量不足时补零到固定的36个。编码器联合编码文字与区域特征，解码器根据多模态上下文及已有答案前缀预测后续词元。训练以目标答案的逐词交叉熵更新参数，在ChartQA上微调30轮，批大小96，学习率$10^{-4}$；使用时以束宽5的束搜索解码。评价采用\ref{sec:vision-understanding-challenges}节所说的放宽准确率。

对于前述差值问题，目标字符串可以是“6”；模型的输出空间不受某个聚合操作集合限制，但内部生成“6”的过程也不暴露18、12和减法。开放表达与准确计算是两种性质：生成接口可以输出完整解释，却仍可能在已经给对的数据上算错。判断视觉特征是否有贡献，还应固定数据表条件：若表头与数字足以独立回答，正确生成并不证明模型利用了图表图像；问题引用“蓝色柱”时，才可能需要颜色与系列的视觉对应。

四种答案形式的对照见表\ref{tab:vision-understanding-answer-compare}。从LoRRA到DocVQA基线，把“复制单个OCR词”换成“抽取连续跨度”，换来了多词答案，代价是依赖OCR序列化顺序。从抽取到VisionTaPas的单元选择与聚合，答案第一次可以是页面上没有的数，代价是只能表达预设的单步操作。VL-T5放开了输出形式，代价是答案与证据之间不再有显式对应。可追溯性与表达范围此消彼长，选择时应看任务更需要哪一项。

\begin{table}[htbp]
  \centering
  \small
  \begin{tabularx}{\linewidth}{@{}p{12mm}XXXX@{}}
    \toprule
    比较项 & LoRRA & DocVQA基线 & VisionTaPas & VL-T5\\
    \midrule
    答案形式 & 选择一项 & 连续跨度 & 单元加聚合 & 自由生成\\
    输出空间 & 固定词表与OCR词 & 上下文中的跨度 & 所选单元的运算结果 & 任意词元序列\\
    文字来源 & 外部OCR & 外部OCR & 数据表 & 数据表\\
    证据追溯 & 指向OCR候选 & 指向跨度位置 & 指向所选单元 & 无显式对应\\
    主要局限 & 只能复制一词 & 依赖序列化顺序 & 只有单步操作 & 计算过程不可见\\
    \bottomrule
  \end{tabularx}
  \caption{四种答案形式的对照。前两者原用于场景文字和文档问答，后两者用于同一ChartQA设置；比较时应固定文字来源与评价口径。}
  \label{tab:vision-understanding-answer-compare}
\end{table}

\subsection{答案的核验}
\label{sec:vision-understanding-answer-check}

答案形成后，还要确认它与证据和问题一致。答案可拆成对象、属性、字符串、数值、单位和关系等判断，分别与来源对照。若输出“总额是128.50美元”，原票据只写人民币符号，则数字正确而单位错误；若正确数值来自另一页的同名表格，来源也不成立。

核验可以改变相关观测后重新作答。例如，把甲2月金额从18改成13，差值应从6变成1；保持图像不变而改问乙行，答案应改为2元。前一种检查对输入变化的响应，后一种检查查询绑定，二者都需要重新确定正确答案，不能只观察模型是否“回答不同”。

缺少证据时回到相应获取步骤；来源冲突时保留冲突并限制结论。多轮对话中，“那一行”“他后来怎样”应绑定到已经明确的单元或对象，新增裁剪、页面或视频帧则登记为新的观测。上一轮回答本身不应替代尚未核对的原始事实。

\section{理解结果的核验}
\label{sec:vision-understanding-evaluation}

理解系统可能答对而找错证据，也可能位置正确但读错字符。验收要回答的问题是：最终结果之外，系统依据的中间结构、来源和输入是否都可靠。一个总分把这些情况混在一起，无法指出错误发生在哪里，也无法比较信息条件不同的系统。

本节按核验的维度展开，三个维度分别对应前文的不同层次。中间结构的核验检查检索候选、定位结果、字符串和表格等成果本身是否正确。来源证据的核验检查最终答案或描述是否真正由所报告的区域和时刻支持。模态贡献的核验检查答案依赖的是哪一种输入，以及比较的系统是否在相同的信息条件下工作。三者可以叠加使用，任一维度通过都不能代替其他维度。

\subsection{中间结构的核验}
\label{sec:vision-understanding-output-metrics}
\label{sec:vision-understanding-grounding-output}

匹配评价应固定候选单位，报告召回和排序；定位要检查区域、点或区间是否满足查询；读取和结构解析则应拆开字符、行列、阅读关系及连接错误。描述与问答还需核验事实覆盖和答案要求。表\ref{tab:vision-understanding-evaluation}给出沿本章例子的对应。

\begin{table}[htbp]
  \centering
  \small
  \begin{tabularx}{\linewidth}{@{}p{19mm}XX@{}}
    \toprule
    成果 & 直接检查 & 仍需单独检查\\
    \midrule
    检索候选 & 完整正例的召回和名次 & 关系是否由同一对象满足\\
    框或区间 & 与目标范围的吻合 & 指代、身份与空目标\\
    OCR字符串 & 字符编辑距离、整串正确 & 数字含义、单位与来源\\
    表格结构 & 行列、跨度及内容归属 & 单元数值和后续运算\\
    描述或答案 & 内容、数值与请求一致 & 支持它的区域和时刻\\
    \bottomrule
  \end{tabularx}
  \caption{多模态理解的分层核验。每行的直接检查只覆盖一种成果，不能自动证明最后一列的条件。}
  \label{tab:vision-understanding-evaluation}
\end{table}
\FloatBarrier

定位结果要能被复核。一个可复核的定位结果至少包含查询、源图或页面编号、时间信息、坐标约定及输出范围。框可以继续作为分割提示，点可以交给界面执行器，但新的输出要按新的标准检查。返回多个合理候选时，应说明差别来自参照不明、重复实例，还是目标部分遮挡；这种歧义反馈本身就是结果的一部分，不应被强行压成唯一答案。视频中“刚才拿杯子的人”还要求身份跨帧延续。可以复用\ref{sec:vision-video-tracking}节的对应和存在状态，区分进入、持续可见与离开；离开画面的对象不应被强行重新绑定到一个相似人物。2026年公开的TempoGround将跨帧对应、对象存在及二维和相机坐标系中的三维框联系起来，可作为这一问题的扩展案例\scite{vision-tempoground2026}。其二维定位、身份保持和三维几何仍是分别受约束的结果，论文设置不构成任意视频的长期一致性保证。

检索与定位中的关系绑定需要专门的诊断。Winoground固定两句文字的词汇集合，只改变词序，并要求与两幅图正确配对，可用于检查角色和关系绑定\scite{vision-winoground2022}。前述“人搬起箱子”和“箱子挡住人”也体现了相同诊断思路：不能只检测是否出现所有名词。基准论文中对当时模型的观察应限于所测设置，不外推为所有新模型的结论。

字符串的读取用字符错误率（式\eqref{eq:vision-cer}）衡量，并需事先规定空格、标点和大小写的处理。将“128.50”读成“123.50”只有一次替换，CER为$1/6$，但金额字段已经错误，因此业务字段还应报告整串或数值正确性。页面级解析则应分项报告文字、版面、表格和公式的误差，例如OmniDocBench的模块评价（\ref{sec:vision-understanding-page}节），并按页面类型保留结果，以发现报纸密集跨栏和论文公式等不同困难。

\subsection{来源证据的核验}
\label{sec:vision-understanding-grounded-evaluation}

答案正确不等于依据正确。来源证据的核验要求答案与它所报告的区域和时刻同时成立。若要求答案、时间和空间都正确，可以把单题成功写成三个条件的交集。令$A_i$表示答案正确，$T_i$表示返回时间区间满足阈值，$B_i$表示目标框满足阈值，则在明确标注协议下，联合成功率为
\begin{equation}
  \frac1n\sum_{i=1}^{n}\bm 1[A_i\land T_i\land B_i]\text{。}
  \label{eq:vision-grounded-success}
\end{equation}
这里$\bm1[\cdot]$为指示函数。这个教学指标的目的在于展示交集关系，不代称某个基准的全部计分规则。100题中若80题答案正确，其中50题时间也正确，这50题中又只有30题空间正确，联合成功率就是30\%，不是把三个环节的边际分数任意相乘。

2026年公开的VideoZeroBench通过不同层级分别收紧答案、时间和空间证据要求，具体体现了答案与依据可能分离的问题\scite{vision-videozerobench2026}。报告时还需说明模型是端到端寻找证据，还是已获得人工给定区间或框；额外给定证据改变了任务难度。该预印本的测试范围不能代表所有长视频应用。

描述和答案中不能出现无依据的内容。POPE以对象是否存在的询问检查视觉幻觉\scite{vision-pope2023}。应同时观察错误肯定和错误否定，避免总回答“没有”获得表面的保守。存在性正确仍不保证颜色、数量、动作、字符串或长段自由描述都正确，这些判断需要各自证据。空间判断还要保留观察视角：VSI-Bench考察视频中的空间信息保存与推断\scite{vision-vsibench2025}。视角改变后，“左侧”可能需要从当前相机坐标转换到场景关系；只记住一句“椅子在桌子左边”可能丢失当时参照。

% 完整作者边注不可跨页，为论述开头和来源共同预留阅读范围。
\Needspace{11\baselineskip}
解释长度也不是证据量。More Thinking, Less Seeing?研究了所测多模态推理模型中推理长度与视觉忠实性的关系，并提出相应诊断\scite{vision-morethinking2025}。由此需要检查的是延长输出后，新增判断是否仍有视觉依据；不能概括成“解释越长一定越好”或“一定越差”。若模型多写了五段推理，却始终使用读错的13，更多文字仍不能把正确差值恢复为6。

\Needspace{14\baselineskip}
Video-MME-v2用关联问题的成组评价补充单题准确率，使相互矛盾的回答更难被总分掩盖\scite{vision-videommev22026}。例如，模型既说“第二次敲门在讲话后”，又说“讲话发生在两次敲门之后”，在上述事件记录下至少一项错误。成组一致性是必要诊断，却不能独立证明内容正确：模型也可能一致地记错所有时间。其2026年版本应与所用字幕、音频和帧采样条件一起记录。

\subsection{模态贡献的核验}
\label{sec:vision-understanding-input-protocol}

多模态系统答对一个问题，未必用到了视觉。要判断视觉是否贡献了答案，可以移除、替换或错配图像，同时固定问题和其他输入。相同方法也适用于声音、字幕、OCR和数据表。移除模态可能产生训练分布之外的输入，性能变化不宜直接解释为纯粹因果贡献；多种受控替换与具体失败检查能提供更可靠的诊断。\ref{sec:vision-alignment-missing}节讨论过模型在推理时缺少某个模态怎样预测，\ref{sec:vision-alignment-imbalance}节讨论过联合训练中强势模态压制弱势模态的现象；这里的检查正是在评价端暴露这两种情形：若去掉图像后答案几乎不变，说明视觉在当前任务上几乎没有被利用，这与\ref{sec:vision-understanding-challenge-language-prior}节所说的语言捷径一致。

输入协议决定了系统实际承担的工作。若字幕已经明确写出答案，无图条件下答对并不意外。若真实数据表包含全部行列和数字，系统就不再承担从图像恢复表格的工作：ChartQA中同一模型在真实数据表和恢复数据表两种条件下的结果，衡量的是不同任务。报告需把原始像素、外部OCR、字幕、真实结构以及检索裁剪预算逐项固定，才能比较两个系统是否在相同信息条件下工作。

\section{本章小结}
\label{sec:vision-understanding-summary}

开篇的三个请求都要求把一句话落到特定观测上。本章按证据由粗到细的顺序给出了回答：检索先用双编码器从大库中召回候选，再以配对判定或逐条件核验重排；定位把文字短语、声音、音画事件和界面指令分别对应到区域、时间和点击位置；读取与解析恢复字符串，再按页面、表格、表单、图表和图解各自的结构组织关系；描述从中选择内容，问答则经过答案要求的确定、证据的获取与组合、答案的形成和核验。“右边拿杯子的人”由文字定位给出框，“总额是多少”由读取、字段解析和抽取给出带来源的数值，“第二次敲门在讲话之前还是之后”由时间定位与事件列表的比较给出。

第一节提出的五项挑战在各节分别得到处理。多粒度对应由召回、重排和四种定位分层解决；弱配对与噪声监督由间隔目标、困难负例和带置信度的伪标注缓解，但伪标注的自我强化仍是风险；高分辨率与长输入通过可变分辨率、裁剪搜索和按需获取证据处理，代价是更多的计算或调用；证据组合可以交给外显的操作或结构化的答案头，代价是表达范围受限；语言先验没有被某个模型消除，只能通过保留证据和分层核验暴露出来。贯穿各节的比较显示，同一个设计选择常同时带来收益和代价。双编码便于召回，却看不见关系绑定；联合计算能判断关系，却只能用于少量候选。无OCR的页面模型省去外部读取，却要求像素中保留小字；生成式答案表达自由，却不再暴露所依据的单元。可检查的结果应同时保留内容和来源：正确数字可能取自错误单元，正确描述可能遗漏关键时序，流畅解释也可能围绕错误观察展开。按中间结构、来源证据和模态贡献分层核验，才能确定问题发生在哪里。

仍未解决的问题集中在证据的完整性与组合推断上。关键证据可能在压缩或缩放时已经丢失，关系绑定错误与多步空间推断相互影响。Video-MSR在2026年将约束定位、链式指代、路线及反事实物理问题组织为多跳空间推理测试，可用于区分单步识别和组合推断\scite{vision-videomsr2026}。基准内训练带来的改善仍需放回具体任务评估，不能直接推论真实导航可靠性。Geo3R讨论了透视、物体朝向和视角变化引起的空间回答偏差，并在推理时引入几何信息及结构化三维关系\scite{vision-geo3r2026}。作为2026年公开案例，它说明可把空间判断重新落到几何对象上检查；若几何估计本身错误，增加三维表示也可能传播错误，而不是普遍消除视觉幻觉。

这些方向共同要求保留可核验的中间对象。读取越多并不总是越有用，关键在于缺少哪项事实、增加的观测是否补上它，以及不同来源冲突时怎样限定答案。下一章转向多模态生成：依据给定条件创作新内容时，需要保持的主体、空间和时间关系仍须明确，但未观测细节的生成与本章对既有证据的解释承担不同职责。

\paragraph{扩展阅读}
在大规模对比预训练普及之前，视觉语言理解多以检测区域特征加Transformer预训练的方式完成，可与本章基于检测区域注意力的描述模型和VL-T5对照阅读。ViLBERT用两路Transformer分别处理文字和区域，再以共同注意力层交换信息；LXMERT同样分设语言、对象关系和跨模态三个编码器，并加入掩码对象预测和图像问答等预训练任务；UNITER改为把文字和区域放入同一个Transformer，并用最优传输显式约束词与区域的对齐\scite{vision-vilbert2019,vision-lxmert2019,vision-uniter2020}OFA进一步把定位、描述、问答和图像生成都写成带指令的序列到序列任务，区域坐标也离散成词元输出，与VL-T5把任务统一为文字生成的思路一致，但覆盖的任务更广\scite{vision-ofa2022}

定位与读取方向，GLIP把目标检测改写成短语定位，用检测、定位和带伪标注的图文数据共同预训练，可与MDETR和Grounding DINO比较怎样扩展开放词表；Kosmos-2让多模态大语言模型在生成的文字中直接插入位置词元，把描述中的短语链接到图像区域\scite{vision-glip2022,vision-kosmos2-2023}Donut不调用外部OCR，由图像编码器和文本解码器直接从页面生成结构化输出，与本章先用TrOCR读取、再交给LayoutLMv3解析的流程正好相反，与Pix2Struct同属无OCR路线；Nougat沿用这种结构，把学术论文页面直接转换为包含公式和表格的标记文本\scite{vision-donut2022,vision-nougat2023}

问答与核验方向，OK-VQA收集了只靠图像内容无法回答、必须调用外部知识的问题，补充了VQA v2和GQA不考察的知识来源；EgoSchema用约三分钟的第一人称视频构造多项选择题，并以回答所需的最短视频长度衡量问题对长时证据的依赖，可与NExT-QA和LongVideoBench对照\scite{vision-okvqa2019,vision-egoschema2023}CHAIR在生成的图像描述中统计提到但图中并不存在的物体比例，是POPE之前常用的对象幻觉指标；POPE改用是非询问，两者分别从自由描述和受控提问两端暴露同一类错误\scite{vision-chair2018}

\chapter{多模态生成}
\label{chap:vision-multimodal-generation}

“让玩具狗绕花瓶走一圈，碰到花瓶时发出轻响”看似只有一句要求，交付时却要同时回答几个问题。玩具狗在转身后是否仍是原来的样子，花瓶有没有真的受到碰撞，声音是否发生在接触时刻？如果还要把玩具狗做成可以转动观察的三维资产，背面和侧面也必须共同解释同一个对象。单幅画面自然，只完成了其中一部分。

这句请求的困难来自两类关系。一类在单个输出内部：视频的各帧要描述同一个持续存在的对象，三维内容的各个视角要描述同一个表面。另一类在不同输出之间：声音要落在画面事件上，说出的回答要与看到的证据一致。读完本章，读者应能说明每种成果需要哪些条件、模型怎样把这些条件接入生成计算、怎样把生成结果解码成可交付的媒体，以及怎样核验上述两类关系。

要回答这些问题，先得知道“做得好不好”怎样度量，以及难点具体落在哪里。本章第一节把开篇请求拆成几项可以度量的挑战，并说明经典方法和学习方法分别能处理到哪一步；之后各节按“哪些内容已经固定、哪些内容由模型生成”逐步改变设定，前一节固定的内容在后一节变成可变量，这就是各节相连的原因。阅读本章需要第一部分的视频运动与时间连续性度量、多模态对齐与融合以及相机几何，也需要\ref{chap:vision-multimodal-models}章的生成模型接口。

生成过程沿用\ref{sec:vision-models-media-generation}节的潜表示、条件连接和采样接口，流匹配的速度目标见\ref{sec:vision-models-flow-matching}节，扩散模型的前向加噪与反向去噪可参考\BookChapterRef{chap:pgm-generative}{卷六“生成模型”章}。下文用$s\in[0,1]$表示流匹配的生成路径参数，用$\ell$表示扩散模型的噪声水平，用$t$或$\tau$表示媒体中的时间，避免把“第几次去噪”与“视频第几秒”混在一起。单幅图像的创建、编辑与观测场景的新视角合成见\ref{chap:vision-image-generation}章；本章增加跨时间、跨输出和跨视角的约束。

统一模型也要按实际输出入口使用。例如\ref{sec:vision-models-showo2}节的Show-o2可以在共享主干中形成文字与视觉输出，但视频能力取决于具体变体的训练范围，也没有因此获得音频输出。共享主干和多模型组合都要接受同样的成果核验，本章因此按成果而不是按主干组织。

\section{多模态生成的问题与挑战}
\label{sec:vision-multigen-challenges}

多模态生成的输入是创作条件，输出是一件或几件媒体成果。输入可以是文字；可以再加首帧、参考图或几帧已知画面；可以是待修改的原片段及修改指令或掩码；可以是一段已经固定、等待配声的视频；也可以是一组带相机参数的观测，或者用户对一段视频提出的问题。输出相应地是帧序列、与视频等长的声轨波形、带声轨的视频、可以从任意相机渲染的三维表示（辐射场、网格或高斯集合），以及文字与语音形式的回答。与\ref{chap:vision-multimodal-understanding}章的理解任务相比，这里最大的不同是同一请求有许多合法输出：“玩具狗绕花瓶走一圈”可以从左边绕，也可以从右边绕，可以走得快些或慢些。评价因此很少能把输出与唯一的参考逐像素比较，常用的度量分成三类：把一批生成结果与一批真实媒体在特征分布上比较，把每个结果与它的输入条件比较，以及检查一件成果内部或几件成果之间的关系。下面先给出本章用到的主要指标，再逐项讨论挑战。

视频的分布级指标中最常用的是\term{Fréchet视频距离}（Fréchet video distance, FVD）。单个生成视频没有唯一的正确答案，但一个好的生成模型产生的一批视频，其总体统计应当接近同类真实视频。FVD先把每段视频送入在Kinetics-400动作识别数据上训练的I3D网络（结构见\ref{sec:vision-video-i3d}节），取末层输出作为该视频的特征向量；再把真实视频的特征和生成视频的特征各自拟合成多元高斯分布，均值和协方差分别为$(\bm\mu_{\mathrm r},\bm\Sigma_{\mathrm r})$与$(\bm\mu_{\mathrm g},\bm\Sigma_{\mathrm g})$；最后计算两个高斯之间的Fréchet距离\scite{vision-fvd2018}计算式为
\begin{equation}
  \mathrm{FVD}
  =\lVert\bm\mu_{\mathrm r}-\bm\mu_{\mathrm g}\rVert_2^2
   +\operatorname{Tr}\!\left(
     \bm\Sigma_{\mathrm r}+\bm\Sigma_{\mathrm g}
     -2(\bm\Sigma_{\mathrm r}\bm\Sigma_{\mathrm g})^{1/2}
    \right)\text{。}
  \label{eq:vision-multigen-fvd}
\end{equation}
这个式子是两个多元高斯之间2-Wasserstein距离的平方。2阶Wasserstein距离的平方是把一个分布“搬运”成另一个分布所需的最小平均平方位移；其定义及多元高斯的闭式可参考\BookChapterRef{chap:information-theory-and-statistical-geometry}{卷2《数学准备：概率、因果与决策》的“概率分布的分析与优化”章}。本章把这一闭式用于视频特征。第一项衡量两批特征的中心偏了多远；第二项衡量它们的散布形状差多少。矩阵$\bm\Sigma_{\mathrm r}\bm\Sigma_{\mathrm g}$一般不对称，它的平方根按$\operatorname{Tr}[(\bm\Sigma_{\mathrm r}\bm\Sigma_{\mathrm g})^{1/2}]=\operatorname{Tr}[(\bm\Sigma_{\mathrm r}^{1/2}\bm\Sigma_{\mathrm g}\bm\Sigma_{\mathrm r}^{1/2})^{1/2}]$理解，右侧是对称半正定矩阵的平方根，迹为非负实数。当两个协方差都是对角阵时，计算可以化到逐维：设第$k$维的标准差为$\sigma_{\mathrm r,k}$和$\sigma_{\mathrm g,k}$，则$(\bm\Sigma_{\mathrm r}\bm\Sigma_{\mathrm g})^{1/2}$的第$k$个对角元为$\sigma_{\mathrm r,k}\sigma_{\mathrm g,k}$，迹项变为$\sum_k(\sigma_{\mathrm r,k}^2+\sigma_{\mathrm g,k}^2-2\sigma_{\mathrm r,k}\sigma_{\mathrm g,k})=\sum_k(\sigma_{\mathrm r,k}-\sigma_{\mathrm g,k})^2$。

用一个二维特征的教学构造检查各项含义。设真实特征均值为$(0,0)$、两维标准差都是1；生成特征均值为$(0.3,0.4)$、标准差为$(0.5,1)$，协方差都取对角。均值项为$0.3^2+0.4^2=0.25$，迹项为$(1-0.5)^2+(1-1)^2=0.25$，FVD为0.5。第二部分的0.25来自生成视频在第一维上散布过窄，这对应“生成结果彼此太像”的多样性不足；即使把均值完全对准，这一项也不会消失。FVD越低越好，但它有几条使用边界。多元高斯只是特征分布的近似；协方差要从样本估计，样本少时估计误差可能大于模型之间的差别；所用网络版本、帧数、分辨率和样本数改变，数值就不可比较。它还是整批视频的单一数字，不能指出哪一段视频在哪一帧出了问题。

为了定位失败，\ref{sec:vision-multigen-vbench}节的VBench把视频质量拆成16个维度，每个维度用一种专门工具给出逐视频分数。以主体一致性为例，设$\bm d_t$是第$t$帧的DINO特征并已归一化为单位长度，$T$为帧数，分数为
\begin{equation}
  S_{\mathrm{subject}}
  =\frac{1}{T-1}\sum_{t=2}^{T}
    \frac12\left(\langle\bm d_1,\bm d_t\rangle+\langle\bm d_{t-1},\bm d_t\rangle\right)\text{。}
  \label{eq:vision-multigen-subject-consistency}
\end{equation}
与首帧比较约束长时间的外观保持，与前一帧比较约束相邻帧的连续变化，两项平均\scite{vision-vbench2024}。三帧教学构造中，若$\langle\bm d_1,\bm d_2\rangle=0.9$、$\langle\bm d_1,\bm d_3\rangle=0.7$、$\langle\bm d_2,\bm d_3\rangle=0.8$，第2帧的项为$(0.9+0.9)/2=0.9$，第3帧为$(0.7+0.8)/2=0.75$，分数为0.825。一段几乎不动的视频也能得到接近1的分数，所以这一维必须与动态程度一起看。

声音的分布级指标沿用同一思路。\term{Fréchet音频距离}（Fréchet audio distance, FAD）把音频切成1秒的窗，每0.5秒取一窗，送入VGGish音频分类网络，取分类层之前的128维激活作为嵌入，再对生成音频与参考音频的嵌入按式\eqref{eq:vision-multigen-fvd}计算Fréchet距离\scite{vision-fad2019}。VGGish的输入是由1秒音频形成的64维对数梅尔特征序列，相位信息不进入嵌入，因此只改变相位的失真可能检测不到。MMAudio等后续工作还用PaSST、PANNs等音频网络的嵌入计算同一距离，统称Fréchet距离（FD），不同嵌入的数值不能互相比较\scite{vision-mmaudio2025}。分布距离不看配对关系：把所有声轨随机交换给别的视频，FAD不变。检查“这段声音是否属于这段视频”需要逐片段的指标。一种做法是用预训练音频分类器把生成音频和该视频的原始声轨分别变成类别概率分布$\bm p_i^{\mathrm{gen}}$和$\bm p_i^{\mathrm{ref}}$，再对$n$个测试片段平均二者的KL散度，例如$\frac1n\sum_i\operatorname{KL}(\bm p_i^{\mathrm{ref}}\|\bm p_i^{\mathrm{gen}})$；散度的方向和概率的归一化方式随所沿用的实现而定，比较模型时须保持一致。另一种做法不需要原声轨：用\ref{sec:vision-models-imagebind}节的ImageBind分别编码视频与生成音频，取两个嵌入的余弦相似度，称为IB分数。两者都在语义层面判断“是不是这种声音”，不回答“是不是在这一刻响”。

音画同步需要时间上的指标。逐事件的检查用\ref{sec:vision-multigen-audio-check}节式\eqref{eq:vision-event-sync-error}的起声偏移：先把视频接触与音频起声配对，再计算偏移及其平均绝对值。逐事件标注成本高时，可以用自动估计的整段偏移：DeSync用在音画偏移检测上训练的Synchformer预测整段音频相对视频的时间偏移，取绝对值，单位为秒；八秒音频时，MMAudio论文对开头和结尾两段4.8秒的截取分别预测再平均。DeSync越小越同步，但它只给出整段的一个偏移，不能指出哪次碰撞出错。

三维内容的指标取决于有没有真实形状。没有参考形状时，只能检查渲染结果是否符合文字。DreamFusion使用CLIP R-Precision：从固定的仰角和若干方位角渲染生成对象，用CLIP判断渲染图在一组文字描述中最匹配哪一条，统计取到正确描述的比例\scite{vision-dreamfusion2023}。设共有$N$条描述$c_1,\dots,c_N$，第$i$个对象的渲染图为$\bm x_i$，图像与文字编码器为$f_{\mathrm{img}},f_{\mathrm{txt}}$，则
\begin{equation}
  \mathrm{R@1}
  =\frac1N\sum_{i=1}^N
   \mathbb 1\!\left[
    \argmax_{j}\cos\bigl(f_{\mathrm{img}}(\bm x_i),f_{\mathrm{txt}}(c_j)\bigr)=i
   \right]\text{。}
  \label{eq:vision-multigen-clip-rprecision}
\end{equation}
论文同时在带颜色的渲染图和去掉纹理的几何渲染图上计算这一比例，因为把纹理画在平板几何上，颜色渲染也可能得到高分。有参考形状时可以直接比较几何。Wonder3D在Google Scanned Objects中的30个对象上，用倒角距离与体积交并比比较重建表面和真实表面，用PSNR、SSIM与LPIPS比较生成的颜色视图与真实视图\scite{vision-wonder3d2024}后三者的定义见\ref{sec:vision-restoration-quality}节。设$P,Q$是从两个表面上采样的点集，\term{倒角距离}（Chamfer distance）常写作
\begin{equation}
  d_{\mathrm{CD}}(P,Q)
  =\frac{1}{|P|}\sum_{\bm p\in P}\min_{\bm q\in Q}\lVert\bm p-\bm q\rVert_2^2
   +\frac{1}{|Q|}\sum_{\bm q\in Q}\min_{\bm p\in P}\lVert\bm q-\bm p\rVert_2^2,
  \label{eq:vision-multigen-chamfer}
\end{equation}
两项分别检查一个表面上的点能否在另一个表面附近找到对应点，二者合起来同时惩罚多出的部件和缺失的部件。教学构造中取$P=\{(0,0),(1,0)\}$、$Q=\{(0,0),(1,0.2)\}$，两个方向各有一个点的最近平方距离为0.04，平均后各为0.02，倒角距离为0.04。不同实现可能用距离而不是平方距离，或先把形状缩放到单位尺度，数值只在同一实现内可比。\term{体积交并比}（volume intersection over union, volume IoU）把两个形状体素化后计算交集体素数与并集体素数之比，1表示完全重合。生成的多个视角之间是否互相一致，目前没有一个被普遍采用的单一指标；常用的检查是在已知相机下把一个视角中的点按深度投影到另一个视角，比较对应位置的颜色或法线，或者干脆把多视角图重建成三维再看重建误差。

表\ref{tab:vision-multigen-metrics}把这些指标按比较对象归在一起。表中每一行回答的问题不同，一个指标的高分不能代替另一行。

\begin{table*}[tbp]
 \centering
 \small
 \begin{tabularx}{\linewidth}{@{}p{20mm}p{30mm}p{34mm}X@{}}
  \toprule
  成果 & 指标 & 比较对象 & 不能说明的内容 \\
  \midrule
  视频 & FVD & 生成与真实视频的I3D特征分布 & 单个视频的哪一帧失败；是否符合文字 \\
  视频 & VBench各维度 & 每段视频的帧间特征、光流、检测结果等 & 事件是否合理；不同维度之间的取舍 \\
  音频 & FAD或FD & 生成与参考音频的嵌入分布 & 声音是否属于该视频；何时发声 \\
  音频 & KL、IB分数 & 与原声轨或视频的逐片段语义 & 起声时刻 \\
  音画 & 起声偏移、DeSync & 事件时间或整段偏移 & 声源是否正确；DeSync不能定位到事件 \\
  三维 & CLIP R-Precision & 渲染图与文字描述 & 几何是否正确；背面是否合理 \\
  三维 & 倒角距离、体积交并比 & 重建表面与真实表面 & 纹理质量；需要真实形状 \\
  \bottomrule
 \end{tabularx}
 \caption{本章使用的主要评价指标。分布级指标针对一批成果，逐片段指标针对单件成果；各指标依赖的特征网络或工具改变时，数值不可比较。}
 \label{tab:vision-multigen-metrics}
\end{table*}
\FloatBarrier

有了这些度量，开篇请求的难点可以拆成五项。每一项都先在玩具狗与花瓶的例子中出现，然后看经典方法能处理到哪一步，学习方法改变了什么。

\subsection{帧间依赖带来的时间一致性困难}
\label{sec:vision-multigen-challenge-temporal}

一段五秒、每秒16帧的视频有80幅画面。每一幅单独看都可能是合理的玩具狗图像，但相邻两幅还必须是同一个连续世界的两个瞬间：背上的花纹不能逐帧换位置，腿的摆动要有一致的节奏，玩具狗贴近花瓶时两者不能互相穿透。逐帧独立生成会同时破坏这三点，常见表现是画面闪烁、外观漂移和运动抖动。运动是否合理更难判断：步幅忽大忽小、碰到花瓶却没有任何反应，都可能出现在每帧清晰的视频里。

经典视频处理有两类工具与此相关。一类是运动估计：\ref{sec:vision-videofound-motion}节的块匹配与光流给出相邻帧之间每个位置移动到哪里。另一类是\ref{sec:vision-videofound-consistency}节的时间连续性度量：闪烁度量比较静止区域相邻帧的亮度变化；扭曲误差先用光流把前一帧变形到后一帧的位置，再在非遮挡区域比较二者的差异，差异大说明外观没有沿运动保持。在复原或增强已有视频时，这些工具可以直接构成约束，例如沿运动轨迹做时间滤波。生成任务却没有原视频，运动本身就是要生成的内容；光流只能在生成之后用来检查，不能告诉模型玩具狗应当怎样走。

学习方法把“整段视频由一个模型共同决定”作为出发点。\ref{sec:vision-multigen-create-video}节的Lumiere在网络内部同时下采样时间和空间，使远隔的帧在最粗层级共同参与计算；Wan先用视频自编码器压缩相邻帧的冗余，再让注意力覆盖全部潜位置。运动先验来自大规模视频数据，而不是手工的平滑假设。经典度量并没有被淘汰，而是以学习工具的形式回到评价中：VBench的运动平滑度借插帧模型检查帧间运动，动态程度借RAFT光流衡量运动幅度。视频修改面对同一困难的另一面：\ref{sec:vision-multigen-edit-video}节要求新内容逐帧一致，未修改的部分还要保持原样。运动在物理上是否合理，最终还要交给\ref{sec:vision-multigen-vbench2}节按事件提问的评价。

\subsection{模态时间尺度差异带来的同步困难}
\label{sec:vision-multigen-challenge-sync}

玩具狗在1.0秒碰到花瓶，碰撞声就应在1.0秒附近开始。困难在于两种媒体的时间单位相差很大。每秒24帧的视频，一帧覆盖约41.7毫秒；每秒8帧的语义特征，一个位置覆盖125毫秒；16 kHz的音频每秒有16000个采样点，压缩到潜空间后每秒仍有约31个位置。若生成网络把视频第$k$个位置与音频第$k$个位置当作同一时刻，就会产生成倍的时间错位。更细的困难是事件粒度：一次碰撞在画面里可能只占一两帧，平均到整段的特征里几乎看不出来。

经典的多模态对齐处理的是两条已经存在的序列。\ref{sec:vision-alignment-dtw}节的动态时间规整在两条序列之间寻找一条单调的对应路径，使累积的逐点代价最小，允许一方局部快一些、另一方局部慢一些；\ref{sec:vision-alignment-alignment}节还区分了对齐的粒度，即在帧、片段还是整段层面建立对应。在配声的生成任务中，这两点各自起作用。粒度的区分直接决定评价方式：DeSync给的是整段粒度的偏移，起声偏移是事件粒度的检查。动态时间规整则需要两条序列都已存在，而生成时至少有一条尚未产生，它不能把同步“写进”生成过程；它仍可在核验时用来匹配事件，但它的单调路径只能拉伸时间，不能表达多出或漏掉的一次碰撞，这些要另行计数。

学习方法把时间对应放进生成的条件接入中。MMAudio按比例缩放视觉序列的旋转位置编码，并加入每秒24个位置的同步特征逐位置调制音频表示；Movie Gen Audio让每个音频潜位置读取时间上最近的视频帧嵌入；Ovi和LTX-2在两个分支的每一层用带时间位置的交叉注意力双向交换信息，这是\ref{sec:vision-alignment-intermediate-fusion}节中间融合在生成模型中的形态。位置编码只让“同一时刻”可以被表达，正确的对应仍要从配对数据中学到。画面固定时的配声由\ref{sec:vision-multigen-audio}节回答，画面与声音都可变时由\ref{sec:vision-multigen-joint}节回答，同步的核验在\ref{sec:vision-multigen-audio-check}节。

\subsection{生成长度增长带来的状态保持困难}
\label{sec:vision-multigen-challenge-long}

本章的视频模型一次覆盖约五秒，音频模型一次最多约30秒。一个一分钟的玩具狗故事因此要由十几段视频接成。每段内部可以很连贯，段与段之间却可能出现三类问题：玩具狗被花瓶遮挡后重新出现时换了花纹，即身份不保持；花瓶的位置每段漂移几个像素，累积起来离开了原位，即场景不保持；背景环境声在段落交界处突然变化。时间越长，需要记住的状态越多，而模型一次能看到的范围是固定的。

经典视频分析为长视频提供了结构单位。\ref{sec:vision-videofound-structure}节的镜头边界检测把视频切成镜头，关键帧选择用少数画面代表一个镜头，背景建模把静止的场景与运动的前景分开；\ref{sec:vision-videofound-consistency}节的身份连续性度量检查同一对象在不同时间是否仍被识别为同一对象。这些方法描述已有视频的结构。在生成任务里，它们转而成为规划和检查的单位：先按镜头规划，再以关键帧作为每段的条件，最后用身份度量检查跨段一致。但经典方法本身不生成任何新的帧。

学习方法目前的主要手段是把已知内容作为条件接到下一段。\ref{sec:vision-multigen-extension}节用Wan的已知帧掩码续接视频；Movie Gen Audio把前一段末尾的音频作为下一段的上下文，或者在每个生成步对重叠窗口共同更新；Lumiere的空间超分辨率也用重叠窗口协调相邻片段。这些做法改善了交界处的衔接，但都没有保存一个持久的场景状态，误差会沿续接链条传递。这项挑战在本章只得到部分回答，本章小结会回到它。

\subsection{视角变化带来的几何一致性困难}
\label{sec:vision-multigen-challenge-geometry}

把玩具狗做成三维资产，要求从正面、侧面和背面看到的画面都来自同一个表面。二维图像生成模型分别画这三幅图，常常各自合理却互相矛盾：耳朵形状不同，或者背面也长出一张脸。只有一幅正面照片时，背面根本没有被观测；玩具狗在运动时，不同时刻拍到的腿部位置又不能拟合成同一个静止表面。

经典的多视图几何在有观测的条件下给出严格的一致性。\ref{sec:vision-formation-two-view}节的对极几何说明，同一个三维点在两幅图中的投影必须满足对极约束，已知两个相机时可以三角化出它的位置；\ref{sec:vision-formation-multi-view}节的运动恢复结构与光束法平差把这一点推广到多幅图，同时估计相机和三维点。\ref{sec:vision-generation-novel-view}节的NeRF与3D Gaussian Splatting则用可微渲染把一组带位姿的照片拟合成场景表示，表示本身是共享的，所以各视角自动一致。这些方法的前提是观测覆盖了要重建的表面，并且图像之间能找到对应；它们不会为没有拍到的背面创作内容，也不能从一句文字出发。

学习方法把二维生成先验和三维一致性结合起来，结合的位置有三种。DreamFusion保留一个共享的三维表示，用冻结的二维扩散模型为每幅渲染图提供修改方向；Wonder3D让二维扩散模型一次生成多个视角，在生成时就用跨视角注意力协调；4D Gaussian Splatting在有观测时不再需要生成先验，用形变场把多时刻观测连成一个动态场景。三者分别对应\ref{sec:vision-multigen-3d}节的三种观测条件。

\subsection{迭代采样带来的制作预算困难}
\label{sec:vision-multigen-challenge-budget}

本章的生成模型都通过多步迭代得到结果。Wan的报告通常采样约50步，若每步做无分类器引导，还要多算一次去条件预测，一段视频就需要约100次主干前向；LTX-2分别调节文字与模态引导时每步三次前向；DreamFusion为一条文字请求优化15000步，论文在4块TPUv4芯片上用约1.5小时。单次成本之外，还有候选数量：一条请求往往要生成多个候选，再挑选、局部重做。只比较最终最好的样本，会把这些成本都藏起来。

经典方法降低计算的主要思路是由粗到细与去除冗余。\ref{sec:vision-transforms-pyramid}节的图像金字塔先在低分辨率上求解全局结构，再逐级补充细节；视频编码利用相邻帧的相似性只存储变化部分。这两种思路都在学习方法中延续：Lumiere先在$128\times128$生成整段运动，再用空间超分辨率补细节；Wan的视频自编码器把时间压缩4倍、空间各压缩8倍，生成网络面对的位置数降到原来的约$1/256$，再在压缩后的潜网格上迭代；LTX-2还在基础潜变量生成后再做潜空间的空间上采样与分块细化。压缩减少了每步的计算，却引入了自编码器的重建误差；级联减少了基础模型的计算，却把协调问题留给了超分辨率阶段。

成本因此不是附带信息，而是比较模型时必须固定的条件。各节在介绍模型时给出论文报告的步数、耗时与规模，\ref{sec:vision-multigen-budget}节说明怎样在同一预算口径下比较制作流程。

这五项挑战决定了本章的展开顺序。时间一致性在只有画面的设定中最纯粹，所以\ref{sec:vision-multigen-video}节先处理视频的创建、修改和延伸，延伸部分同时触及长时状态保持。画面固定之后，同步困难只剩一个方向，由\ref{sec:vision-multigen-audio}节在配声中回答；画面和声音都可变时，同步要求两路在生成中互相读取，这是\ref{sec:vision-multigen-joint}节的问题。几何一致性把一致性的对象从时间轴换到视角，由\ref{sec:vision-multigen-3d}节按观测多少依次讨论。\ref{sec:vision-multigen-response}节的语音回答把同步要求推到播放时刻：说出的内容要等证据到达，又不能拖得太久。最后，\ref{sec:vision-multigen-evaluation}节用本节的指标和预算口径汇总验收，并说明哪些挑战仍未解决。

\section{视频生成与编辑}
\label{sec:vision-multigen-video}

视频生成与编辑的输出是一段帧序列，输入则随任务不同：创建视频时只有文字，或再加一幅首帧；修改视频时还有原片段，以及说明改哪里的指令或掩码；延伸视频时有一段已知帧，待生成的是它之前、之后或之间的时间范围。三类任务共同的难点是时间一致性：同一只玩具狗要在每一帧保持外观，运动要连贯，而一段五秒、每秒16帧的视频就有80幅画面，网络很难在显存允许的范围内同时看到全部像素。修改和延伸还多一个难点：模型必须知道哪些内容属于输入、应当保留，哪些内容允许改变。

本节按任务类型划分小节。创建视频时，两条路线对“怎样让整段视频共同参与计算”给出不同回答：Lumiere在像素空间的网络内部压缩时间，Wan先用视频自编码器把整段视频压成潜网格，再在潜网格上用Transformer处理全部位置。修改视频时，两条路线对“怎样告诉模型改哪里”给出不同回答：Movie Gen Edit从文字指令中学习修改关系，Lumiere视频修补用时空掩码直接指定区域。延伸与补全使用Wan的已知帧掩码接口，它把“哪些帧已知”编码为条件，与修补中“哪些像素已知”的做法相对应。

\subsection{视频的创建}
\label{sec:vision-multigen-create-video}

创建视频的输入是文字，有时还有一幅首帧；输出是规定时长、帧率和画幅的帧序列。文字规定场景和事件，首帧还规定故事开始时的外观与布局。若请求包含精确轨迹，则需要模型接收随时间变化的位置或运动条件；仅支持文字和首帧的模型没有获得这条数值轨迹。相同的“镜头绕行”，也可能由相机移动、对象旋转或两者共同运动形成，不能只检查画面里出现了转动。

开始生成前，制作记录要把时长换算成帧数。一个教学记录规定四秒、每秒24帧、固定画幅，以及“玩具狗从花瓶左侧走到右侧，花瓶保持不动”。按每帧占据一个播放时间间隔的约定，需要96帧，第$i$帧的起始时刻为$(i-1)/24$秒。若只规定首尾采样时刻相距四秒，则包含两端会得到97帧。模型自身还可能要求帧数符合特定分组规则，制作记录应把采样范围、容器播放时长和模型输入长度分别写清。

两个代表模型都要解决同一个矛盾：时间一致性要求网络同时看到整段视频，而计算预算不允许在全分辨率下这样做。Lumiere保留像素表示，在网络内部同时对空间和时间做下采样，在最粗的时空网格上处理整段运动，再用分窗口的空间超分辨率补细节。Wan把压缩前移到独立训练的视频自编码器中，生成网络只面对压缩后的潜网格，并用流匹配训练。前者的代价落在级联的超分辨率阶段，后者的代价落在自编码器的重建误差和潜空间序列长度上。

\subsubsection{Lumiere}
\label{sec:vision-multigen-lumiere}

Lumiere之前的文字生成视频系统常采用级联：基础模型先生成时间上稀疏的关键帧，再由时间超分辨率模型在关键帧之间插入画面。这种做法在关键帧间距较大时会遇到具体困难。设玩具狗的腿每0.25秒摆动一个周期，关键帧若每0.5秒取一幅，相邻两幅中腿的位置几乎相同；插帧模型只看到两幅相似画面，无法判断其间是静止、摆了一次还是摆了两次，生成的步态因此可能变慢或抖动。这是时间采样不足造成的混叠，补帧阶段缺少决定运动所需的信息\scite{vision-lumiere2024}

Lumiere的主要想法是让基础模型一次处理整段视频的全部帧，在网络内部对时间维也做下采样和上采样。时间被压短后，远隔帧的信息能在较小表示上共同参与计算，整段运动在最粗的时空网格上就已决定；随后逐层上采样，恢复完整帧率。由于基础模型直接输出完整帧率序列，后续阶段只需补空间细节，不再需要时间插帧，上述混叠来源就被移除了。作者把这种网络称为\term{时空U-Net}（space-time U-Net, STUNet），即在空间和时间两个方向上同时做多尺度下采样与上采样的U形去噪网络。

模型从预训练的Imagen文字生成图像模型出发。Imagen在像素空间工作，由一个低分辨率基础模型和空间超分辨率级联组成。Lumiere基础模型的输入是含噪视频$\bm X_\ell\in\mathbb R^{80\times128\times128\times3}$、噪声水平$\ell$和文字表示，输出同形状的噪声预测；80帧按每秒16帧对应五秒。论文给出的结构信息有三条。第一，时间下采样和上采样模块紧跟在预训练网络的每个空间缩放模块之后，所以空间每缩小一级，时间也同时缩短。第二，除最粗层级外，每个层级插入分解的时空卷积，即先做空间卷积、再沿时间做一维卷积；分解的好处是新增参数只有一维时间核，原有空间卷积可以原样保留。第三，最粗层级插入时间注意力，让每个空间位置在全部剩余时间位置之间交换信息；注意力的成本随序列长度平方增长，只有在时间已被压短的最粗层级才负担得起。新增模块以残差方式接入，其输出投影初始化为零映射。论文没有给出层级数、每一级的时间下采样倍数和参数量。为了看清张量形状怎样变化，以每个层级时空各减半作为教学假设，形状依次为$80\times128^2\to40\times64^2\to20\times32^2\to10\times16^2$，最粗层级只剩10个时间位置；表\ref{tab:vision-multigen-lumiere-stunet}按这个假设列出各层级的组成，其中形状一列不是论文发布的配置。

\begin{table*}[tbp]
 \centering
 \small
 \begin{tabularx}{\linewidth}{@{}p{24mm}p{30mm}p{34mm}p{26mm}X@{}}
  \toprule
  层级组 & 预训练空间部分 & 新增时间层 & 时空采样 & 输出形状（教学假设） \\
  \midrule
  编码第1级 & 空间卷积块，冻结 & 分解时空卷积 & 空间、时间各下采样 & $40\times64\times64$ \\
  编码第2级 & 空间卷积块，冻结 & 分解时空卷积 & 空间、时间各下采样 & $20\times32\times32$ \\
  编码第3级 & 空间卷积块，冻结 & 分解时空卷积 & 空间、时间各下采样 & $10\times16\times16$ \\
  最粗层级 & 空间卷积与注意力，冻结 & 时间注意力 & 无 & $10\times16\times16$ \\
  解码各级 & 空间卷积块，冻结，接收跳跃连接 & 分解时空卷积 & 空间、时间各上采样 & 逐级恢复到$80\times128\times128$ \\
  \bottomrule
 \end{tabularx}
 \caption{Lumiere时空U-Net的层级组成。时间层的种类与位置依据论文；论文未给出层级数、每级下采样倍数与参数量，表中层级数和输出形状是便于手算的教学假设，通道数随原图像U-Net而定，论文未单独列出。}
 \label{tab:vision-multigen-lumiere-stunet}
\end{table*}
\FloatBarrier

训练使用约3000万段带文字描述的视频，每段80帧、每秒16帧。原图像模型的权重保持固定，只更新新增的时间参数。论文没有单列改写后的目标函数；按其沿用的扩散概率模型训练，可用常见的噪声预测形式写成
\begin{equation}
  L_{\mathrm{Lum}}
  =\mathbb E_{\bm X,c,\ell,\bm\epsilon}
   \left[
    \left\lVert
     \widehat{\bm\epsilon}_{\bm\theta}(\alpha_\ell\bm X+\sigma_\ell\bm\epsilon,\ell,c)-\bm\epsilon
    \right\rVert_2^2
   \right],
  \label{eq:vision-lumiere-objective}
\end{equation}
其中$\bm X$是训练视频，$c$是文字条件，$\bm\epsilon$是与视频同形状的标准高斯噪声，$\alpha_\ell,\sigma_\ell$由噪声日程给出，$\bm\theta$只包含新增时间参数。冻结图像权重保留了大规模图文训练得到的外观先验，新增参数只需学习时间关系。时间下采样会混合或选择帧，网络初始化不能简单解释为“每帧恰好等于原图像模型”；原方法把时间采样模块初始化为最近邻下采样和上采样，再加上零初始化的输出投影，使训练开始时的网络行为接近逐帧运行图像模型。

推理时从整段噪声视频出发，反复运行时空网络并执行采样更新。一次前向能看完整段，不表示一次前向已经完成采样。基础模型输出$80\times128\times128$的视频后，空间超分辨率模型把每帧提升到$1024\times1024$。高分辨率网络受显存限制，只能处理较短时间窗口。若各窗口独立完成再拼接，边界两侧可能出现不同的毛色或光照。Lumiere采用MultiDiffusion的做法，在每一步采样中处理重叠窗口，再对重复覆盖的位置融合预测，使下一步各窗口从同一段全局状态继续采样；原文的时间窗口重叠两帧，窗口长度论文未给出。

设窗口$i$选择全局视频中的若干连续帧，选择算子为$\bm P_i$；把当前一步的窗口预测展平成向量$\bm y_i$。希望新的全局状态$\bm z$在每个窗口内都接近相应预测，可以求解
\begin{equation}
  \min_{\bm z}\sum_i\lVert\bm P_i\bm z-\bm y_i\rVert_2^2,
  \qquad
  \widehat{\bm z}
  =\left(\sum_i\bm P_i^\top\bm P_i\right)^{-1}
     \sum_i\bm P_i^\top\bm y_i\text{。}
  \label{eq:vision-window-consensus}
\end{equation}
这里假设所有全局位置至少被一个窗口覆盖，且$\bm P_i$只负责取出对应坐标。对$\bm z$求导并令梯度为零，就得到右侧的正规方程。矩阵$\sum_i\bm P_i^\top\bm P_i$是对角矩阵，每个对角元素等于该位置被覆盖的次数，因此解就是对重叠预测逐位置求平均。

例如八帧教学序列采用窗口1至5和4至8，第4、5帧被预测两次。若第4帧某个颜色分量的两次预测为0.4和0.6，合并值为0.5；只被覆盖一次的分量直接保留。图\ref{fig:vision-window-consensus}展示了覆盖计数的来源。这个平均发生在每步采样的状态协调中，不能用“最终把两段成片淡入淡出”代替。

把这一步写成算法时，还可以给每个窗口内的位置一个非负权重$w_{i,k}$，等权平均是全部权重取1的特例；\ref{sec:vision-multigen-movie-audio}节的Movie Gen Audio就使用中间高、两端低的三角形权重。若把式\eqref{eq:vision-window-consensus}中的平方误差按$w_{i,k}$加权，同样令梯度为零，每个全局位置的解就是覆盖它的各个预测按权重的加权平均，算法\ref{alg:vision-multigen-window-consensus}给出完整步骤。

\begin{algorithm}[重叠窗口的逐步共同更新]
  \label{alg:vision-multigen-window-consensus}
  \begin{algorithmic}[1]
    \Require 全局含噪状态$\bm z^{(0)}$；覆盖全部位置的窗口集合$\{\bm P_i\}$及窗口内权重$w_{i,k}\geq0$；单窗口去噪网络；采样步数$K$
    \Ensure 全局生成结果$\bm z^{(K)}$
    \For{$k=0,1,\dots,K-1$}
      \For{每个窗口$i$}
        \State 取出$\bm P_i\bm z^{(k)}$，运行网络与一步采样更新，得到窗口预测$\bm y_i$
      \EndFor
      \For{每个全局位置$u$}
        \State 累加分子$a_u\gets\sum_{(i,j):\,\bm P_i\text{的第}j\text{个坐标为}u}w_{i,j}\,y_{i,j}$
        \State 累加分母$b_u\gets\sum_{(i,j):\,\bm P_i\text{的第}j\text{个坐标为}u}w_{i,j}$
        \State 若$b_u=0$则报告窗口布置错误；否则$z^{(k+1)}_u\gets a_u/b_u$
      \EndFor
    \EndFor
  \end{algorithmic}
\end{algorithm}

用上面八帧的例子复核算法。若两个五帧窗口都采用三角形权重$(1,2,3,2,1)$，第4帧在窗口一中是第4个位置，权重为2，在窗口二中是第1个位置，权重为1。两个预测仍为0.4和0.6时，合并值为$(2\times0.4+1\times0.6)/(2+1)\approx0.467$，比等权平均更靠近权重较大的窗口一。权重的作用是让窗口边缘的预测少占比例，因为边缘位置只看到单侧的上下文。

\begin{figure}[htbp]
  \centering
  \begin{tikzpicture}[x=1cm,y=1cm,font=\figurefont\small,
  cell/.style={draw=BookInk,line width=0.7pt,minimum width=7mm,
    minimum height=7mm,inner sep=0pt},
  arr/.style={-{Stealth[length=2mm]},BookInk,line width=0.8pt}]
  \node[anchor=east] at (1.65,2.45) {帧位置};
  \node[anchor=east] at (1.65,1.6) {窗口一};
  \node[anchor=east] at (1.65,0.65) {窗口二};
  \node[anchor=east] at (1.65,-0.45) {覆盖次数};
  \foreach \i in {1,...,8} {
    \pgfmathsetmacro{\xx}{1.5+0.95*\i}
    \node at (\xx,2.45) {\i};
    \ifnum\i<6
      \node[cell,fill=FigureInput!15] at (\xx,1.6) {};
    \fi
    \ifnum\i>3
      \node[cell,fill=FigureOutput!15] at (\xx,0.65) {};
    \fi
    \ifnum\i=4
      \node at (\xx,-0.45) {2};
    \else
      \ifnum\i=5
        \node at (\xx,-0.45) {2};
      \else
        \node at (\xx,-0.45) {1};
      \fi
    \fi
  }
  \draw[BookAmber,dashed,line width=0.8pt] (4.85,0.15) rectangle (6.7,2.02);
  \draw[arr] (5.3,0.1) -- (5.3,-0.1);
  \node[anchor=west] at (1.65,-1.3) {同一位置：$(0.4+0.6)/2=0.5$};
\end{tikzpicture}

  \caption{重叠窗口的教学构造。两个五帧窗口在第4、5帧重叠，覆盖次数决定式\eqref{eq:vision-window-consensus}的逐位置平均。图中帧数用于手算，不是Lumiere的发布配置；原文空间超分辨率窗口在时间上重叠两帧。}
  \label{fig:vision-window-consensus}
\end{figure}
\FloatBarrier

让玩具狗绕花瓶走一圈时，可分别检查起终点是否连接、步态是否连贯、遮挡后外观是否保持。重叠融合能协调窗口边界，却不能证明模型理解了接触或三维形状。80帧是该设计所覆盖的时间范围，不能由此推出任意长度的叙事都能一次处理；短片段内部连贯与多镜头之间的身份保持，仍是不同的验收对象。

\Needspace{60mm}
\subsubsection{Wan}
\label{sec:vision-multigen-wan}

Wan的视频自编码器、扩散Transformer主干和条件接入方式已在\ref{sec:vision-models-wan}节说明，主干各规模的层数与宽度等配置也以该节为准，本节不再重复。这里只回顾使用时需要的结论：自编码器把视频在时间上压缩4倍、在高和宽上各压缩8倍，潜通道为16，第一帧只做空间压缩，作为时间起点单独保留；主干把潜网格切成时空块序列，用自注意力处理全部位置，文字经umT5编码后由交叉注意力接入；训练采用式\eqref{eq:vision-flow-loss}的流匹配速度目标\scite{vision-wan2025}。图像条件版本把首帧与其后补零的帧组成条件视频，经同一自编码器得到条件潜网格，再与含噪潜网格和重排后的已知位置掩码沿通道拼接；首帧还经CLIP图像编码器提取整体特征，通过独立的交叉注意力接入。

与Lumiere相比，Wan把时间一致性问题交给两处处理。自编码器先把冗余的相邻帧压缩，生成网络面对的位置数减少到原来的约$1/(4\times8\times8)$；主干的自注意力在压缩后的全部位置之间直接交换信息，不需要像STUNet那样在网络内部逐级缩短时间。因此Wan不需要级联的空间超分辨率，由自编码器解码直接得到目标分辨率的画面。报告给出1.3B和14B两种规模，训练使用每秒16帧、五秒长的片段，分辨率逐阶段提高到720p。

使用时首先组织输入。帧数要满足自编码器的分组约定，即第一帧之后的帧数是4的倍数。以$T=17$、$256\times256$画面为尺寸示例，目标潜网格为$5\times32\times32\times16$。文字生成只需要文字和目标尺寸；图像条件生成还需要首帧，首帧的宽高比应与目标画幅一致，否则缩放或裁剪会先改变参考内容。条件视频的其余帧填零，零表示未提供内容，不能把它当成“未来画面应为黑色”的监督目标。文字生成与图像条件生成使用不同的检查点，报告中图像编码器分支是在任务微调阶段才加入的，不能把文字生成权重直接配上首帧使用。

一般使用不再微调主干，主要调节采样设置。报告的推理通常使用约50步采样，并采用无分类器引导；引导强度越大，结果越贴近文字，但动作幅度和画面多样性可能下降。随机种子决定初始噪声，同一请求应保留多组种子的候选。采样结束后，潜网格由自编码器解码为帧序列；解码器的因果结构允许分块解码较长序列，但这不表示生成网络已经学会长时间状态保持。报告还指出，14B模型在单块高端GPU上未经优化的推理约需30分钟，这一成本应计入制作预算。

固定同一首帧，分别要求玩具狗转头、向前走和镜头环绕，可以观察首帧条件留下了什么自由度。转头要补出原来看不到的侧面，行走要决定腿部运动和接触，镜头环绕还要补背景与背面。首帧只给了起点，这些未知内容不会因“图生视频”这个任务名称而变成已知。常见失败包括：大幅运动时细节丢失，这也是报告列出的局限；转身后耳朵或花纹改变；文字与首帧冲突时，例如文字写“白色玩具狗”而首帧是棕色，模型可能在前几帧逐渐改色。冲突应在生成前消除，而不是期待模型自行选择。

两个模型可以在同一“玩具狗绕花瓶走一圈”的请求下比较，见表\ref{tab:vision-video-creation-compare}。Wan把时间压缩从生成网络移到自编码器，换来了单阶段的高分辨率输出和更多条件入口；代价是生成质量受自编码器重建误差限制，长序列的注意力成本随潜位置数的平方增长。Lumiere保留像素表示，冻结图像模型使训练只需学习时间参数，但五秒、$128\times128$的基础输出必须依靠分窗口超分辨率补细节，窗口之间只能靠逐步平均协调。

\begin{table*}[tbp]
 \centering
 \small
 \begin{tabularx}{\linewidth}{@{}p{24mm}XX@{}}
  \toprule
  比较项 & Lumiere & Wan \\
  \midrule
  生成空间 & 像素空间，基础模型$128\times128$，再级联空间超分辨率 & 视频自编码器的潜网格，时间压缩4倍、空间各8倍 \\
  整段时间的处理 & 网络内部时空下采样，最粗层级用时间注意力 & 自编码器先压缩，主干自注意力覆盖全部潜位置 \\
  训练目标 & 噪声预测扩散目标，冻结图像权重，只训练时间参数 & 流匹配速度回归，训练整个主干 \\
  条件入口 & 文字 & 文字；图像条件版本另有首帧潜网格、掩码和CLIP整体特征 \\
  输出范围 & 80帧、每秒16帧，超分辨率后$1024\times1024$ & 每秒16帧、五秒片段，最高720p \\
  主要代价 & 超分辨率窗口只能靠逐步平均协调 & 自编码器重建误差，注意力成本随潜位置数增长 \\
  \bottomrule
 \end{tabularx}
 \caption{视频创建两条路线的比较。各项依据两篇报告的说明，不是同一实验条件下的质量排名；两者都只覆盖约五秒的短片段。}
 \label{tab:vision-video-creation-compare}
\end{table*}
\FloatBarrier

\subsection{视频的修改}
\label{sec:vision-multigen-edit-video}

修改视频的输入是原片段和修改要求，输出是改动后的片段。难点有两个。第一，模型要区分哪些变化是任务要求，哪些变化应算误改。把花瓶换成金属杯主要改变对象外观；让玩具狗跳上桌子则改变动作、遮挡和接触，原背景中某些区域也会随之显露。保留范围必须根据请求判断，不能一律要求整段视频的所有像素不变。第二，真实的“原视频—编辑后视频”配对极少，模型很难直接从数据中学到修改关系。

两条路线对“怎样指定改哪里”给出不同回答。Movie Gen Edit只接收文字指令，由模型自己判断修改范围，监督则通过分阶段构造的训练配对逐步获得。Lumiere视频修补让用户提供时空掩码，直接规定允许改变的像素，训练配对可以从任意完整视频中遮去一部分自动产生。前者使用方便，但修改范围依赖模型理解；后者范围明确，但掩码必须随对象运动逐帧给出。

\Needspace{60mm}
\subsubsection{Movie Gen Edit}
\label{sec:vision-multigen-movie-edit}

只有“文字—视频”配对时，模型学到的是描述什么画面，尚未学到根据原视频执行修改。图像编辑已有较多“原图—指令—目标图”配对，但把它们逐帧用于视频，相邻帧的编辑会各自独立，杯子的形状可能逐帧跳变。Movie Gen Edit的出发点是：在没有成对视频编辑数据的条件下，借助视频生成模型和若干可构造的代理任务，逐步缩小单帧编辑与真实视频编辑之间的差距\scite{vision-moviegen2024}

主要想法有两点。结构上，把原视频作为额外输入通道接入已经训练好的视频生成主干，使编辑模型继承生成模型的运动先验。监督上，分三个阶段构造训练关系：先用图像编辑配对学单帧修改，再用带运动的合成序列学多帧修改，最后用回译把合成缺陷移到输入侧，让目标回到真实视频。

编辑模型从预训练的文字生成视频模型Movie Gen Video初始化。该模型在时间自编码器（temporal autoencoder, TAE）的潜空间中用流匹配生成视频。TAE在时间、高、宽三个方向各压缩8倍，潜通道为16，潜网格再用核为$1\times2\times2$的三维卷积切成词元。报告给出的最大配置是300亿参数的Transformer（不含文字编码器与TAE），共48层，宽度6144，48个注意力头，前馈层宽度16384，见表\ref{tab:vision-multigen-edit-config}；编辑一节只说明从“预训练的文字生成视频模型”初始化，没有单独写明编辑模型采用哪一个规模。以最大上下文为例，256帧、$768\times768$的视频经TAE压成$32\times96\times96$的潜网格，再按$2\times2$空间块切分，得到$32\times48\times48=73728$个视频词元，约为报告所说的73K。

主干的每一块沿用LLaMa3的Transformer块，使用RMSNorm和SwiGLU前馈层，并作三处修改：注意力改为双向，在自注意力和前馈层之间加入读取文字的交叉注意力，用自适应层归一化接入生成时刻。记第$l$块的输入词元为$\bm Z$、文字表示为$\bm C$、生成时刻嵌入为$\bm e_s$，一块的计算顺序为
\begin{equation}
 \begin{aligned}
  \bm Z'&=\bm Z+\operatorname{SA}\bigl(\operatorname{AdaNorm}(\bm Z;\bm e_s)\bigr),\\
  \bm Z''&=\bm Z'+\operatorname{CA}\bigl(\operatorname{RMSNorm}(\bm Z'),\bm C\bigr),\\
  \bm Z'''&=\bm Z''+\operatorname{SwiGLU}\bigl(\operatorname{AdaNorm}(\bm Z'';\bm e_s)\bigr),
 \end{aligned}
 \label{eq:vision-multigen-moviegen-block}
\end{equation}
其中$\operatorname{SA}$与$\operatorname{CA}$分别是多头自注意力和交叉注意力，$\operatorname{AdaNorm}$表示由$\bm e_s$预测缩放与平移的归一化。报告只说明加入自适应层归一化和交叉注意力的位置，交叉注意力前的归一化形式与调制系数的确切接法这里按常见做法写出，属于示意。位置编码采用可学习的分解形式，分别为时间、高、宽学习嵌入，并在每一层都加入，报告指出这样做减少了时间上的变形。文字由UL2、ByT5和长文本版MetaCLIP三个编码器处理，各自投影到6144维后拼接为$\bm C$，其中按字符编码的ByT5只用于画面中需要出现的文字。

\begin{table*}[tbp]
 \centering
 \small
 \begin{tabularx}{\linewidth}{@{}p{28mm}Xccccc@{}}
  \toprule
  模型 & 输入单元 & 层数 & 宽度 & 头数 & 前馈宽度 & 参数量 \\
  \midrule
  Movie Gen Video & TAE潜网格的$1\times2\times2$块，最多约73K个 & 48 & 6144 & 48 & 16384 & 300亿 \\
  \bottomrule
 \end{tabularx}
 \caption{Movie Gen Edit所依托主干的配置。数值为报告中Movie Gen Video最大规模的配置，参数量只计Transformer本身；编辑模型的规模报告未单独给出。本章各Transformer配置表采用相同的列。}
 \label{tab:vision-multigen-edit-config}
\end{table*}
\FloatBarrier

编辑模型在块嵌入层增加输入通道，把原视频潜网格与含噪目标潜网格沿通道拼接。以教学尺寸为例，64帧、$256\times256$的视频压成$8\times32\times32\times16$的潜网格，两者拼接后为$8\times32\times32\times32$，再切成$8\times16\times16=2048$个词元进入主干。每种编辑任务有一个学习到的任务嵌入：一个线性变换把它映射成四个向量，拼接到文字编码器的隐状态之后；另一个线性变换的输出加到生成时刻嵌入上。所有新增权重初始化为零，其余权重来自预训练的文字生成视频模型，因此训练开始时模型的行为与原生成模型相同。

训练目标仍是流匹配速度回归，只是速度网络多了原视频条件$\bm z^{\mathrm{src}}$和任务嵌入$k$：
\begin{equation}
  L_{\mathrm{edit}}
  =\mathbb E
   \left[
    \left\lVert
     \bm v_{\bm\theta}(\bm z(s),s,c,\bm z^{\mathrm{src}},k)
     -(\bm z_1-\bm z_0)
    \right\rVert_2^2
   \right],
  \label{eq:vision-edit-objective}
\end{equation}
其中$\bm z_1$是目标视频的潜网格，$c$是编辑指令，其余记号同式\eqref{eq:vision-flow-loss}。三个阶段的区别只在于$(\bm z^{\mathrm{src}},c,\bm z_1)$从哪里来。

第一阶段把图像编辑配对视为单帧视频编辑，与文字视频生成交替训练；纯视频生成样本以黑色视频作为空输入条件。论文在约六分之五的批次中采样图像编辑，训练3万步。混合这两类任务，使模型既能利用编辑关系，又不忘记视频运动。单帧训练的时间位置嵌入随机选择，避免所有图像编辑都只出现在序列起点。

单帧训练缺少“原片段连续输入”的经验。第二阶段因此使用两类多帧构造。一类对原图和编辑后图逐步施加相同的仿射变换，形成成对序列。由于两侧使用同样的变换，模型可以学习在相同运动下保留编辑差异，但这类平面移动不能提供真实关节运动。另一类从真实视频定位并分割对象，将相应区域涂成指定颜色，把分割任务转成“把这个对象标成某颜色”的编辑配对，补充真实时间变化。两类构造与文字生成视频按$3/5$、$1/5$、$1/5$的比例采样批次。论文报告这一阶段只训练约1000步，结果变得不那么模糊，但编辑后的元素运动偏少，有时颜色过饱和。

这些目标仍有一部分来自模型生成，可能携带失真。第三阶段使用\term{回译}（backtranslation）：先用前一阶段的模型把真实视频$V$按指令$r$编辑为$\widetilde V$，再由语言模型根据前后描述写出逆向指令$r^{-}$，用$(\widetilde V,r^{-},V)$训练。回译一词借自机器翻译，指把结果翻回原文来构造训练对。这样目标回到真实视频，合成缺陷主要留在输入侧，模型不必把自己的缺陷当作目标学习。例如真实片段里是红花瓶，前向请求为“改成金属杯”，合成输入是带杯子的片段，逆向请求为“把金属杯改回红花瓶”，训练目标仍是原始花瓶视频。这里的$r^{-}$是语义上构造的修改要求，不是保证可逆的数学算子；删除对象可能同时丢失背景细节，逆向指令也可能不准确。报告没有给出这一阶段的训练步数。

推理时，原视频经TAE编码，与指令及任务嵌入一起进入模型，目标潜网格从噪声逐步形成，再解码为编辑片段。检查金属杯的材质之外，还要检查杯体是否在遮挡后恢复为同一形状、玩具狗是否被连带改变。连续编辑时保存每轮输入与结果，才能发现哪一轮开始出现身份或背景漂移。Movie Gen的生成、编辑和配声分别使用相应路线，不能把模型族名称当成一个同时完成全部输出的网络。

\subsubsection{Lumiere视频修补}
\label{sec:vision-multigen-video-inpainting}

文字指令无法精确说明“只改第3至40帧中花瓶占据的像素”。Movie Gen Edit由模型自己推断范围，误改背景时用户没有直接的约束手段。\term{视频修补}（video inpainting）换一种条件：用户给出时空掩码，指出每帧哪些像素允许改变，模型根据其余可见内容和文字补出被遮区域。Lumiere的修补模型沿用前面的时空U-Net，整段时空处理让各帧补出的杯子共享上下文，避免逐帧修补时形状跳变。

结构上只改动输入层。含噪目标视频、保留可见区域的条件视频和二值掩码沿通道拼接：RGB目标3通道、RGB条件3通道、掩码1通道，输入层从3通道扩展到7通道，其余结构与基础模型相同。为便于讨论，本节令$m_{t,i,j}=1$表示允许生成的位置，0表示要保留的位置。训练时从完整视频$\bm X_t$构造条件
$\bm C_t=(\bm1-\bm M_t)\odot\bm X_t$，
掩码沿颜色通道广播；完整视频的含噪版本仍作为待预测对象，损失与式\eqref{eq:vision-lumiere-objective}相同，只是网络多接收了$\bm C_t$与$\bm M_t$。由于掩码可以从任意完整视频随机生成，这种监督不需要人工编辑配对。其他实现可能把掩码语义反过来，接入时须转换，不能把同一个0或1在不同接口中按相同含义解释。

推理时输入原视频、覆盖花瓶的逐帧掩码和替换要求，采样后检查边缘、遮挡及未指定区域。掩码若没有随花瓶移动，模型有时会在旧位置补出杯子，却在新位置留下原花瓶。区域跟踪与掩码传播可利用\ref{chap:vision-video-recognition}章的方法，但传播质量也属于编辑条件的一部分。

若交付明确要求掩码外像素完全相同，可以在制作阶段规定硬合成
\begin{equation}
  \bm X_t^{\mathrm{out}}
  =(\bm1-\bm M_t)\odot\bm X_t+
     \bm M_t\odot\widehat{\bm X}_t\text{。}
  \label{eq:vision-edit-composite}
\end{equation}
该式在掩码为0处代数上保留原像素；这是明确的输出规则，不是模型训练自动提供的保证。硬边界可能暴露接缝，改变光照或物体形状也可能需要更大的编辑范围。若使用软边缘，混合区域就不再逐像素相同，必须相应修改验收范围。

同一接口还能组织动态照片：把一张静态图复制到各个时间位置，只开放局部区域变化。例如让玩具狗摇尾巴而背景固定，掩码应包含尾巴摆动需要扫过的空间。复制帧提供静态背景，没有给出尾巴的唯一真实运动。

两种修改路线的差别见表\ref{tab:vision-video-edit-compare}。Lumiere视频修补用显式掩码换来了明确的修改范围和可自动构造的监督，代价是用户要准备逐帧掩码，而且掩码外的内容只由模型“倾向”保留，严格保持仍需式\eqref{eq:vision-edit-composite}。Movie Gen Edit去掉了掩码，能处理风格替换这类全局修改，代价是监督要分三阶段构造，修改范围依赖模型对指令的理解。

\begin{table*}[tbp]
 \centering
 \small
 \begin{tabularx}{\linewidth}{@{}p{24mm}XX@{}}
  \toprule
  比较项 & Movie Gen Edit & Lumiere视频修补 \\
  \midrule
  修改范围的指定 & 文字指令，由模型推断 & 用户提供逐帧二值掩码 \\
  条件接入 & 原视频潜网格沿通道拼接，任务嵌入接入文字与时刻嵌入 & 条件视频与掩码沿通道拼接，输入层扩展到7通道 \\
  监督来源 & 图像编辑、合成多帧编辑、回译三阶段构造 & 从完整视频随机遮挡自动构造 \\
  适合的修改 & 对象替换、风格变化、全局修改 & 局部区域替换、动态照片 \\
  典型失败 & 误改指令未提及的背景，多轮编辑累积漂移 & 掩码未跟随对象，边缘接缝 \\
  \bottomrule
 \end{tabularx}
 \caption{视频修改两条路线的比较。两者都以原片段为条件；差别在于修改范围由谁指定，以及训练配对从哪里来。}
 \label{tab:vision-video-edit-compare}
\end{table*}
\FloatBarrier

\subsection{视频的延伸与补全}
\label{sec:vision-multigen-extension}
\label{sec:vision-multigen-wan-extension}

延伸与补全的输入是一段或几段已知帧，输出是尚未给定的时间范围。给定历史片段向后续接，只有已发生的状态；给定首尾两侧，中间过程还受到终点限制；给定零散几帧，途中必须经过若干指定状态。难点在于新内容既要与已知帧在边界处衔接，又要在条件不足的方向上给出合理的新动作。创作允许新增动作与事件，而\ref{chap:vision-video-restoration}章的插帧侧重依据已采样观测估计中间时刻；两者可能使用相似掩码，评价证据却不相同。

这三种任务不需要三套模型。Wan的图像条件接口本来就是“已知帧放入条件视频、掩码标出已知位置”的形式，首帧生成只是其中已知帧最少的一种情形。把已知位置换成开头一段，就得到续接；换成首尾两帧，就得到首尾转换；换成零散几帧，就得到其他时间补全。接口的结构与训练见\ref{sec:vision-models-wan}节；报告在预训练中把这些掩码任务混合训练，再对具体任务分别微调。这里只讨论使用。

图\ref{fig:vision-time-masks}使用八个抽样位置展示条件差别。续接给出前三个状态，后面可以朝多种方向发展；首尾条件只固定两端，中间有多种合法运动；零散条件还限制途中必须经过的状态。图中的格子是教学时间位置，不是把八帧直接当作Wan自编码器的合法输入长度。

\begin{figure}[htbp]
  \centering
  \begin{tikzpicture}[x=1cm,y=1cm,font=\figurefont\small,
  cell/.style={draw=BookInk,line width=0.8pt,minimum width=7mm,
    minimum height=7mm,inner sep=0pt},
  known/.style={cell,fill=BookBlue,text=white}]
  \node[anchor=east] at (1.65,2.9) {时间位置};
  \node[anchor=east] at (1.65,2.1) {续接};
  \node[anchor=east] at (1.65,1.1) {首尾转换};
  \node[anchor=east] at (1.65,0.1) {零散条件};
  \foreach \i in {1,...,8} {
    \pgfmathsetmacro{\xx}{1.5+0.95*\i}
    \node at (\xx,2.9) {\i};
    \ifnum\i<4
      \node[known] at (\xx,2.1) {1};
    \else
      \node[cell] at (\xx,2.1) {0};
    \fi
    \ifnum\i=1
      \node[known] at (\xx,1.1) {1};
    \else
      \ifnum\i=8
        \node[known] at (\xx,1.1) {1};
      \else
        \node[cell] at (\xx,1.1) {0};
      \fi
    \fi
    \ifnum\i=2
      \node[known] at (\xx,0.1) {1};
    \else
      \ifnum\i=5
        \node[known] at (\xx,0.1) {1};
      \else
        \ifnum\i=8
          \node[known] at (\xx,0.1) {1};
        \else
          \node[cell] at (\xx,0.1) {0};
        \fi
      \fi
    \fi
  }
  \draw[-{Stealth[length=2mm]},BookInk,line width=0.8pt]
    (2.1,-0.7) -- (9.5,-0.7) node[below left] {时间递增};
\end{tikzpicture}

  \caption{时间条件的三种组织。实心格和数字1表示已提供内容，空心格和数字0表示待生成内容。位置按时间排序，编码与潜空间重排发生在此后；同一掩码形式不保证任意发布权重都支持相应任务。}
  \label{fig:vision-time-masks}
\end{figure}
\FloatBarrier

组织输入时，已知帧按其在目标序列中的时间位置放入条件视频，其余帧填零，再经自编码器编码。掩码保留帧级信息，再按自编码器的时间分组重排到通道维，所以已知帧最好与分组边界对齐：若一个四帧组中只有一帧已知，压缩后的条件潜变量会把已知帧与零帧混在同一组里，模型拿到的是部分信息。续接时通常把上一段末尾若干帧作为新一段的开头。报告构造续接训练数据时，用前1.5秒与后3.5秒画面的SigLIP相似度筛选内容连续的视频；构造首尾转换数据时，则提高首尾差别明显的样本比例。这说明两种任务需要的训练分布不同，使用时应选对应任务微调后的权重。

让玩具狗离开画面后继续前行，模型可以决定离开后的路线；给定它最终位于桌右侧的终帧，就要求生成过程连接两个状态。判断连接质量时，除了比较边界帧外观，还可以记录图像位置的有限差分。若边界前两次横向位移为2、2像素，新片段第一次位移为9像素，即使边界图像接得上，也出现了明显的速度变化。透视和相机运动会影响像素速度，所以此数值只用于固定相机的教学诊断。只给一帧作为续接条件时，模型无法从单帧推断速度，这类速度跳变更常见；给出多帧历史可以提供速度信息。

反复续接时，新片段成为下一轮条件，误差可能进入后续每一段。花瓶位置每次只漂移少量像素，长序列里也可能明显偏离原场景。首尾相接、局部运动连贯与全程状态一致应分别检查。把两端之间的创作画面用于说明一个故事是合理的；把它当成实际未拍到过程的观测证据，则超出了条件所能支持的结论。

\section{视觉条件下的声音生成}
\label{sec:vision-multigen-audio}

本节固定已有视频，待生成的对象是声轨。输入是视频，可选文字说明，有时还有一段已知声音；输出是与视频等长的音频波形。难点有三个：声音内容要符合画面中的声源，即“是什么声音”；起声要落在画面事件上，即“何时发声”；视频可能长达数分钟，而生成网络一次只能处理有限长度。模型可以在训练时同时使用视频、文字与音频，但只要推理输出只改变声音，就仍是视觉条件下的配声任务，与同时创造画面和声音使用不同的可修改范围。

本节按制作步骤展开：先把画面事件和创作选择整理成声音要求，再生成声轨，最后核验并调整。生成步骤中，MMAudio和Movie Gen Audio对上述后两个难点给出不同回答。MMAudio加入专门的高帧率同步特征来解决“何时发声”，处理的是八秒左右的短片段；Movie Gen Audio只用逐帧语义特征，但接收音频上下文，从而能延伸和修补长声轨。

\subsection{声音要求的确定}
\label{sec:vision-multigen-audio-conditions}

画面中一次接触可以对应轻敲、摩擦或破碎，具体声音还取决于材料、力度和录音位置。玩具狗移动也可以配脚步或机械声。视频往往没有唯一决定这些声学属性，因此需要把可见事件与创作选择分开记录。

可以用事件记录$(\tau_i^{\mathrm{start}},\tau_i^{\mathrm{end}},e_i,r_i)$描述第$i$个声音要求，其中时间以秒计，$e_i$表示对应事件，$r_i$说明声源、内容及保持要求。某次敲击记录来自视频中的接触；画外鸟鸣记录来自创作说明。两者都能成为有效条件，但后者不是模型从像素直接测得的事实。

同一四秒教学片段中，假定两次接触发生于1.0秒和2.5秒，可规定每次接触配一个短促轻响，其他时间保留低音量环境声。这份记录不等于模型一定支持逐事件数值控制；如果入口只有文字和视频，它只是验收依据。只有模型实际接收时间标记、局部掩码或已知音频上下文时，这些字段才直接进入生成计算。

\subsection{声轨的生成}
\label{sec:vision-multigen-audio-alignment}

语义特征帮助决定“什么声源”，时间特征帮助决定“何时发声”。一段视频的平均特征可以保留花瓶和桌面，却难以说明两次敲击之间的停顿；每秒只取几帧的语义特征，也可能恰好错过只持续几十毫秒的接触瞬间。生成器必须同时利用场景内容与局部变化，才能把声轨安排到画面事件上。

两个模型都在音频潜空间中用流匹配生成，也都把视觉特征按时间对齐到音频潜序列。它们的差别在于时间信息的来源和长度的处理。MMAudio在每秒8帧的语义特征之外，另加每秒24帧、专门为音画同步训练的特征；Movie Gen Audio的视觉特征只有逐帧语义嵌入，但它把已知音频作为上下文条件，用同一个模型完成整段生成、延伸和修补。

\subsubsection{MMAudio}
\label{sec:vision-multigen-mmaudio}

视频配声的配对数据有限，常用的VGGSound约有18万段十秒视频，MMAudio训练时取每段前8秒；只用这些数据训练，模型对声音种类和文字描述的覆盖不足。一种做法是先训练文字生成音频模型，再加入视觉条件微调，但这样视觉和文字在不同阶段学习，二者之间的交互较弱。MMAudio的主要想法是从头联合训练：同一个网络同时接收音视频配对和音频文字配对，缺失的模态用学习到的空条件替代，这样大规模音频文字数据也能帮助视频配声\scite{vision-mmaudio2025}。为了解决时间对应，它再加入帧级同步特征，以逐位置调制的方式影响音频生成。

音频先转为梅尔频谱，再由预训练的自编码器压成潜序列。梅尔频谱的计算已在\ref{sec:vision-models-encoding}节说明：把波形切成重叠短窗，对每窗求频谱，再按接近听觉尺度的梅尔刻度分组汇总能量，得到“时间窗$\times$频率组”的二维数组；这里只需记住它的一维是时间，所以压缩后的潜序列仍按时间排列。默认16 kHz版本的潜序列为每秒31.25个位置、每个位置20维；44.1 kHz版本为每秒约43.07个位置、40维。视觉帧和文字经CLIP提取特征，视觉特征为每秒8帧、1024维，文字为77个词元、1024维。生成主干先有$N_1$个多模态Transformer块，再接$N_2$个只更新音频表示的单模态块，最后预测音频潜变量的速度。

多模态块大体沿用SD3的MM-DiT设计：文字、视觉、音频三条序列各自保留一条计算支路，只在注意力处汇合。设三条序列在某一块中的查询、键、值分别为$(\bm Q_{\mathrm t},\bm K_{\mathrm t},\bm V_{\mathrm t})$、$(\bm Q_{\mathrm v},\bm K_{\mathrm v},\bm V_{\mathrm v})$、$(\bm Q_{\mathrm a},\bm K_{\mathrm a},\bm V_{\mathrm a})$，联合注意力把它们沿序列方向拼接后做一次注意力，
\begin{equation}
 [\bm O_{\mathrm t};\bm O_{\mathrm v};\bm O_{\mathrm a}]
 =\operatorname{softmax}\!\left(
   \frac{[\bm Q_{\mathrm t};\bm Q_{\mathrm v};\bm Q_{\mathrm a}]
         [\bm K_{\mathrm t};\bm K_{\mathrm v};\bm K_{\mathrm a}]^\top}{\sqrt{d_{\mathrm h}}}
  \right)[\bm V_{\mathrm t};\bm V_{\mathrm v};\bm V_{\mathrm a}],
 \label{eq:vision-multigen-joint-attention}
\end{equation}
再把输出按原长度切回三条序列，其中$[\cdot;\cdot]$表示沿序列方向拼接，$d_{\mathrm h}$为每个注意力头的维数。这样每个音频位置可以同时读取文字、视觉和其他音频位置。块内其余部分是：由全局条件经自适应层归一化提供的逐通道缩放与偏移；视觉和音频支路用核大小为3的一维卷积代替前馈层中的线性层（论文称为ConvMLP），使每个位置还能看到相邻时间位置；文字支路仍用普通前馈层。单模态块去掉文字与视觉支路，联合注意力退化为音频自注意力。论文给出的四种规模见表\ref{tab:vision-multigen-audio-config}，默认小模型取$N_1=4$、$N_2=8$、隐藏宽度448，约1.57亿参数；最大的L版本把块数加到7和14、宽度加到896，约10.3亿参数。论文没有给出注意力头数和前馈层宽度。

三种序列不具有相同的时间分辨率。音频潜序列为每秒31.25个位置，视觉语义特征为每秒8个位置，把二者相同的序号直接当作同一时刻会发生错位。模型在联合注意力之前对视觉和音频支路的查询与键施加旋转位置编码，并把视觉支路的频率乘以$31.25/8$：第$k$个视觉位置的有效位置变为$31.25k/8$，例如第8个视觉位置（第1秒）对应有效位置31.25，正好与第1秒的音频位置相同。文字没有时间含义，不加旋转位置编码。

更细的同步信息来自在音画偏移检测上预训练的Synchformer视觉编码器。它读取短时视频，提供与事件发生时间相关的特征，有效特征率为每秒24个位置、768维；附录的提取过程从每秒25帧输入形成重叠短片段，不能把24直接当作所有输入长度下不变的原始取样率。同步特征经过局部卷积处理，并用最近邻插值对齐到音频潜位置，再加上由文字、视觉及生成时刻得到的全局条件。令对齐后的第$j$个条件为$\bm c_j$，该位置音频隐状态为$\bm h_j\in\mathbb R^d$，调制写成
\begin{equation}
  \widetilde{\bm h}_j
  =\bm\gamma(\bm c_j)\odot
     \operatorname{LayerNorm}(\bm h_j)
    +\bm\beta(\bm c_j)\text{。}
  \label{eq:vision-audio-sync-modulation}
\end{equation}
两个学习映射给出逐通道缩放与偏移。全局条件会在所有$j$处重复，同步条件则随$j$变化，因此第1秒附近可以得到与第2.5秒不同的调制。这里没有把同步编码器的偏移分类结果直接当作波形，它提供的是生成条件。

在默认设置中，八秒音频包含128000个波形采样点，压缩后有250个潜位置；同一段视频的语义特征有64个位置，同步特征有192个位置。对齐不是要求这三个长度相等，而是让它们按共同时间关系进入交互。插值不会新增视频中未观察到的敲击事件，特征本身漏掉动作时，调制也不能凭空恢复证据。

训练数据包括VGGSound的音视频配对，以及AudioCaps、Clotho和WavCaps等音频文字数据，后者合计约95万对。训练从随机初始化开始，把真实音频潜变量与噪声连接，回归条件速度，复用式\eqref{eq:vision-flow-loss}的目标。视觉条件和文字条件各以10\%的概率替换为空条件；这样同一网络既学会了带条件的速度，也学会了不带某个条件的速度，推理时可以做无分类器引导。音频文字数据本来没有视频，视觉与同步入口始终使用空条件。

推理固定视频和文字，从音频噪声出发用欧拉法积分25步，无分类器引导强度取4.5，得到干净潜表示，再由自编码器恢复频谱、由\term{声码器}（vocoder）恢复波形。声码器是把频谱转换成可播放波形采样的网络，它不负责判断画面中的碰撞是否发生。论文报告小模型生成八秒声音约需1.23秒。固定花瓶落桌片段，分别给出“轻轻敲击”和“陶瓷破碎”的文字，可以区分内容选择与时间对应。若画面没有破碎而声轨出现碎裂声，作品可能满足了一个夸张音效要求，却不符合写实碰撞要求。MMAudio的音效和环境声结果也不能证明长段对白可懂、说话人稳定或歌词准确，这些输出需要自己的训练范围与评价。

\subsubsection{Movie Gen Audio}
\label{sec:vision-multigen-movie-audio}

MMAudio一次处理八秒左右，长片段若逐段独立生成，相邻段的背景声、响度和混响可能突变；它也没有接收已知声音的入口，不能只替换声轨中的一小段。Movie Gen Audio的主要想法是把已知音频作为上下文条件：训练时随机遮去目标音频的一部分，让模型根据视频、文字及其余已知声音预测被遮部分\scite{vision-moviegen2024}。没有上下文对应整段生成，已知一侧对应延伸，两侧已知对应局部修补，长声轨则通过反复延伸或多窗口协调形成。

模型是约130亿参数的扩散Transformer，用流匹配训练，输出48 kHz音频。波形先由单独训练的DAC-VAE压成每秒25帧、每帧128维的连续潜序列。上下文音频同样编码，保留内容的位置作为条件，其余位置填零，再与含噪目标潜变量逐位置沿通道拼接。视频每帧经长提示版MetaCLIP得到1024维嵌入，每个音频潜位置取时间上最近的视频帧，投影后逐位置加到音频表示上，而不是沿时间拼接成更长的序列。文字经T5-base编码为最多512个词元的768维特征，在每个Transformer块中通过交叉注意力接入。以十秒片段为例，目标潜序列为$250\times128$，与上下文拼接后每个位置有256维输入，另有250个对齐的视觉嵌入。

主干共36层，宽度4608，前馈层宽度18432，即四倍于宽度；130亿参数不含MetaCLIP、T5和DAC-VAE，注意力头数论文未给出。每一块依次是自注意力、读取文字的交叉注意力和前馈层，生成时刻的接入方式与\ref{sec:vision-models-wan}节回顾的Wan相同：一个所有层共享的小网络读取生成时刻嵌入，预测四个缩放和两个偏置共六个调制参数，每层再加上自己的一组偏置。把式\eqref{eq:vision-multigen-moviegen-block}中的自适应归一化理解为这六个参数对自注意力与前馈层输入输出的缩放平移，就得到这里一块的计算顺序。视频嵌入经带门控的线性投影后，与音频表示逐位置相加。表\ref{tab:vision-multigen-audio-config}把它与MMAudio的四种规模放在同一组列下。

\begin{table*}[tbp]
 \centering
 \small
 \begin{tabularx}{\linewidth}{@{}p{28mm}Xccccc@{}}
  \toprule
  模型 & 输入单元 & 层数 & 宽度 & 头数 & 前馈宽度 & 参数量 \\
  \midrule
  MMAudio S-16kHz & 每秒31.25个音频潜位置，20维 & 4+8 & 448 & 未给出 & 未给出 & 1.57亿 \\
  MMAudio S-44.1kHz & 每秒约43.07个音频潜位置，40维 & 4+8 & 448 & 未给出 & 未给出 & 1.57亿 \\
  MMAudio M-44.1kHz & 同上 & 4+8 & 896 & 未给出 & 未给出 & 6.21亿 \\
  MMAudio L-44.1kHz & 同上 & 7+14 & 896 & 未给出 & 未给出 & 10.3亿 \\
  Movie Gen Audio & 每秒25个DAC-VAE潜位置，128维，与上下文拼接为256维 & 36 & 4608 & 未给出 & 18432 & 130亿 \\
  \bottomrule
 \end{tabularx}
 \caption{两种视觉条件声音生成模型的配置。MMAudio的层数写成“多模态块数+单模态块数”；“未给出”表示论文没有报告该项。Movie Gen Audio的参数量不含文字、视频编码器与DAC-VAE。}
 \label{tab:vision-multigen-audio-config}
\end{table*}
\FloatBarrier

训练先在约百万小时量级的音频上预训练，单次长度不超过30秒，再用带视频的高质量数据微调。目标仍是式\eqref{eq:vision-flow-loss}形式的速度回归，只是速度网络多了视频、文字和音频上下文条件。论文明确说明模型不覆盖语音和带人声的音乐。推理用中点法求解64步，无分类器引导权重为3，并为每条请求生成8个候选再重排挑选。一次最多生成约30秒，即750个潜位置，更长的声轨用两种延伸方式。一种逐段自回归：把上一段末尾的音频作为下一段的已知上下文。若两个片段都包含同一次碰撞，就不能在最终结果里各保留一次，必须按共同时间位置合并重叠区域。另一种是多窗口共同更新：每次生成步把全局音频潜序列分成重叠窗口，各自预测，再按三角形窗口权重归一化合并，之后进入下一步。论文报告后者略好，作为默认方式。

多窗口共同更新就是算法\ref{alg:vision-multigen-window-consensus}，只是把\ref{sec:vision-multigen-lumiere}节的等权平均换成三角形加权平均。用中间高、两端低的三角形权重，窗口边缘的预测在合并中占比较小，可以减少从单窗口突然切到等权混合造成的突变；若某个位置的权重总和为零，则无法归一化，窗口布置必须保证覆盖。论文用窗口长度$n_{\mathrm{win}}$、相邻窗口终点之差$n_{\mathrm{hop}}$和重叠长度$n_{\mathrm{ctx}}=n_{\mathrm{win}}-n_{\mathrm{hop}}$描述窗口布置，以帧为单位，默认取$40$、$30$和$10$；在音效与音乐延伸实验中，它又以秒为单位取$n_{\mathrm{hop}}=n_{\mathrm{ctx}}=5.5$秒。

教学上，先为固定视频形成环境声，再只替换第二次敲击。条件声轨在第二次敲击附近留空，其余部分作为上下文；生成后检查新敲击是否出现、周围环境是否连续。即使潜空间有上下文条件，解码器仍可能影响邻近波形，所以要求严格保持的区间要在最终波形上核验。

两个模型的比较见表\ref{tab:vision-audio-gen-compare}。Movie Gen Audio相对MMAudio改变了两处设计：去掉专门的同步特征，改用每秒25帧的逐帧语义嵌入；增加音频上下文入口。前者使时间对应依赖模型从语义帧的变化中学习，后者换来了延伸、修补和长声轨能力。代价是模型规模大一到两个数量级，而且对快速连续的短促事件，没有专门同步特征时的对应精度需要单独核验。

\begin{table*}[tbp]
 \centering
 \small
 \begin{tabularx}{\linewidth}{@{}p{24mm}XX@{}}
  \toprule
  比较项 & MMAudio & Movie Gen Audio \\
  \midrule
  规模与输出 & 1.57亿至10.3亿参数，16 kHz或44.1 kHz & 约130亿参数，48 kHz \\
  语义条件 & CLIP视觉特征每秒8帧，联合注意力接入 & MetaCLIP逐帧嵌入，按最近帧对齐后相加 \\
  时间对应 & Synchformer同步特征每秒24帧，逐位置调制 & 无专门同步特征，依赖逐帧嵌入 \\
  长度与上下文 & 约八秒片段，无音频上下文入口 & 单次约30秒；音频上下文支持延伸与修补 \\
  训练 & 音视频与音频文字数据从头联合训练 & 大规模音频预训练，再带视频微调 \\
  不覆盖 & 长段对白、歌词等未作为目标 & 语音与带人声的音乐 \\
  \bottomrule
 \end{tabularx}
 \caption{两种视觉条件声音生成模型的比较。数值来自各自论文的说明，不是同一测试集上的质量比较。}
 \label{tab:vision-audio-gen-compare}
\end{table*}
\FloatBarrier

\subsection{声轨的核验}
\label{sec:vision-multigen-audio-check}

声音可以“听起来像敲击”，却早于画面半秒出现。对短促事件，可以先匹配声源与次数，再比较已匹配事件的起始时刻。设视频接触时刻为$\tau_i^{\mathrm v}$，音频起声时刻为$\tau_i^{\mathrm a}$，定义
\begin{equation}
  \delta_i=\tau_i^{\mathrm a}-\tau_i^{\mathrm v},
  \qquad
  e_{\mathrm{sync}}=\frac1m\sum_{i=1}^{m}|\delta_i|\text{。}
  \label{eq:vision-event-sync-error}
\end{equation}
$\delta_i>0$表示声音较晚。这里的$m$是已正确配对事件数，漏声、多声与声源错误另行计数，不能只对少量成功配对求平均而隐藏其余失败。

计算偏移之前要先回答“哪次起声对应哪次接触”。这正是\ref{sec:vision-alignment-alignment}节讨论的跨模态对齐问题，只是对齐的粒度是事件，而不是逐帧或整段。事件多、间隔不均匀时，可以借用\ref{sec:vision-alignment-dtw}节的动态时间规整：把接触时刻序列与起声时刻序列看成两条序列，逐对代价取时间差的绝对值，在保持先后顺序的前提下求累积代价最小的对应路径。路径的单调性保证了第一次接触不会与第二次起声配对、第二次接触再与第一次起声配对。但动态时间规整要求每个元素都有对应，它会把多出的一声强行并到某次接触上；所以在配对之后，应当把时间差超过事先规定的匹配窗口的对应拆开，计为多声或漏声；匹配窗口回答“算不算同一事件”，比下文判断同步是否合格的容差宽。教学上，若接触为1.0秒和2.5秒，起声为1.12秒、1.30秒和2.48秒，规整路径会把前两次起声都对应到第一次接触。取0.25秒的匹配窗口，1.30秒那一声与接触相差0.30秒，应计为多出的声音，而不是第一次接触的另一个偏移；1.12秒那一声保留为配对，再进入下面的同步判断。

沿用1.0秒和2.5秒的接触，若起声时刻为1.12秒和2.48秒，偏移为120毫秒和$-20$毫秒，平均绝对偏移70毫秒。若教学验收预先规定每次不得超过80毫秒，则只有第二次通过。这是便于计算的任务阈值，不是普适的人类听觉阈值。图\ref{fig:vision-event-sync}把两种方向的偏移放到共同时间轴上。

\begin{figure}[htbp]
  \centering
  \begin{tikzpicture}[x=1cm,y=1cm,font=\figurefont\small,
  axis/.style={BookInk,line width=0.8pt},
  contact/.style={BookBlue,line width=1.2pt},
  sound/.style={BookAmber,line width=1.2pt}]
  \node[anchor=east] at (1.6,3.4) {接触};
  \node[anchor=east] at (1.6,2.6) {起声};
  \draw[axis] (2,3.4) -- (10,3.4);
  \draw[axis] (2,2.6) -- (10,2.6);
  \foreach \t in {0,...,4} {
    \pgfmathsetmacro{\xx}{2+2*\t}
    \draw[axis] (\xx,3.35) -- (\xx,3.45);
    \draw[axis] (\xx,2.55) -- (\xx,2.65);
    \node[below] at (\xx,2.45) {\t};
  }
  \node[anchor=west] at (7.45,1.78) {媒体时间（秒）};
  \foreach \xx in {4,7} {
    \draw[contact] (\xx,3.2) -- (\xx,3.65);
    \fill[BookBlue] (\xx,3.4) circle (1.5pt);
  }
  \foreach \xx in {4.24,6.96} {
    \draw[sound] (\xx,2.6) -- ++(-0.1,-0.16) -- ++(0.2,0) -- cycle;
  }
  \node at (4,3.95) {1.00};
  \node at (7,3.95) {2.50};
  \node[anchor=west] at (0.3,1.3) {第一次：声音较晚};
  \node[anchor=west] at (5.75,1.3) {第二次：声音较早};
  \draw[axis,-{Stealth[length=2mm]}] (0.6,0.1) -- (5,0.1);
  \draw[axis,-{Stealth[length=2mm]}] (6,0.1) -- (10.4,0.1);
  \draw[contact] (1.7,-0.08) -- (1.7,0.35);
  \draw[sound] (3.9,-0.08) -- (3.9,0.35);
  \draw[sound] (7.4,-0.08) -- (7.4,0.35);
  \draw[contact] (8.35,-0.08) -- (8.35,0.35);
  \draw[BookInk,-{Stealth[length=2mm]},line width=0.8pt]
    (1.7,0.58) -- (3.9,0.58);
  \draw[BookInk,-{Stealth[length=2mm]},line width=0.8pt]
    (8.35,0.58) -- (7.4,0.58);
  \node[above] at (2.8,0.62) {$+120$ ms};
  \node[above] at (7.88,0.62) {$-20$ ms};
  \node[below] at (1.7,-0.12) {1.00};
  \node[below] at (3.9,-0.12) {1.12};
  \node[below] at (7.25,-0.12) {2.48};
  \node[below] at (8.5,-0.12) {2.50};
  \node[anchor=west] at (0.3,-0.85) {局部放大分别使用自己的时间尺度，刻度单位为秒。};
\end{tikzpicture}

  \caption{两次碰撞的教学时间对齐。上排为可见接触，下排为音频起声，局部放大给出偏移方向与大小；不是模型实测波形。平均值可能掩盖第一次事件超出既定容差，漏声和多声还须单独记录。}
  \label{fig:vision-event-sync}
\end{figure}
\FloatBarrier

逐事件标注成本较高时，可以使用自动估计。MMAudio论文用DeSync指标评价同步：由Synchformer预测整段音画的时间偏移，取其绝对值，单位为秒。这种估计不需要事件标注，但它给出的是整段的单一偏移，不能指出哪一次碰撞出了问题；而且评价用的Synchformer与MMAudio生成时用的同步编码器相同，模型可能恰好迎合了该编码器的偏好。因此自动指标适合在模型之间做大批量比较，最终交付仍应对关键事件做式\eqref{eq:vision-event-sync-error}的逐事件检查。

持续环境声没有唯一冲击起点，应看场景与前后连贯性；对白还需要转写内容、说话人和口型关系。近景同步测量也要说明是否扣除音频设备延迟，远距离声源则可能具有合理的传播延迟。声画不同时出现并不总是错误，判断取决于所描述的物理或创作条件。

调整要针对失败来源。声源不合适可改文字，局部时间错误可用上下文修补重做相应区间，长背景突变可改用上下文延伸或多窗口协调。每次保存改动范围和候选，重新检查与邻段的连接。修改、延伸与修补是不同条件下的选择，不是所有配声都必须经历的固定步骤。

\section{音视频的协同生成}
\label{sec:vision-multigen-joint}

当画面与声音都尚未固定，输入只有文字（有时加一幅首帧），输出是一段带声轨的视频。难点在于两种媒体要描述同一组事件：玩具狗碰到花瓶的那一帧，正是碰撞声开始的时刻；人物开口说话时，嘴型要对应发出的音节。前一节的配声模型只能让声音去适应已经固定的画面，如果画面里的碰撞本身不清楚，声音无法补救。

两条路线对“画面和声音在什么时候互相约束”给出不同回答。分阶段条件传递先生成一种媒体，再以它为条件生成另一种，约束是单向的。联合采样把两种媒体放进同一个生成过程，每一步都让双方读取对方的当前状态，约束是双向的。后一条路线由Ovi和LTX-2实现，二者结构相近，差别在于训练组织和推理时的引导方式。

\subsection{基于分阶段条件传递的生成}
\label{sec:vision-multigen-staged}

先制作玩具狗碰倒花瓶的视频，再按成片配声，可以直接复用前两节的模型。视频阶段交付的不只是帧文件，还应包括实际时长、帧率、事件时间、可见声源以及声音可自由设计的部分。若碰撞从原计划的1.0秒变成1.4秒，配声条件应采用成片中的时间。

制作过程中可把理想化的分布关系写成
\begin{equation}
  p(V,A\mid c)=p(V\mid c)\,p(A\mid V,c),
  \label{eq:vision-staged-factorization}
\end{equation}
其中$V$为视频，$A$为音频，$c$为创作条件。这是概率链式分解，本身没有损失一般性；实际误差来自两个模型的近似、条件传递以及第二阶段受已选视频约束。不能仅从分阶段形式推出它在数学上无法表示联合分布。问题出在第一阶段：$p(V\mid c)$的模型在生成时并不知道之后要配什么声音，它可能生成一个碰撞不清楚的画面，而第二阶段没有办法修改它。

在固定成片上配声时，声音模型只能适应已有碰撞。若画面没有真正接触花瓶，补一声碰撞不会让动作成立。此时应回到视频阶段修改，再更新配声，或者选择允许共同调整的模型。分阶段路线的好处是每个阶段都可以选用最合适的专用模型，并能单独替换和核验；不同模型间共享时间单位和事件记录，可以组成一致的制作流程，但不等于它们共享参数或共同训练。说话场景是这条路线最难处理的情形：先生成的人物嘴部运动没有对应的语音内容，后配的语音很难与嘴型吻合。

\subsection{基于联合采样的生成}
\label{sec:vision-multigen-joint-mechanism}

联合采样把画面和声音都放在待采样状态中。文字应同时说明可见动作与声音要求，训练配对要来自同一时间范围，否则分支之间学到的联系可能是错误的。两个代表模型都采用“每种媒体一个分支、分支之间双向交叉注意力”的结构。Ovi的两个分支结构完全相同，音频分支先单独预训练，再与视频分支融合；LTX-2的两个分支宽度可以不同，结构和训练已在\ref{sec:vision-models-ltx2}节说明，这里只讲它在推理时怎样分别调节文字和另一模态的影响。

\subsubsection{Ovi}
\label{sec:vision-multigen-ovi}

分阶段路线无法让先生成的画面考虑声音。若直接训练一个大网络同时输出两种潜变量，视频和音频的潜空间、序列长度都不同，从头训练还要同时学会两种媒体的生成，成本很高。Ovi的主要想法是让两个分支采用完全相同的结构：视频分支取自已经训练好的视频生成模型，音频分支用同样的结构单独学会生成声音，然后在每一层用双向交叉注意力把两者连接起来，只训练连接和注意力部分\scite{vision-ovi2025}。结构对称使两个分支每一层的隐藏宽度相同，交叉注意力可以直接逐层配对。

视频分支由Wan2.2的50亿参数模型初始化，音频分支结构相同、从头训练，合计约110亿参数。每个分支有30个块，隐藏宽度3072，24个注意力头、每头128维，前馈层宽度14336，配置见表\ref{tab:vision-multigen-ovi-config}。每块包含自注意力、文字交叉注意力、音视频交叉注意力和前馈层。论文列出了每块的这几类子层，但没有写明它们的先后顺序和生成时刻的接入方式。为了说明数据怎样流动，下面按“先自注意力、再两种交叉注意力、最后前馈层”的示意顺序写出视频分支第$l$块，记视频词元为$\bm Z^{\mathrm v}$、同一块中音频分支的隐藏表示为$\bm Z^{\mathrm a}$、文字表示为$\bm C$：
\begin{equation}
 \begin{aligned}
  \bm Z_1&=\bm Z^{\mathrm v}+\operatorname{SA}(\operatorname{Norm}(\bm Z^{\mathrm v})),\\
  \bm Z_2&=\bm Z_1+\operatorname{CA}(\operatorname{Norm}(\bm Z_1),\bm C),\\
  \bm Z_3&=\bm Z_2+\operatorname{CA}(\operatorname{Norm}(\bm Z_2),\bm Z^{\mathrm a}),\\
  \bm Z_4&=\bm Z_3+\operatorname{FFN}(\operatorname{Norm}(\bm Z_3)),
 \end{aligned}
 \label{eq:vision-multigen-ovi-block}
\end{equation}
音频分支对称地以视频表示作为音视频交叉注意力的键和值。式中的顺序与归一化位置是示意，不是论文给出的配置。两个分支宽度相同，第三行的交叉注意力不需要额外的维度投影就能逐层配对；30个块各有一个音视频交叉注意力，所以两种媒体在整个网络深度上反复交换信息。

\begin{table*}[tbp]
 \centering
 \small
 \begin{tabularx}{\linewidth}{@{}p{28mm}Xccccc@{}}
  \toprule
  模型 & 输入单元 & 层数 & 宽度 & 头数 & 前馈宽度 & 参数量 \\
  \midrule
  Ovi视频分支 & Wan2.2视频潜网格的时空块，121帧对应31个潜时间位置 & 30 & 3072 & 24 & 14336 & 约50亿 \\
  Ovi音频分支 & MMAudio 16 kHz自编码器潜序列，五秒157个位置 & 30 & 3072 & 24 & 14336 & 约50亿 \\
  \bottomrule
 \end{tabularx}
 \caption{Ovi两个分支的配置。论文给出全模型约110亿参数、每个分支约50亿参数，两个分支之和与全模型规模之间的差额论文未逐项说明；融合阶段约57亿参数参与更新。}
 \label{tab:vision-multigen-ovi-config}
\end{table*}
\FloatBarrier

输入的视频是121帧、每秒24帧、约$720\times720$像素的片段，约5秒；视频自编码器把它压成31个潜时间位置。音频用MMAudio的16 kHz自编码器，五秒音频按每512个采样一个位置，得到157个潜位置；生成后先解码为频谱，再由BigVGAN声码器恢复波形。文字由一个冻结的T5编码器处理，两个分支共享同一条文字条件。提示中用\code{<S>}与\code{<E>}包住要说出的台词，用\code{<AUDCAP>}与\code{<ENDAUDCAP>}包住对声音的描述，放在提示末尾。

交叉注意力读取另一模态的隐藏特征，需要知道哪些位置对应同一时刻。视频有31个时间位置，音频有157个，二者覆盖同样五秒。若直接用序号作旋转位置编码，第31个音频位置会被当作与第31个视频位置同时，而前者实际只在约1秒处。Ovi把音频的旋转位置编码频率乘以$31/157\approx0.197$：第$j$个音频位置的有效位置变为$0.197j$，第157个位置约为31，与最后一个视频位置对齐。这一缩放只建立时间尺度的对应，正确的敲击同步仍要从配对数据中学习。

训练分三个阶段。第一阶段只训练音频分支：用数十万小时、以语音为主的音频，长度最多12秒，学习文字到声音的生成；论文报告训练5万步，批量2880，学习率$10^{-4}$。第二阶段把音频长度固定为约5秒，并加入音效数据微调，使音频分支与视频片段长度一致。第三阶段融合：新建的音视频交叉注意力从随机初始化开始，两个分支的前馈网络全部冻结，只训练自注意力和交叉注意力，110亿参数中约57亿参与更新，训练4万步，批量768，学习率$5\times10^{-5}$。配对音视频提供两种干净目标$\bm z_1^{\mathrm v}$和$\bm z_1^{\mathrm a}$，各自抽取独立噪声，但使用同一个生成时刻$s$。损失是两个流匹配目标的加权和：
\begin{equation}
  L_{\mathrm{joint}}
  =w_{\mathrm v}\left\lVert\bm v^{\mathrm v}-(\bm z_1^{\mathrm v}-\bm z_0^{\mathrm v})\right\rVert_2^2
   +w_{\mathrm a}\left\lVert\bm v^{\mathrm a}-(\bm z_1^{\mathrm a}-\bm z_0^{\mathrm a})\right\rVert_2^2,
  \label{eq:vision-ovi-objective}
\end{equation}
其中$w_{\mathrm v},w_{\mathrm a}$由论文设定。共享$s$的原因是推理时两个分支要沿同一日程前进；若训练时两边噪声水平不同，交叉注意力学到的是“干净程度不同的两种媒体”之间的关系，与推理时的情况不符。

推理时，把联合状态视为一个较长的向量，其演化可写为
\begin{equation}
 \frac{\mathrm d}{\mathrm ds}
 \begin{pmatrix}\bm z^{\mathrm v}\\ \bm z^{\mathrm a}\end{pmatrix}
 =
 \begin{pmatrix}
 \bm v^{\mathrm v}(\bm z^{\mathrm v},\bm z^{\mathrm a},s,c)\\
 \bm v^{\mathrm a}(\bm z^{\mathrm a},\bm z^{\mathrm v},s,c)
 \end{pmatrix}\text{。}
 \label{eq:vision-joint-flow}
\end{equation}
每边速度都依赖两边当前状态。采用欧拉法说明时，两个下一步都应由同一个当前状态计算，再共同更新；若先把视频推进一步，然后用新视频与旧音频计算音频速度，就已换成另一种分步更新。论文用UniPC求解器共同积分两个分支，并报告它比欧拉法稳定；无分类器引导在每一步增加一次去条件前向。

教学上取单维潜状态$z^{\mathrm v}=0.2$、$z^{\mathrm a}=-0.1$，假设当前预测速度分别为0.6和0.4，步长0.25。一次共同欧拉更新得到0.35和0。这个计算只说明数值更新的同步，没有证明媒体里的敲击已同步。生成时刻$s$走到1以后，还需分别解码视频和音频，按媒体时间组合播放。

请求中可以同时规定“玩具狗接触花瓶”“短促金属碰撞声”和“安静房间环境”。结果应检查声源、动作及起止是否属于同一事件。说话场景还要检查台词、音素和嘴部运动，不能用敲击同步代替口型验证。Ovi原版聚焦约五秒片段；将多段分别成功的片段拼成一分钟故事，仍需另外解决跨镜头身份、场景与叙事状态。

\subsubsection{LTX-2}
\label{sec:vision-multigen-ltx}

LTX-2的视频与音频分支、双向交叉注意力及联合解码已在\ref{sec:vision-models-ltx2}节说明，两个分支的层数与宽度等配置以该节为准：视频和音频各由自己的自编码器压成潜表示，两个分支通过双向交叉注意力交换信息，跨模态连接主要使用时间位置，训练用流匹配监督两路速度，文字条件来自冻结的Gemma3隐表示\scite{vision-ltx22026}。与Ovi相同，它在每个生成步共同更新两路状态。这里讨论使用时的一项差别：推理时可以分别增强文字条件和另一模态带来的影响，使用差异发生在每一步的预测组合中，不是生成完成后才为声音与画面打分。

固定当前模态的潜状态及生成时刻，记$\bm f(c,m)$为同时给定文字$c$和另一模态信息$m$时的预测，$\bm f(\varnothing,m)$为去掉文字的预测，$\bm f(c,\varnothing)$为去掉另一模态信息的预测。报告的引导组合可写为
\begin{equation}
 \widehat{\bm f}
 =\bm f(c,m)
  +\lambda_c[\bm f(c,m)-\bm f(\varnothing,m)]
  +\lambda_m[\bm f(c,m)-\bm f(c,\varnothing)]\text{。}
 \label{eq:vision-modality-guidance}
\end{equation}
两项差分别刻画加入文字与加入另一模态后当前预测如何变化。这里把完整条件预测作为基准，故$\lambda_c=\lambda_m=0$仍保留两种条件，不能误读为无条件生成。去条件分支也应使用方法规定的缺省入口，不是任意填一幅黑图就等价。每一步因此需要三次前向，计算量约为不加引导时的三倍。

假设某个潜分量三次预测依次为0.6、0.2、0.5，取$\lambda_c=2$、$\lambda_m=1$，组合值为$0.6+2(0.6-0.2)+(0.6-0.5)=1.5$。它是潜空间预测，可以超出三次预测的范围，不是必须在0至1之间的像素。提高引导会改变采样轨迹，过大时画面可能过饱和或动作僵硬，是否改善当前请求需比较完整结果。

调节时，可以在同一“玩具狗靠近花瓶并发出机械声”请求下固定初始噪声、长度与采样步数，只改变$\lambda_m$，查看声音是否更贴近动作、画面是否出现额外变化。随后再在多组噪声下重复，避免单个候选主导判断。若文字与首帧冲突，增强两个条件也不自动解决冲突，需要先明确优先保留哪一个。高分辨率制作可以先生成基础潜变量，再做空间潜上采样，并对重叠时空区块细化后合并。两路潜变量最终分别恢复为视频帧与音频频谱，频谱经声码器形成波形。容器层仍要写入正确时长、采样率与时间戳；如果导出时错误改变帧率，模型内部的时间联系也会被破坏。分块降低资源压力，不提供无限长叙事的一致性保证。

三种协同方式的比较见表\ref{tab:vision-joint-compare}。Ovi相对分阶段路线把单向约束改成每步的双向读取，换来了画面能考虑声音、说话场景能联合生成嘴型与语音，代价是需要大规模配对音视频数据和双倍的主干计算。LTX-2保留了联合采样，又在推理时把两种条件的强度拆开调节，代价是每步三次前向。

\begin{table*}[tbp]
 \centering
 \small
 \begin{tabularx}{\linewidth}{@{}p{22mm}XXX@{}}
  \toprule
  比较项 & 分阶段条件传递 & Ovi & LTX-2 \\
  \midrule
  约束方向 & 视频约束音频，单向 & 每层双向交叉注意力 & 每层双向交叉注意力 \\
  分支结构 & 两个独立模型 & 两个结构相同的分支，各约50亿参数 & 宽度可以不同的两个分支 \\
  时间对应 & 共享事件记录与时间单位 & 音频旋转位置编码按$31/157$缩放 & 跨模态连接使用时间位置 \\
  训练 & 各自训练，可选专用模型 & 音频分支先预训练，融合时冻结前馈网络 & 两路速度共同监督 \\
  推理调节 & 可单独重做任一阶段 & 共同求解，标准无分类器引导 & 文字与模态引导分别加权 \\
  主要局限 & 画面无法考虑声音，口型难对 & 约五秒片段，需配对数据 & 每步三次前向 \\
  \bottomrule
 \end{tabularx}
 \caption{音视频协同生成方式的比较。三者面对同一文字请求，差别在于两种媒体何时、以何种方向互相约束。}
 \label{tab:vision-joint-compare}
\end{table*}
\FloatBarrier

\section{多视角与三维内容生成}
\label{sec:vision-multigen-3d}

空间内容生成的输出是一个可以从不同方向观察的对象或场景，形式可以是辐射场、网格或高斯集合。难点是多视角一致：各个方向的画面必须由同一个三维表面产生。二维生成模型从海量图像学到了很强的外观先验，却没有被要求让不同视角互相一致；而带三维真值的数据远少于图像数据。另一难点是区分证据：哪些表面由观测支持，哪些由先验补出。

本节用到两组已有材料，这里只回顾结论。第一组是\ref{chap:vision-image-formation}章的相机几何：针孔相机用内参和外参把三维点投影到像素；\ref{sec:vision-formation-two-view}节说明同一三维点在两个视图中的投影受对极约束，相机已知时可以三角化出点的位置；\ref{sec:vision-formation-multi-view}节的运动恢复结构与光束法平差在多幅图中同时估计相机与三维点。这些结论告诉我们，“多视角一致”有一个可检查的几何含义：同一表面点在各视图中的投影必须落在相机几何规定的位置上。第二组是\ref{sec:vision-generation-novel-view}节的场景表示：NeRF用网络表示空间中每一点的密度和颜色，沿相机射线做体渲染得到像素（\ref{sec:vision-generation-nerf}节）；3D Gaussian Splatting用一组带中心、协方差、不透明度和颜色的三维高斯表示场景，投影到图像平面后按深度做透明度合成（\ref{sec:vision-generation-gaussians}节）。两种表示都是可微的，所以可以用渲染图与目标图之间的误差反向更新三维参数。本节的三种方法都建立在这两组材料上，差别在于目标图从哪里来：来自二维生成模型的判断、来自二维生成模型一次生成的多视角图，还是来自真实观测。

本节按输入能提供多少观测划分小节，每种条件对应一种代表方法，它们对“怎样把二维先验和三维一致性结合”给出不同回答。没有参考观测时，DreamFusion保留一个三维表示，反复从随机视角渲染，用冻结的二维扩散模型给每幅渲染图提供修改方向，一致性由共享的三维表示保证。只有单幅参考时，Wonder3D先让二维扩散模型一次生成多个视角的颜色和法线图，在生成阶段就用跨视角注意力协调一致性，再把这些图融合成表面。有多幅或连续观测时，4D Gaussian Splatting不再需要生成先验，直接用观测拟合一个随时间变化的高斯场景。三者都可能输出渲染图，但展示用的图像、可编辑网格和物理模拟模型有不同的验收要求。

\subsection{无参考观测下的生成}
\label{sec:vision-multigen-no-reference}
\label{sec:vision-multigen-dreamfusion}

文字“陶瓷玩具狗”没有给出目标对象的真实样子。若用二维生成模型分别画出正面、侧面和背面，每幅都可能画出一张脸，或者耳朵形状各不相同，因为每次生成都独立进行。训练一个直接输出三维的生成模型又缺少足够的带文字三维数据。DreamFusion的主要想法是：为每条文字请求从随机初始化的三维辐射场出发，反复渲染随机相机视图，用冻结的二维扩散模型判断渲染图应向哪个方向改，再把这个方向传回三维参数\scite{vision-dreamfusion2023}。不同视角共享同一个辐射场，一致性由表示本身保证；外观知识则全部来自二维模型，不需要任何三维训练数据。它的\term{分数蒸馏采样}（score distillation sampling，SDS）就是把二维先验偏好的变化传回三维参数的计算方法。

三维表示基于mip-NeRF 360的辐射场：一个MLP读取空间位置，输出体密度和反照率，即材料本身的颜色倾向，不含光照。这个MLP由5个残差块组成，隐藏宽度128，使用Swish激活，块间加层归一化；输入位置先经最高频率为$2^{7}$的8级正弦位置编码，并按mip-NeRF的做法对每个射线区间积分。密度经指数函数、反照率经sigmoid输出，密度上还加一个以原点为中心的高斯形偏置，使优化初期物体集中在场景中央。另有一个小MLP读取射线方向，输出环境背景，按累积不透明度与前景合成。渲染时先由密度的梯度求表面法线，再用随机采样的点光源对反照率做着色，体渲染见\ref{sec:vision-generation-nerf}节。着色迫使几何本身显出凹凸，若模型只把凹凸“画”在颜色上，换一个光照方向就会露馅。二维先验是Imagen的$64\times64$基础模型，以T5-XXL文字嵌入为条件，不使用其超分辨率级联，参数在整个过程中固定。这里设场景参数为$\bm\theta$，相机和光照条件为$\pi$，渲染结果展平为$\bm x=g_{\bm\theta}(\pi)\in\mathbb R^d$。噪声水平用$\ell$表示，加入噪声得到$\bm z_\ell=\alpha_\ell\bm x+\sigma_\ell\bm\epsilon$，其中$\bm\epsilon\sim\mathcal N(\bm0,\bm I)$。冻结模型在文字$c$下预测噪声$\widehat{\bm\epsilon}_\phi$。

SDS给出的参数梯度估计为
\begin{equation}
 \bm g_{\mathrm{SDS}}
 =w(\ell)\bm J_g^\top
     [\widehat{\bm\epsilon}_\phi(\bm z_\ell,\ell,c)
       -\bm\epsilon],
 \qquad
 \bm J_g=\frac{\partial g_{\bm\theta}(\pi)}
                 {\partial\bm\theta}\text{。}
 \label{eq:vision-sds-gradient}
\end{equation}
若参数数目为$p$，$\bm J_g\in\mathbb R^{d\times p}$，噪声差为$d$维，结果是$p$维参数梯度。白话地说，二维模型看到加噪后的渲染图，预测“其中哪些是噪声”；预测噪声与真正加入的噪声之差，指出渲染图中不像“陶瓷玩具狗”的成分，渲染Jacobian再把这个像素方向换算成参数方向。优化器沿其负方向更新场景。二维模型参数$\phi$保持固定，计算中也不对噪声预测网络的输入Jacobian反向传播。

后一限制不是普通噪声平方误差求导后的自然结果。若直接最小化
$\lVert\widehat{\bm\epsilon}_\phi(\bm z_\ell,\ell,c)-\bm\epsilon\rVert_2^2$，
链式法则还会出现噪声网络对$\bm z_\ell$的导数，它需要穿过整个大网络反传，计算昂贵且在小噪声水平下数值不稳定。SDS通过概率分布之间的目标获得另一种方向，可以由下面的简化推导看清。

固定$\pi$和$\ell$，令$q_{\bm\theta}$为均值$\alpha_\ell g_{\bm\theta}(\pi)$、方差$\sigma_\ell^2\bm I$的高斯分布，$p_\phi(\bm z_\ell\mid c)$为扩散先验在该噪声水平的条件分布。我们希望渲染图加噪后的分布落在先验分布的高概率区域，即减小$\operatorname{KL}(q_{\bm\theta}\|p_\phi)$。因为$q_{\bm\theta}$的方差不变，熵不随均值改变，KL散度对$\bm\theta$的梯度只来自交叉熵项$-\mathbb E_{q}[\log p_\phi]$。利用重参数化$\bm z_\ell=\alpha_\ell g_{\bm\theta}(\pi)+\sigma_\ell\bm\epsilon$把期望移到固定分布的$\bm\epsilon$上，再对$\bm\theta$求导，得到
\[
 \nabla_{\bm\theta}\operatorname{KL}(q_{\bm\theta}\|p_\phi)
 =-\alpha_\ell\bm J_g^\top
     \mathbb E_{\bm\epsilon}
       [\nabla_{\bm z_\ell}\log p_\phi(\bm z_\ell\mid c)]\text{。}
\]
用噪声网络提供的分数近似
$\nabla_{\bm z_\ell}\log p_\phi\approx
-\widehat{\bm\epsilon}_\phi/\sigma_\ell$，
再乘以权重$w(\ell)\sigma_\ell/\alpha_\ell$，得到
$w(\ell)\bm J_g^\top\mathbb E[\widehat{\bm\epsilon}_\phi]$。
固定相机和噪声水平时，$\bm J_g$不依赖抽样噪声，而$\mathbb E[\bm\epsilon]=\bm0$，因此可减去$\bm\epsilon$作为零均值控制变量，形成式\eqref{eq:vision-sds-gradient}的估计。减去$\bm\epsilon$不改变期望，但当网络预测接近真实噪声时，单次估计的方差明显变小。

这一推导要求所用噪声水平的$\alpha_\ell,\sigma_\ell$非零，并把学习到的分数当作分布分数的近似。实际引导和网络误差会影响方向，不能因此声称每步都精确降低真实三维分布的KL散度。实现可对$w(\ell)(\widehat{\bm\epsilon}_\phi-\bm\epsilon)$停止梯度，再让它与渲染图像作内积求导；直接对噪声差的平方求导会得到不同计算。

用一个局部线性教学渲染器检查方向。设$\bm\theta=(a,b)^\top$，两个像素分量为$g_{\bm\theta}=(a,a+b)^\top$，则
$\bm J_g=\left(\begin{smallmatrix}1&0\\1&1\end{smallmatrix}\right)$。
若权重为1、噪声差为$(0.2,-0.1)^\top$，参数梯度为$(0.1,-0.1)^\top$。从$(0.5,0.2)^\top$以步长0.1更新，得到$(0.49,0.21)^\top$，渲染变为$(0.49,0.70)^\top$。第二像素同时受$a,b$影响，两者变化抵消，说明像素方向不能不经渲染Jacobian直接当作参数方向。图\ref{fig:vision-sds-path}标出了梯度经过和停止的位置。

\begin{figure*}[tbp]
  \centering
  \begin{tikzpicture}[x=1cm,y=1cm,font=\figurefont\small,
  block/.style={draw=FigureModel,fill=FigureModel!10,line width=0.8pt,
    minimum height=10mm,minimum width=23mm,align=center},
  input/.style={block,draw=FigureInput,fill=FigureInput!10},
  arr/.style={-{Stealth[length=2mm]},BookInk,line width=0.8pt,
    rounded corners=1.2mm},
  grad/.style={-{Stealth[length=2mm]},FigureLoss,dashed,line width=1pt,
    rounded corners=1.2mm}]
  \node[input] (theta) at (1.2,2) {三维参数\\$\bm\theta$};
  \node[block] (render) at (4.5,2) {渲染器\\$g_{\bm\theta}(\pi)$};
  \node[block] (noise) at (7.8,2) {加噪\\$\alpha_\ell\bm x+\sigma_\ell\bm\epsilon$};
  \node[block] (prior) at (11.1,2) {冻结二维先验\\$\widehat{\bm\epsilon}_\phi$};
  \node[block] (res) at (14.6,2) {噪声差\\$\widehat{\bm\epsilon}_\phi-\bm\epsilon$};
  \draw[arr] (theta) -- (render);
  \draw[arr] (render) -- node[above] {$\bm x$} (noise);
  \draw[arr] (noise) -- (prior);
  \draw[arr] (prior) -- (res);
  \node (camera) at (4.5,3.5) {相机与光照$\pi$};
  \node (eps) at (7.8,3.5) {噪声$\bm\epsilon$};
  \node (text) at (11.1,3.5) {文字$c$};
  \draw[arr] (camera) -- (render);
  \draw[arr] (eps) -- (noise);
  \draw[arr] (eps.north) -- (7.8,4.25) -- (14.6,4.25) -- (res.north);
  \draw[arr] (text) -- (prior);
  \node[block,draw=FigureLoss,fill=FigureLoss!8,minimum width=35mm]
    (jac) at (4.5,-0.15) {通过渲染器求导\\$w(\ell)\bm J_g^\top(\widehat{\bm\epsilon}_\phi-\bm\epsilon)$};
  \draw[grad] (res.south) -- (14.6,-0.15) -- node[above] {噪声差停止梯度} (jac.east);
  \draw[grad] (jac.west) -- (1.2,-0.15) -- (theta.south);
  \node[anchor=west] at (7.1,-0.8) {不反传噪声网络的输入Jacobian};
\end{tikzpicture}

  \caption{SDS把冻结二维先验的噪声差传回三维参数。实线为渲染、加噪与预测，虚线为通过渲染器的参数梯度。噪声差在求参数梯度时视为常量，不反传冻结噪声网络的输入Jacobian；每次重新抽取相机，使多个方向约束同一个场景。}
  \label{fig:vision-sds-path}
\end{figure*}
\FloatBarrier

完整的优化过程对每条文字请求单独进行，没有跨请求的训练阶段。每次迭代随机抽取相机位置、焦距和点光源，渲染$64\times64$图像，从$[0.02,0.98]$中均匀抽取噪声时刻，计算式\eqref{eq:vision-sds-gradient}并更新辐射场参数；论文取权重$w(\ell)=\sigma_\ell^2$，并报告均匀权重效果相近。相机仰角在$-10^\circ$到$90^\circ$之间、方位角在$0^\circ$到$360^\circ$之间抽取，相机到中心的距离在1到1.5之间。论文共优化15000步，用Distributed Shampoo优化器，在4块TPUv4芯片上约需1.5小时。除SDS外，还加入表面朝向正则项，惩罚背向相机的法线，以及控制累积不透明度、抑制半透明浮游物的正则项。论文使用很大的无分类器引导权重100，这使渲染结果更贴近文字，但也使颜色偏饱和、几何偏平滑。模型还按相机方位角在文字后附加“正面视角”“侧面视角”“背面视角”等描述，减少把背面也拉向正面外观的冲突。

优化完成后，辐射场可以从不同相机渲染；若要交付网格，还须提取表面并整理材质。多个正面、几何过平滑、细结构缺失或依赖视角的伪外观，都可能在环绕观察中出现。多个正面的根源在于二维先验对每幅渲染图独立判断“像不像玩具狗”，而训练图像中正面照最多，视角描述只能部分抑制这一偏好。每条请求都要重新优化，一次生成按论文设置约需1.5小时。可渲染性没有证明表面可直接编辑，生成的背面也不是任何真实玩具的背面。

\subsection{单幅参考下的生成}
\label{sec:vision-multigen-single-reference}
\label{sec:vision-multigen-wonder3d}

单幅照片增加了可见轮廓、颜色和主体身份的约束，但背面和遮挡区域仍未知。创作前可约定正面保持原耳朵、眼睛与花纹，允许设计背面纹理。这样才能把偏离参考的错误，与合理补全但无法验证的内容区分开。

DreamFusion式的逐请求优化在这里有两个不足：每个对象都要优化很久，而且二维先验对每幅渲染图独立打分，多视角一致性要靠长时间优化慢慢形成。另一种做法是先让二维扩散模型一次生成多个固定视角的颜色图，再从这些图重建三维。但只生成颜色图，重建时仍可能把纹理变化解释成凹凸，因为颜色不直接说明表面朝向。Wonder3D的主要想法是同时生成多视角颜色图和\term{法线图}（normal map），让外观与表面朝向在生成阶段就互相约束，再用一个快速的表面优化把它们融合\scite{vision-wonder3d2024}。法线图的每个有效像素给出该处表面的朝向向量，它不直接给出深度或三维点位置。

生成部分以单幅$256\times256$参考图为条件，输出前、后、左、右、右前、左前六个固定视角的颜色图和法线图，共12幅。去噪网络由Stable Diffusion Image Variations模型初始化，它本来就以图像为条件生成图像。输出域由一个一维的域标记说明当前生成的是颜色还是法线，该标记经位置编码后与生成时刻嵌入拼接，进入U-Net。原来的自注意力被扩展为跨视角注意力：一个视角的查询可以读取全部六个视角的键和值，使不同相机下的表示互相协调。每个Transformer块中还在交叉注意力之前加入跨域注意力，把同一视角的颜色与法线表示的键和值合并。仅有域标记能告诉网络输出哪类数据，却不能单独保证两类结果描述同一个表面，跨域注意力正是为此加入的。

两种注意力的差别在于键和值从哪里取。把12幅图在U-Net某一层的特征记为$\bm H_{k,u}$，其中$k=1,\dots,6$为视角，$u\in\{\mathrm{rgb},\mathrm{n}\}$为颜色或法线域，每幅图有$m$个空间位置。多视角注意力中，视角$k$、域$u$的查询读取同一域全部六个视角的键和值：
\begin{equation}
 \operatorname{MVAttn}(\bm H_{k,u})
 =\operatorname{Attn}\bigl(\bm Q_{k,u},
   [\bm K_{1,u};\dots;\bm K_{6,u}],
   [\bm V_{1,u};\dots;\bm V_{6,u}]\bigr),
 \label{eq:vision-multigen-wonder-mv}
\end{equation}
键值序列长$6m$。跨域注意力中，查询读取同一视角两个域的键和值：
\begin{equation}
 \operatorname{CDAttn}(\bm H_{k,u})
 =\operatorname{Attn}\bigl(\bm Q_{k,u},
   [\bm K_{k,\mathrm{rgb}};\bm K_{k,\mathrm n}],
   [\bm V_{k,\mathrm{rgb}};\bm V_{k,\mathrm n}]\bigr),
 \label{eq:vision-multigen-wonder-cd}
\end{equation}
键值序列长$2m$。前者没有显式使用对极约束，相机之间的对应由注意力权重从数据中学出；后者让同一视角的颜色与法线互相参照。以$32\times32$的特征层为例，$m=1024$，多视角注意力的键值长度为6144，普通自注意力只有1024，代价随之增加。论文说明跨域注意力放在U-Net每个Transformer块的交叉注意力之前，多视角注意力由原自注意力扩展而来；注意力头数等其余配置沿用初始化模型，论文未单列。

训练使用Objaverse中经过清理的约三万个三维对象。对每个对象从固定视角渲染颜色图和法线图，以某一视图作为输入参考，其余视图和法线作为目标，用标准的去噪扩散目标微调网络。论文报告训练3万步，总批量512，在8块A800 GPU上约需三天；训练沿用初始化模型的噪声预测目标与优化器设置。采样时从12幅图的噪声共同生成结果，未观测的背面由生成先验补充。六幅看起来各自合理的图，仍可能在细耳朵、尾巴和遮挡位置互相矛盾。

形成几何时，方法优化一个\term{有符号距离场}（signed distance field，SDF）：空间点到表面的距离带有内外符号，零值集合定义表面。若约定物体外部为正，在可微且梯度非零处，
$\bm n(\bm p)=\nabla f(\bm p)/\lVert\nabla f(\bm p)\rVert_2$
给出向外单位法线。法线由SDF对空间的一阶导数得到；优化涉及法线的损失时，还要对场参数继续求导，不能把“求法线”本身称为计算SDF的二阶空间导数。SDF用instant-NGP式的多分辨率哈希网格表示，以加快优化。

从各生成视图抽取像素射线，用SDF渲染颜色、法线和前景占据，再与相应生成图比较。颜色用平方误差，前景掩码用二元交叉熵，法线用下面的加权损失；梯度范数接近1的正则项保持距离场性质，稀疏与平滑项限制浮游表面和局部噪声。论文把总目标写成这六项之和，各项权重未在文中给出。这些是对生成多视图的融合，不是把六张图直接粘成六面体。

生成法线也可能错误，因此原方法依据观察方向赋权。设$\bm v_i$是从相机射入表面的单位方向，$\bm n_i$是生成的向外单位法线，$\widehat{\bm n}_i$是当前渲染法线。几何法线损失为
\begin{equation}
 L_{\mathrm n}
 =\frac{\sum_i w_i(1-\widehat{\bm n}_i^\top\bm n_i)}
         {\sum_i w_i},
 \quad
 w_i=
 \begin{cases}
  0,&\bm v_i^\top\bm n_i>\eta,\\
  \exp(|\bm v_i^\top\bm n_i|),&\bm v_i^\top\bm n_i\leq\eta,
 \end{cases}
 \label{eq:vision-wonder-normal}
\end{equation}
其中$\eta$为接近0的负阈值，且当前批次要求$\sum_iw_i>0$。相机看到的向外表面法线应大致逆着入射视线，点积应为负；接近掠射的法线获得较低权重，朝向明显矛盾的结果被排除。全被排除时不能除以零，应跳过该项或重新取样。

取两个教学射线，方向点积为$-1$和$-0.5$，法线差异$1-\widehat{\bm n}_i^\top\bm n_i$分别为0.1和0.3，阈值取$-0.01$。两项权重为$e$和$e^{0.5}$，归一化损失约0.17551，比等权平均0.2更靠近正视方向的误差。它说明赋权的含义，不证明正视法线一定正确。颜色和掩码误差还采用剔除一部分最大残差的策略，降低局部不一致的影响，但可能同时忽略真实细节，须观察最终表面。论文报告整个流程约两分钟即可得到带纹理的网格，远快于逐请求的SDS优化。

输入玩具狗正面时，可以检查输入侧轮廓是否保持，再环绕检查耳朵是否连接、纹理是否有接缝。六个固定视角之间的斜后方没有生成图约束，那里的表面只由SDF的平滑性插出。若没有真实尺度或额外视图，背面纹理与绝对尺寸都不具有测量依据。用于展示的成果也许已经足够，用于动画绑定或制造则还要检查拓扑、厚度与材质等要求。

\subsection{多幅与连续观测下的生成}
\label{sec:vision-multigen-observed-space}
\label{sec:vision-multigen-4dgs}

多幅观测能限制更多表面，前提是相机和时间关系正确。静态场景的拟合与新视角合成已在\ref{sec:vision-generation-novel-view}节展开。若玩具狗正在运动，就不能把不同时间的腿部位置全部拟合为同一静态表面；若为每个时刻单独拟合一组高斯，每个时刻只有少数相机看到，表面会严重欠约束，存储也随帧数线性增长。需要一个随时间变化、又能在时刻之间共享信息的场景表示。

Wu等人的4D Gaussian Splatting的主要想法是：保留一组所有时刻共享的基准三维高斯，用一个紧凑的形变场根据位置和时间预测每个高斯的变化，再渲染相应时刻\scite{vision-4dgs2024}。共享基准使每个时刻都能借用其他时刻的观测，形变场则只需表示变化部分。这里的“四维”指三维空间随时间变化，采用的是该文的时空平面与形变路线，不与其他同名的四维高斯参数化混用。

静态高斯的中心、旋转、尺度、不透明度和颜色，以及投影与透明度合成，见\ref{sec:vision-generation-gaussians}节。形变模块读取基准中心$\bm\mu_i=(x_i,y_i,z_i)^\top$及查询时间$\tau$，在$(x,y)$、$(x,z)$、$(y,z)$、$(x,\tau)$、$(y,\tau)$和$(z,\tau)$六种二维平面上取特征。每个平面用四个邻近网格点双线性插值，因此不用存储完整四维稠密网格。论文的平面基础分辨率为64，并在此基础上再上采样2倍和4倍形成多分辨率（D-NeRF合成数据只用2倍一级）。若每个方向都取64个网格点，四维稠密网格有$64^4\approx1.7\times10^7$个格点，六个平面合计只有$6\times64^2=24576$个；这里假设时间方向也取64个格点只是为了比较，论文没有单独给出时间方向的分辨率，也没有给出平面特征的通道数。

在同一分辨率下，六个平面特征逐元素相乘；不同分辨率的结果拼接，再由一个小型MLP $\phi_{\mathrm d}$合并为形变特征$\bm f_{\mathrm d}$，最后由三个分开的MLP头分别输出位置、旋转和尺度的变化。设$\mathcal P$为上述六组坐标对，$r$标记分辨率，可写
\begin{equation}
 \begin{aligned}
 \bm h_i(\tau)
 &=\mathop{\operatorname{concat}}_r
   \left[
    \mathop{\bigodot}_{(a,b)\in\mathcal P}
       \operatorname{interp}
        (\bm F_{a,b,r},a_i,b_i)
   \right],\\
 \bm f_{\mathrm d}&=\phi_{\mathrm d}(\bm h_i(\tau)),\\
 \Delta\bm\mu_i(\tau)&=\phi_{\mathrm x}(\bm f_{\mathrm d}),\quad
 \Delta\bm r_i(\tau)=\phi_{\mathrm r}(\bm f_{\mathrm d}),\quad
 \Delta\bm s_i(\tau)=\phi_{\mathrm s}(\bm f_{\mathrm d})\text{。}
 \end{aligned}
 \label{eq:vision-4d-planes}
\end{equation}
当坐标对含$\tau$时，式中的相应坐标取查询时间；$\bm F_{a,b,r}$表示带特征通道的网格数组，$\Delta\bm r_i$与$\Delta\bm s_i$是旋转和尺度参数的变化量。论文没有给出$\phi_{\mathrm d}$与三个头的层数和宽度，只在HyperNeRF数据上说明把解码器宽度改为128。相邻位置会共享网格参数，使附近高斯的形变具有共同表示，但不会强制它们属于同一个刚体。

中心按$\bm\mu_i(\tau)=\bm\mu_i+\Delta\bm\mu_i(\tau)$更新，旋转和尺度也在相应参数化中加入预测变化。渲染前还须把旋转参数恢复为合法旋转、尺度恢复为正值，不能把任意残差直接当作正定协方差。原路线的高斯不透明度与颜色参数仍来自基准集合，查询相机时再进行投影与透明度合成。

教学上，一个中心原为$(0,0,2)^\top$，在两个时刻预测的位移分别为$(0,0,0)^\top$和$(0.2,0,0)^\top$。若固定针孔相机焦距为100像素、主点横坐标为0，投影横坐标从0变为$100\times0.2/2=10$像素。纵深、旋转和尺度变化还会影响投影足迹，不能只把二维图中的点平移而忽略高斯形状。

训练数据是带相机参数和时间戳的多视角视频或单相机连续视频。训练先不使用形变场，用3000次迭代优化静态高斯获得初始场景，再联合学习基准表示与形变场；D-NeRF合成数据共训练20000次迭代，Neu3D多相机数据为14000次。每次选择观测相机及时间，渲染当时的高斯，以对应图像的L1像素误差作为监督，并对时空平面施加基于网格的全变差正则，使形变在空间和时间上连续。论文报告在单块RTX 3090上，D-NeRF合成场景以$800\times800$分辨率渲染可达每秒82帧，存储约18 MB。使用时输入任意查询相机和时间，生成该位置的视图；基准高斯、平面特征和小型网络共同构成需要保存的场景。

验收应把已观测时间、未观测视角和观测间时间分开。训练相机下拟合良好，可能仍在背面或快速运动时失败。单相机连续观测还把视角与时间交织在一起，某些表面变化难以唯一解释。高斯有基准编号也不等于它们具有真实物体的持久身份，形变网络更不是含质量、力和碰撞约束的物理模拟器。

三种方法的组件配置列在表\ref{tab:vision-multigen-3d-config}中。三者都不是单一的Transformer，表的列因此按组件而不是按层数、宽度排列。

\begin{table*}[tbp]
 \centering
 \small
 \begin{tabularx}{\linewidth}{@{}p{22mm}XXX@{}}
  \toprule
  组件 & DreamFusion & Wonder3D & 4D Gaussian Splatting \\
  \midrule
  三维表示 & mip-NeRF 360式MLP：5个残差块、宽128、8级积分位置编码；另有背景MLP & instant-NGP式哈希网格SDF & 基准三维高斯集合 \\
  生成或形变网络 & 冻结的Imagen $64\times64$基础模型，T5-XXL文字条件 & 由Stable Diffusion Image Variations初始化的U-Net，加多视角与跨域注意力 & 六平面特征（基础分辨率64，再上采样2倍、4倍）、合并MLP与三个形变头 \\
  分辨率 & 渲染$64\times64$ & 参考与生成图$256\times256$，6视角$\times$2域 & 随数据集，例如D-NeRF为$800\times800$ \\
  训练或优化 & 每条请求15000步 & 生成网络3万步，批量512；表面融合每个对象单独优化 & 静态预热3000次，再联合优化至14000至20000次 \\
  报告耗时 & 4块TPUv4约1.5小时 & 训练8块A800约三天；单对象约两分钟 & 单块RTX 3090，渲染每秒30至82帧 \\
  \bottomrule
 \end{tabularx}
 \caption{三种空间内容生成方法的组件配置。数值取自各论文；注意力头数、平面特征通道数、形变MLP层数等论文未给出的项不列。}
 \label{tab:vision-multigen-3d-config}
\end{table*}
\FloatBarrier

三种条件下的方法比较见表\ref{tab:vision-3d-compare}。从DreamFusion到Wonder3D，改变的是一致性在哪里形成：前者在优化中靠共享表示逐步形成，后者在生成时就靠跨视角注意力协调，换来了从小时到分钟的速度，代价是需要三维资产渲染数据训练，且只覆盖六个固定视角。4D Gaussian Splatting去掉了生成先验，结果有观测支持，但它只能重现被拍到的内容，不能为未拍到的背面创作外观。

\begin{table*}[tbp]
 \centering
 \small
 \begin{tabularx}{\linewidth}{@{}p{22mm}XXX@{}}
  \toprule
  比较项 & DreamFusion & Wonder3D & 4D Gaussian Splatting \\
  \midrule
  输入条件 & 文字 & 单幅参考图 & 多视角或连续观测及相机 \\
  三维表示 & 辐射场（密度与反照率） & 有符号距离场 & 基准高斯加时空平面形变场 \\
  一致性来源 & 共享表示，随机视角反复优化 & 跨视角与跨域注意力，再融合表面 & 观测像素与共享基准 \\
  先验与训练 & 冻结二维扩散模型，无三维数据 & 用三维资产渲染微调二维扩散模型 & 无生成先验，逐场景拟合 \\
  单次耗时 & 逐请求优化，以小时计 & 约两分钟 & 逐场景优化 \\
  未见区域依据 & 全部来自先验 & 背面来自先验 & 无观测处缺乏依据 \\
  \bottomrule
 \end{tabularx}
 \caption{三种观测条件下空间内容生成的比较。耗时依据各论文报告，硬件不同，只说明量级。}
 \label{tab:vision-3d-compare}
\end{table*}
\FloatBarrier

\section{视觉条件下的文字与语音响应}
\label{sec:vision-multigen-response}
\label{sec:vision-multigen-omni-response}

当用户询问“哪个物体先碰到桌面”，要生成的是基于视觉证据的回答，以及说出回答的声音。答案是否有依据属于\ref{chap:vision-multimodal-understanding}章的问题；这里关注文字内容如何转成语音，以及用户何时能听到它。输入是视频、声音和问题，输出是文字与语音两路回答；难点在于两路内容必须一致，而且语音要尽早开始播放，又不能因为抢先开口而说出证据尚不支持的结论。这与给视频中的物体配音效具有不同目标，与前几节的媒体生成相比，生成内容由理解结果决定。

Qwen2.5-Omni的结构与训练见\ref{sec:vision-models-qwen-omni}节：Thinker读取视觉、声音和问题，生成文字及其隐表示；Talker接收这些隐表示和已采样的文字词元，自回归地产生语音编码；流式解码器把语音编码分块转换为梅尔频谱，再由声码器恢复波形\scite{vision-qwenomni2025}。Talker同时接收文字词元，帮助区分语义接近但发音不同的词，而不只依据模糊的语义向量猜测音节。

使用时一般不微调模型，而是组织输入和设置输出方式。视频按实际时间与声音一起送入，模型用对齐的时间位置表示两者；问题以文字或语音给出。输出端要选择是否需要语音，只需文字时可以不运行Talker，减少计算。语音编码需要解码才能播放：原方法把相邻编码分成块，解码当前块时读取有限的前后上下文，报告中的设置包含前面两块、当前块和后面一块。因此“流式”并不意味着拿到第一个文字词元就能零等待地发声；前看一块本身就需要一定积累。块越大，每块需要等待的生成越多，首次发声越晚；块越小，块边界越多，声音可能出现不连续。

响应延迟至少要区分首次发声和整段完成。设用户输入就绪时刻为0，第$k$块可播放音频准备好于$a_k$，播放时长为$d_k$，实际开始播放的时刻为$b_k$。忽略额外播放器缓冲时，
\begin{equation}
  b_1=a_1,\qquad
  b_k=\max(a_k,b_{k-1}+d_{k-1})\text{，}\quad k>1\text{。}
  \label{eq:vision-stream-playback}
\end{equation}
若音频尚未准备好，就必须等到$a_k$；若准备得很早，也要等前一块播完。这个排程关系与模型的具体参数无关，可以用来检查推理和播放是否衔接。

例如三块音频分别在0.8、1.2、1.7秒就绪，每块长0.5秒，则开始播放时刻为0.8、1.3、1.8秒，整段在2.3秒结束，中间无空隙。若第二块直到1.5秒才就绪，开始时刻变成0.8、1.5、2.0秒，出现0.2秒停顿。全部生成后再播放，在前一设置中要等到1.7秒才开始，并于3.2秒结束。以上均为教学排程，不是Qwen2.5-Omni的性能实测。首次发声时间包含输入编码、文字生成、语音编码预测及首块波形解码；视频越长，输入编码越久，首次发声越晚。

视觉问答还需决定是否可以提前说出某个结论。若视频中第二次接触尚未观察到，提前回答“花瓶先落桌”可能建立在猜测上。已经播放的词句无法像屏幕文字那样悄然撤回，纠正需要明确说出来。流式系统应在证据已经足够时输出相应内容，并把后来获得的信息与已说出的内容保持一致。

常见失败包括：语音漏读否定词，使“没有碰到”变成“碰到”；把“左边”说成“右边”；在数字中间断块造成停顿；以及文字回答正确而语音转写不同。验收时同时查看视觉证据、显示文字和语音转写，自然度、可懂度、首次发声时间和最长停顿分别记录；更早开口不能补偿答错。该模型生成文字与语音回答，不承担视频内容生成。

\section{成果验收}
\label{sec:vision-multigen-evaluation}

前几节在每种成果之后都留下了核验对象：视频的身份与运动、编辑的保留范围、声轨的起声时刻、联合音画的共同事件、三维内容的观测支持、语音回答的时序。本节把它们组织成交付前的验收。媒体成果既要独立可用，也要满足请求中规定的共同关系，还要说明为得到它花了多少代价。三者分别回答“单项好不好”“关系对不对”“代价多大”，按维度划分小节；把它们保留为分项结果，才能判断失败发生在哪里。自动特征分数、人类观看试听和具体任务检查承担不同职责，下文在各维度中分别说明；\ref{sec:vision-multigen-challenges}节定义的分布距离、逐片段指标和同步偏移，在这里按它们回答的问题归入相应维度。

\subsection{单项质量的核验}
\label{sec:vision-multigen-quality}

视频可检查模糊、闪烁、运动和身份变化；声音可检查失真、声源、语言内容和可懂度；空间资产可检查表面、纹理、拓扑和编辑用途。指标应作用于相应交付对象。图像质量分数不包含声音，可旋转展示的辐射场也不自动满足网格编辑要求。

视频的单项质量最难用一个数概括，两个基准给出不同的分解方式。VBench把质量拆成若干相互独立的维度，每个维度用一种专门工具测量画面层面的性质；VBench-2.0在此之上转向事件层面，问生成内容是否符合物理和常识。前者回答“看起来是否稳定、清晰”，后者回答“发生的事是否合理”。

\subsubsection{VBench}
\label{sec:vision-multigen-vbench}

单个总分，例如式\eqref{eq:vision-multigen-fvd}的FVD这类生成视频与真实视频的特征分布距离，会把闪烁、身份变化和文字不符混在一起，无法说明模型在哪方面失败。VBench把视频生成评价拆成16个层次化、相互解耦的维度，包括主体与背景一致性、时间闪烁、运动平滑度、动态程度、画面质量及条件语义等，再为每个维度选择专门的分析工具。其用途是定位差异，避免一个总分遮住主要失败\scite{vision-vbench2024}

十六个维度分为两大类。视频质量一类不看提示内容，又分时间质量与逐帧质量；视频与条件一致性一类检查提示中的语义和风格是否出现，最后另有一个整体一致性维度。表\ref{tab:vision-multigen-vbench-dims}列出每个维度所用的工具和计分方式。主体一致性已在式\eqref{eq:vision-multigen-subject-consistency}中给出，背景一致性用同一公式，只是把DINO特征换成CLIP图像特征。时间闪烁在提示要求静止场景的视频上计算相邻帧的平均绝对像素差$S$，再换算为$(255-S)/255$，越接近1越稳定。运动平滑度去掉奇数帧，用AMT插帧模型从偶数帧重新插出它们，再与原奇数帧比较平均绝对误差，误差小说明运动符合插帧模型所学的平滑先验。动态程度用RAFT估计光流，取每段视频中最大的5\%光流幅值的平均来判断视频是否在动，分数是判为非静止的视频比例。这些工具输出的是代理测量，例如特征相似度高不能证明每根手指都保持正确。

\begin{table*}[tbp]
 \centering
 \small
 \begin{tabularx}{\linewidth}{@{}p{22mm}p{26mm}p{34mm}X@{}}
  \toprule
  类别 & 维度 & 工具 & 计分方式 \\
  \midrule
  时间质量 & 主体一致性 & DINO & 式\eqref{eq:vision-multigen-subject-consistency}的帧间余弦相似度 \\
   & 背景一致性 & CLIP图像编码器 & 同一公式，换用CLIP特征 \\
   & 时间闪烁 & 像素差 & 静止场景相邻帧的平均绝对差，换算为$(255-S)/255$ \\
   & 运动平滑度 & AMT插帧模型 & 由偶数帧重插奇数帧，与原帧比较 \\
   & 动态程度 & RAFT光流 & 判为非静止的视频比例 \\
  逐帧质量 & 美学质量 & LAION美学预测器 & 0至10的逐帧评分线性归一化后平均 \\
   & 成像质量 & 在SPAQ上训练的MUSIQ & 逐帧分数除以100后平均 \\
  语义 & 对象类别 & GRiT检测 & 检出指定类别的帧比例 \\
   & 多个对象 & GRiT检测 & 所有指定对象同时出现的帧比例 \\
   & 人体动作 & 在Kinetics-400上训练的UMT & 均匀取16帧分类，检查目标动作是否在高置信候选中 \\
   & 颜色 & GRiT描述 & 含对象且有颜色描述的帧中，颜色正确的比例 \\
   & 空间关系 & 规则判定，沿用T2I-CompBench & 对象对的左右、上下关系正确率 \\
   & 场景 & Tag2Text描述 & 生成的场景描述是否与提示一致 \\
  风格 & 外观风格 & CLIP & 帧与风格描述的特征相似度 \\
   & 时间风格 & ViCLIP & 视频与镜头运动等时间风格描述的相似度 \\
  整体 & 整体一致性 & ViCLIP & 一般提示下的视频文字相似度 \\
  \bottomrule
 \end{tabularx}
 \caption{VBench十六个维度的评价工具。每个维度配有约100条提示的专用提示集；分数都是工具给出的代理测量，工具的误判计入评价误差。}
 \label{tab:vision-multigen-vbench-dims}
\end{table*}
\FloatBarrier

评价流程从固定提示集合开始，每个维度有自己约100条提示的子集；为每条提示保留规定数量的生成样本，用同一版本的分项脚本处理，再与相应的人类偏好标注比较。VBench用人类标注检查自动测量的对应关系，但某个基准中的对应不保证在所有新场景中成立。工具对生成缺陷的误判也应保留为评价误差。

一个几乎不动的玩具狗片段可能在外观一致性上很好，却没有完成“绕花瓶走一圈”。另一个片段运动丰富，却在转身后变成另一只狗。二者需要分别报告动作完成与身份保持，不能把低动态当成稳定性的无条件优点。模型间比较还须固定帧率、长度、样本数和生成预算，否则评分差异可能来自不同制作条件。

\Needspace{65mm}
\subsubsection{VBench-2.0}
\label{sec:vision-generation-physics-check}
\label{sec:vision-multigen-vbench2}

VBench的维度主要测量画面层面的性质。一个片段可以每帧清晰、运动平滑，却让杯子穿过桌面或让玩具狗凭空多出一条腿。这类错误不在画面层面，而在事件层面。VBench-2.0的2025年v2把评价目标改为内在的真实性，按人体表现、可控性、创造性、物理与常识五个方面组织测试，每个方面再细分为若干能力\scite{vision-vbench22025}。评价工具组合了两类：通用的视觉语言模型和大语言模型，用来回答关于视频内容的问题；以及专项检测器，例如为生成视频设计的异常检测方法。

在物理状态变化中，评价先根据提示形成具体的可见行为描述，再用多个问题核查生成视频。部分流程先排除初始状态不符合要求的片段，使状态变化判断对应原题。解释这种分数时，应同时报告初始条件的通过情况；只看过滤后的得分，会遗漏模型连题目起点都没有生成正确的失败。

对“玩具狗推倒花瓶”，教学核验可以分别问是否接触、花瓶是否倾斜、倾倒后是否保留同一个物体。仅有抬腿动作或一声碰撞都不足以回答这些问题。常识评价中的异常合并、突然复制与消失也需要跨帧观察，逐帧都像真图不表示帧间关系合理。

该基准的多视角一致性使用特征匹配等代理证据，因为生成视频通常没有对应真实三维模型。它不能直接测出真实几何误差。测试分数提高支持的是相应题目和评价流程下的改善，不是已经具备普遍物理模拟能力；视觉语言模型作为裁判也会出错，自动判定仍须与可见事实及人类检查对照。

两个基准的关系可以用同一组玩具狗视频说明。VBench能发现“转身后变成另一只狗”和“几乎不动”，因为这两者在特征相似度和光流上有明显信号；VBench-2.0能发现“玩具狗碰到花瓶，花瓶却没有倒，下一帧又出现在桌下”，因为它针对事件提问。VBench-2.0改变了评价工具的类型，从专用特征提取器换成以问题驱动的视觉语言模型，换来了对事件和常识的覆盖，代价是裁判模型本身的错误更难预测，分数的可重复性依赖裁判模型的版本。两者都只评价视频，不覆盖声轨、三维资产和语音回答。

\subsection{条件与跨输出关系的核验}
\label{sec:vision-multigen-relations}

单项质量合格，不等于成果满足请求。验收先把请求中的对象与关系落到同一份记录上。参考玩具狗的身份、杯子的材质、碰撞次数、声源及起声时刻相互关联，但应分别有可检查对象。表\ref{tab:vision-media-acceptance}列出贯穿示例中的核验方式，表中任务阈值需在生成前确定。

\begin{table*}[tbp]
 \centering
 \small
 \begin{tabularx}{\linewidth}{@{}p{25mm}p{44mm}X@{}}
  \toprule
  要求 & 主要证据 & 容易混淆的成功与失败 \\
  \midrule
  主体保持 & 原参考、转身前后与遮挡后画面 & 颜色相近仍可能改变耳朵、形状或身份 \\
  局部修改 & 原片段、编辑掩码、最终像素 & 指定区域改对，掩码外却有背景漂移 \\
  碰撞配声 & 接触时间、起声时间、事件次数 & 声音自然，但过早、过晚或多出一次 \\
  联合音画 & 共同事件、对白转写与口型 & 两路各自自然，却来自不同动作或声源 \\
  空间内容 & 参考覆盖、相机、表面与材质 & 新视图逼真，背面却无观测依据或不可编辑 \\
  语音回答 & 视觉依据、文字、语音与时序 & 文字正确，语音漏词，或早开口后长时间停顿 \\
  \bottomrule
 \end{tabularx}
 \caption{媒体制作的分项验收示例。证据应对应最终交付格式，掩码、相机、时间和事件来源需保存。各行回答不同问题，不能由某一项高分推断整件成果通过。}
 \label{tab:vision-media-acceptance}
\end{table*}
\FloatBarrier

跨输出关系可能随制作步骤改变。视频编辑把花瓶换成金属杯后，原陶瓷碰撞声可能已经不合适；三维资产改短了尾巴，沿用旧轮廓生成的多视角宣传图可能不再一致。应从受影响的关系出发重新核验，而不是只检查被修改的文件。

空间成果需标明哪些表面有输入观测支持，哪些由先验补出。可以在留出的真实视图上测量渲染误差；没有真实背面时，只能检查创作合理性、内部一致性及用途。把全部生成视角当作“验证视角”，会让同一个生成先验同时提供答案和检验依据。同样，用生成时使用的同步编码器来评价同步，也会让条件模型与评价模型重合，应当用独立的方法或人工检查补充。

\subsection{制作预算的核验}
\label{sec:vision-multigen-budget}

一条成功视频可能来自大量候选与多轮局部重做。比较制作能力时，应固定输入、时长、分辨率、模型版本、候选数量和编辑预算，并记录达到可用成果的总代价。只展示最佳样本会把筛选成本和失败概率都隐藏起来。

假设两个教学流程都交付一个四秒片段。甲生成4个候选，每个占用2分钟计算，再做一次1分钟的局部修改，累计计算量为9设备分钟；乙生成1个候选，占用5设备分钟。两者都通过当前验收，也不能因此说甲更省资源。若甲的4个候选并行生成，其等待时间可能短于9分钟，但设备累计时间仍是9分钟，人工观看与选择时间还要另计。本章的模型在代价结构上差别很大：DreamFusion每个对象要优化很久，Wonder3D约两分钟；LTX-2的模态引导每步需要三次前向；Wan的14B模型未经优化时单次约半小时。这些都应计入同一预算口径。

任务完成率的分母应是事先选定的请求，而不是最后留下的候选。可以记录规定预算内是否至少有一个结果通过，并同时给出每请求候选数、重做次数与失败类型。这样既描述了最终可用性，也保留模型一次生成的不稳定性。

\section{本章小结}
\label{sec:vision-multigen-summary}

开篇的请求要求玩具狗在转身后保持原样、碰撞真实发生、声音落在接触时刻，还要能从各个方向观察。本章逐步改变“哪些内容固定、哪些内容生成”的设定来回答这些要求。时间一致性方面，Lumiere在网络内部压缩时间，Wan先用自编码器压缩再用Transformer处理整段潜网格；修改与延伸则通过掩码、原片段或已知帧把“应当保留什么”写进条件。声音方面，MMAudio用高帧率同步特征决定何时发声，Movie Gen Audio用音频上下文延伸和修补长声轨；画面与声音都可变时，Ovi和LTX-2在每个生成步让两路互相读取，分阶段路线则以单向约束换取模块化。

空间方面，DreamFusion用分数蒸馏把二维先验传回共享的三维表示，Wonder3D在生成阶段就协调多视角颜色与法线，4D Gaussian Splatting用观测拟合随时间形变的高斯场景。三者的差别归结为一致性在哪里形成，以及未见区域由先验还是观测决定。视觉条件下的语音响应还把生成延伸到播放时序，文字正确、语音可懂和及时作答共同决定使用效果。

验收把这些要求拆回单项质量、条件与跨输出关系、制作预算三个维度。只有身份、动作、接触、声音和时间关系共同满足，开篇的片段才算完成，一幅漂亮画面、一条自然声轨或一个汇总分都不能替代。

回到第一节的五项挑战：时间一致性由整段共同建模与自编码器压缩回答，同步由按时间对齐的条件接入和每步双向读取回答，几何一致性由共享三维表示、跨视角注意力和观测拟合回答，预算则通过压缩、级联和统一的验收口径变得可以比较。只有长时状态保持一项，本章的方法仅改善了片段之间的衔接。

仍未解决的部分同样清楚。本章的视频与音视频模型都只覆盖约五秒到三十秒的片段，长时身份与场景漂移需要持续保存状态；局部修改需要协调新旧内容的边界；联合音画已能学习共同事件，但口型和多声源的验证仍依赖人工；多视角创建的背面始终没有观测依据；流式回答要协调证据到达和语言播放。扩大主干或增加条件只提供更多表示和学习机会，究竟解决了哪个问题，仍要通过相应时间范围、输入条件和制作预算下的结果回答。

\paragraph{扩展阅读}
视频生成方向，Video Diffusion Models把图像扩散的U形网络分解为空间层和时间层，图像与视频联合训练，并用重建引导的条件采样延伸片段，可看作Lumiere一类时空网络的早期形式；Make-A-Video从文生图模型出发，只用无文字标注的视频学习运动，说明文字与外观可以先在图文数据上学好；Imagen Video用基础视频模型加多级空间与时间超分模型组成级联，正是Lumiere试图避免的时间级联结构\scite{vision-vdm2022,vision-makeavideo2022,vision-imagenvideo2022}Stable Video Diffusion在潜空间视频扩散上着重讨论数据筛选与图像预训练、视频预训练、高质量微调三阶段训练，可与Wan按阶段提高分辨率的训练安排对照；VideoPoet则把视频和音频离散成词元，用只含解码器的语言模型在同一序列中完成生成、编辑和配声等多种任务，代表扩散与流匹配之外的自回归路线\scite{vision-svd2023,vision-videopoet2024}

声音生成方向，AudioLDM用对比学习得到的音频—文字嵌入作为条件，在潜空间扩散生成声音，只接受文字条件，可用来理解MMAudio为何还要加入视觉同步特征；Diff-Foley先用音画对比预训练得到与时间对齐的视觉特征，再以它们为条件做潜扩散配声，与MMAudio同样把同步放进条件接入\scite{vision-audioldm2023,vision-difffoley2023}三维生成方向，Magic3D沿用分数蒸馏，先优化低分辨率神经辐射场，再转为带纹理的网格在更高分辨率上细化；ProlificDreamer把三维参数视为随机变量，提出变分分数蒸馏，针对DreamFusion式蒸馏常见的过饱和、过平滑和缺乏多样性\scite{vision-magic3d2023,vision-prolificdreamer2023}Zero-1-to-3以相对相机位姿为条件微调图像扩散模型，从单幅图像逐个合成新视角，各视角之间没有共同计算，可与Wonder3D在生成阶段协调多视角的做法对照；LRM则不逐物体优化，用Transformer从单幅图像直接前馈预测三平面表示，几秒内给出三维结果，省去分数蒸馏的逐例优化，前提是有大规模多视角数据用于训练\scite{vision-zero123-2023,vision-lrm2024}评测方面，EvalCrafter用一组提示词和多项客观指标评估视频生成模型，并按用户评分调整各指标的组合权重，可与VBench的维度划分相互参照\scite{vision-evalcrafter2024}

